% La gestion des déchets (le tri des ordures) — Allemand Tle A — Cours signé OnBuch+
% Programme officiel MINESEC. Compiler avec : tectonic AL11-la-gestion-des-dechets-le-tri-des-ordures.tex
\newcommand{\DOCMATIERE}{Allemand}
\newcommand{\DOCNIVEAU}{Tle A}
\newcommand{\DOCMODULE}{Chapitre 3 : Environnement, santé et bien-être}
\newcommand{\DOCLECON}{AL11}
\newcommand{\DOCTITRE}{La gestion des déchets (le tri des ordures)}
% OnBuch+ — préambule « cours signés » ENRICHI (compilé avec tectonic / XeLaTeX)
% Système visuel « Pop » d'OnBuch+ : fond crème (jamais blanc), bordures encre
% épaisses, ombres « dures » décalées, titres Archivo Black en capitales.
% Version MATHÉMATIQUES : graphes (pgfplots/tikz), boîtes pédagogiques
% (définition, propriété, méthode, attention, le savais-tu, exercice résolu,
% exercice, TP, bilan), page de garde et signature OnBuch+ ancrée en bas.
%
% Macros attendues dans main.tex AVANT \input{preamble} :
%   \DOCMATIERE  \DOCNIVEAU  \DOCMODULE  \DOCLECON (numéro)  \DOCTITRE
\documentclass[11pt]{article}

\usepackage{fontspec}
\usepackage[ngerman,french]{babel} % français = langue principale ; l'allemand via \all{..} / textebox
\usepackage[a4paper,margin=2cm,top=3.4cm,bottom=2.8cm,headheight=1.4cm,footskip=1.3cm]{geometry}
\usepackage{amsmath,amssymb,amsfonts}
\usepackage{mathtools} % \xmapsto, \coloneqq... souvent utilisés par les modèles
\usepackage{cancel} % simplifications barrées (bilans de forces)
% \abs{z} (module, valeur absolue) et \norm{u} (norme), utilisés par les modèles
\providecommand{\abs}{}\let\abs\relax\DeclarePairedDelimiter{\abs}{\lvert}{\rvert}
\providecommand{\norm}{}\let\norm\relax\DeclarePairedDelimiter{\norm}{\lVert}{\rVert}
\usepackage[table]{xcolor} % [table] : fournit aussi \rowcolors (alternance de lignes)
\usepackage{pagecolor}
\usepackage{tikz}
\usetikzlibrary{arrows.meta,positioning,calc,shapes.geometric,shapes.misc,shapes.arrows,decorations.pathmorphing,decorations.pathreplacing,decorations.markings,patterns,shadows,mindmap,trees,fit,backgrounds,matrix,intersections,angles,quotes}
\usepackage{pgfplots}
\usepackage[version=4]{mhchem} % équations nucléaires \ce{^{235}_{92}U}
\usepackage[european,siunitx]{circuitikz} % schémas électriques
\pgfplotsset{compat=1.18}
\usepgfplotslibrary{fillbetween,groupplots} % aires entre courbes (calcul intégral)
\usepackage{enumitem}
\usepackage{fancyhdr}
% [most] charge toutes les bibliothèques tcolorbox (listings, minted, raster,
% documentation...) et gonfle inutilement la mémoire TeX sur un document long
% (~300 boîtes) : on ne charge que ce qui est réellement utilisé.
\usepackage[skins,breakable]{tcolorbox}
\usepackage{graphicx}
\usepackage{colortbl}
\usepackage{booktabs}
\usepackage{tabularx}
\usepackage{diagbox} % case d'en-tête diagonale des tableaux à double entrée
% Colonnes extensibles centrée / gauche / droite (C, L, R), souvent utilisées par les modèles
\newcolumntype{C}{>{\centering\arraybackslash}X}
\newcolumntype{L}{>{\raggedright\arraybackslash}X}
\newcolumntype{R}{>{\raggedleft\arraybackslash}X}
\usepackage{array}
\usepackage{multicol}
\usepackage{multirow}
\usepackage{siunitx}
% parse-numbers=false : accepte aussi \SI{2,5 \times 10^{-3}}{...}
\sisetup{locale=FR,output-decimal-marker={,},group-minimum-digits=4,parse-numbers=false}
\DeclareSIUnit\ohm{\ensuremath{\symup{\Omega}}} % U+2126 absent de la police
\DeclareSIUnit\division{div} % réglages d'oscilloscope (ms/div, V/div)
\DeclareSIUnit\tr{tr} % tours (vitesse de rotation, ex. \SI{1500}{\tr\per\minute})
\DeclareSIUnit\molar{M} % concentration molaire (ex. \SI{3}{\molar}, \SI{1}{\micro\molar})
\DeclareSIUnit\an{an} % année (le modèle confond parfois avec \year, un compteur TeX) — ex. \SI{5}{\kilo\meter\per\an}
\DeclareSIUnit\FCFA{FCFA} % franc CFA (ex. \SI{500}{\FCFA\per\kilo\gram})
\DeclareSIUnit\degreeBrix{\textdegree Brix} % teneur en sucre d'un sirop/jus (réfractométrie)
\DeclareSIUnit\month{mois} % le modèle utilise parfois \month, un compteur TeX, comme unité de durée
\usepackage{qrcode}
% \degree : commande fréquemment utilisée par les modèles pour un angle ou une
% température en dehors de \SI{}{} (non fournie par les paquets chargés).
\providecommand{\degree}{^{\circ}}
% \qed (fin de démonstration), utilisé par les modèles sans amsthm
\providecommand{\qed}{\hfill$\square$}
% \barre : double barre des valeurs interdites dans un tableau de variations (tkz-tab)
\providecommand{\barre}{\ensuremath{\|}}
% environnement proof (amsthm non chargé) utilisé par les modèles
\ifdefined\proof\else\newenvironment{proof}[1][Démonstration]{\par\noindent\textit{#1.}\ }{\qed\par}\fi
\usepackage{lastpage}
\usepackage{hyperref}
\hypersetup{hidelinks}

% ============================================================
% Palette « Pop » OnBuch+ — mêmes hex que la classe P (plus_ui.dart)
% ============================================================
\definecolor{poporange}{HTML}{F59321}
\definecolor{popdark}{HTML}{E07A0C}
\definecolor{popcream}{HTML}{FFF6E8}
\definecolor{popink}{HTML}{1C1714}
\definecolor{poppurple}{HTML}{7A5AE0}
\definecolor{popgreen}{HTML}{2E9E6A}
\definecolor{poppink}{HTML}{E0457B}
\definecolor{popblue}{HTML}{2F7DE1}
\definecolor{popgold}{HTML}{C98A12}
\definecolor{popmuted}{HTML}{8A7F76}
\definecolor{popline}{HTML}{EADFD2}
\definecolor{popgray}{HTML}{8A7F76}
\colorlet{popgrey}{popgray}
% Alias tolérants (le modèle nomme parfois les couleurs différemment)
\colorlet{popwhite}{white}
\colorlet{popblack}{popink}
\colorlet{popred}{poppink}
\colorlet{popyellow}{popgold}
% Teintes claires (fonds de tableaux/encadrés) — le modèle nomme parfois
% les couleurs avec un suffixe L (ex. popblueL) sans les avoir définies.
\colorlet{poporangeL}{poporange!20}
\colorlet{popdarkL}{popdark!20}
\colorlet{poppurpleL}{poppurple!20}
\colorlet{popgreenL}{popgreen!20}
\colorlet{poppinkL}{poppink!20}
\colorlet{popblueL}{popblue!20}
\colorlet{popgoldL}{popgold!20}
\colorlet{popredL}{popred!20}
\colorlet{popgrayL}{popgray!20}
% Teintes claires pour fonds de tableaux / zones
\colorlet{popblueL}{popblue!10!white}
\colorlet{poporangeL}{poporange!14!white}
\colorlet{popgreenL}{popgreen!12!white}
\colorlet{poppurpleL}{poppurple!12!white}
\colorlet{poppinkL}{poppink!12!white}
\colorlet{popgoldL}{popgold!14!white}

\pagecolor{popcream}

% ============================================================
% Typographie
% ============================================================
\defaultfontfeatures{Ligatures=TeX}
\setmainfont{PlusJakartaSans-Medium.ttf}[Path=../../../fonts/,BoldFont=PlusJakartaSans-Bold.ttf]
% Maths en sans-serif (Fira Math) avec chiffres et lettres droites de Plus
% Jakarta, pour que les formules se fondent dans le texte.
\usepackage[math-style=ISO,bold-style=ISO]{unicode-math}
\setmathfont{FiraMath-Regular.otf}[Path=../../../fonts/]
\setmathfont{PlusJakartaSans-Medium.ttf}[Path=../../../fonts/,range={up/num,up/{Latin,latin},\mathup}]
\usepackage{icomma} % virgule décimale sans espace en mode math
% Caractères mathématiques Unicode absents des polices de texte (Plus Jakarta,
% Archivo Black) : rendus par leur équivalent mathématique, en texte comme en
% formule — sinon ils s'affichent en carré vide (ex. « Arithmétique dans ℤ »).
\usepackage{newunicodechar}
\newunicodechar{ℝ}{\ensuremath{\mathbb{R}}}
\newunicodechar{ℤ}{\ensuremath{\mathbb{Z}}}
\newunicodechar{ℂ}{\ensuremath{\mathbb{C}}}
\newunicodechar{ℕ}{\ensuremath{\mathbb{N}}}
\newunicodechar{ℚ}{\ensuremath{\mathbb{Q}}}
\newunicodechar{𝔻}{\ensuremath{\mathbb{D}}}
\newunicodechar{α}{\ensuremath{\alpha}}
\newunicodechar{β}{\ensuremath{\beta}}
\newunicodechar{γ}{\ensuremath{\gamma}}
\newunicodechar{δ}{\ensuremath{\delta}}
\newunicodechar{ε}{\ensuremath{\varepsilon}}
\newunicodechar{θ}{\ensuremath{\theta}}
\newunicodechar{λ}{\ensuremath{\lambda}}
\newunicodechar{μ}{\ensuremath{\mu}}
\newunicodechar{σ}{\ensuremath{\sigma}}
\newunicodechar{τ}{\ensuremath{\tau}}
\newunicodechar{φ}{\ensuremath{\varphi}}
\newunicodechar{ψ}{\ensuremath{\psi}}
\newunicodechar{ω}{\ensuremath{\omega}}
\newunicodechar{Σ}{\ensuremath{\Sigma}}
% majuscules produites par \MakeUppercase dans les titres (α-aminés -> Α)
\newunicodechar{Α}{\ensuremath{\alpha}}
\newunicodechar{Β}{\ensuremath{\beta}}
\newunicodechar{Γ}{\ensuremath{\Gamma}}
\newunicodechar{Δ}{\ensuremath{\Delta}}
\newunicodechar{Ε}{\ensuremath{\varepsilon}}
\newunicodechar{Θ}{\ensuremath{\Theta}}
\newunicodechar{Λ}{\ensuremath{\Lambda}}
\newunicodechar{Μ}{\ensuremath{\mu}}
\newunicodechar{Τ}{\ensuremath{\tau}}
\newunicodechar{Φ}{\ensuremath{\Phi}}
\newunicodechar{Ψ}{\ensuremath{\Psi}}
\newunicodechar{Ω}{\ensuremath{\Omega}}
\newunicodechar{↦}{\ensuremath{\mapsto}}
\newunicodechar{∈}{\ensuremath{\in}}
\newunicodechar{∉}{\ensuremath{\notin}}
\newunicodechar{∧}{\ensuremath{\wedge}}
\newunicodechar{⇔}{\ensuremath{\Leftrightarrow}}
\newunicodechar{⟺}{\ensuremath{\Longleftrightarrow}}
\newunicodechar{⇒}{\ensuremath{\Rightarrow}}
\newunicodechar{⟹}{\ensuremath{\Longrightarrow}}
\newunicodechar{∘}{\ensuremath{\circ}}
\newunicodechar{∪}{\ensuremath{\cup}}
\newunicodechar{∩}{\ensuremath{\cap}}
\newunicodechar{≡}{\ensuremath{\equiv}}
\newunicodechar{⊂}{\ensuremath{\subset}}
\newunicodechar{⊥}{\ensuremath{\perp}}
\newunicodechar{∃}{\ensuremath{\exists}}
\newunicodechar{∀}{\ensuremath{\forall}}
\newunicodechar{∛}{\ensuremath{\sqrt[3]{\,}}}
\newunicodechar{∜}{\ensuremath{\sqrt[4]{\,}}}
\newunicodechar{⃗}{\ensuremath{{}^{\to}}}
\newunicodechar{ⁿ}{\ifmmode^{n}\else\textsuperscript{\ensuremath{n}}\fi}
\newunicodechar{ˣ}{\ifmmode^{x}\else\textsuperscript{\ensuremath{x}}\fi}
\newunicodechar{ᵇ}{\ifmmode^{b}\else\textsuperscript{\ensuremath{b}}\fi}
\newunicodechar{ʳ}{\ifmmode^{r}\else\textsuperscript{\ensuremath{r}}\fi}
\newunicodechar{ᵘ}{\ifmmode^{u}\else\textsuperscript{\ensuremath{u}}\fi}
\newunicodechar{ʸ}{\ifmmode^{y}\else\textsuperscript{\ensuremath{y}}\fi}
\newunicodechar{ᵃ}{\ifmmode^{a}\else\textsuperscript{\ensuremath{a}}\fi}
\newunicodechar{ᵉ}{\ifmmode^{e}\else\textsuperscript{\ensuremath{e}}\fi}
\newunicodechar{ᵏ}{\ifmmode^{k}\else\textsuperscript{\ensuremath{k}}\fi}
\newunicodechar{ᵖ}{\ifmmode^{p}\else\textsuperscript{\ensuremath{p}}\fi}
\newunicodechar{ᴺ}{\ifmmode^{N}\else\textsuperscript{\ensuremath{N}}\fi}
\newunicodechar{ᶜ}{\ifmmode^{c}\else\textsuperscript{\ensuremath{c}}\fi}
\newunicodechar{ʰ}{\ifmmode^{h}\else\textsuperscript{\ensuremath{h}}\fi}
\newunicodechar{ᵠ}{\ifmmode^{q}\else\textsuperscript{\ensuremath{q}}\fi}
\newunicodechar{ₙ}{\ifmmode_{n}\else\textsubscript{\ensuremath{n}}\fi}
\newunicodechar{ₐ}{\ifmmode_{a}\else\textsubscript{\ensuremath{a}}\fi}
\newunicodechar{ᵢ}{\ifmmode_{i}\else\textsubscript{\ensuremath{i}}\fi}
\newunicodechar{ᵣ}{\ifmmode_{r}\else\textsubscript{\ensuremath{r}}\fi}
\newunicodechar{ₖ}{\ifmmode_{k}\else\textsubscript{\ensuremath{k}}\fi}
\newunicodechar{ᵦ}{\ifmmode_{\beta}\else\textsubscript{\ensuremath{\beta}}\fi}
\newunicodechar{ₓ}{\ifmmode_{x}\else\textsubscript{\ensuremath{x}}\fi}
\newfontfamily\popdisplay{ArchivoBlack-Regular.ttf}[Path=../../../fonts/]
\newfontfamily\poplogo{SpaceGrotesk-Bold.ttf}[Path=../../../fonts/]
\newfontfamily\popsemi{PlusJakartaSans-SemiBold.ttf}[Path=../../../fonts/]
\newfontfamily\popxbold{PlusJakartaSans-ExtraBold.ttf}[Path=../../../fonts/]
\newfontfamily\popxboldls{PlusJakartaSans-ExtraBold.ttf}[Path=../../../fonts/,LetterSpace=3.0]

\setlist[enumerate]{leftmargin=1.7em,itemsep=5pt,topsep=4pt}
\setlist[itemize]{leftmargin=1.7em,itemsep=4pt,topsep=4pt,label={\color{poporange}\textbullet}}
\setlength{\parindent}{0pt}
\setlength{\parskip}{5pt}
\renewcommand{\arraystretch}{1.25}

% pgfplots : style Pop par défaut
\pgfplotsset{
  every axis/.append style={axis line style={popink,line width=1pt},tick style={popink},
    grid style={popline,line width=0.5pt},label style={font=\small},tick label style={font=\footnotesize}},
  popcurve/.style={poporange,line width=1.8pt,no markers,smooth},
}
% Même style disponible dans un \draw TikZ simple (hors pgfplots)
\tikzset{popcurve/.style={poporange,line width=1.8pt}}
\tikzset{popgrid/.style={popline,very thin}}
\colorlet{popcurve}{poporange} % le nom du style est parfois utilisé comme couleur

% ============================================================
% Pill / squiggle
% ============================================================
\newcommand{\poppill}[3]{%
  \tikz[baseline=-0.55ex]\node[draw=popink,line width=1.2pt,rounded corners=2.6pt,%
    fill=#2,inner xsep=6.5pt,inner ysep=3pt]{%
    \popxboldls\fontsize{7}{7}\selectfont\color{#3}\MakeUppercase{#1}};%
}
\newcommand{\popsquiggle}[1]{%
  \tikz[baseline=-0.4ex]{
    \draw[#1,line width=2.6pt,line cap=round]
      plot [smooth] coordinates {(0,0.09) (0.34,0.01) (0.68,0.21) (1.02,0.01) (1.36,0.10)};
  }%
}
% Mot-clé surligné (marqueur orange sous le texte)
\newcommand{\cle}[1]{{\popxbold\color{popdark}#1}}

% ============================================================
% En-tête / pied
% ============================================================
\pagestyle{fancy}
\fancyhf{}
\renewcommand{\headrulewidth}{0pt}
\renewcommand{\footrulewidth}{0pt}
\lhead{\raisebox{-0.15ex}{\poplogo\fontsize{13}{13}\selectfont\color{popink}OnBuch\color{poporange}+}}
\rhead{\poppill{\DOCMATIERE\ \textbullet\ \DOCNIVEAU\ \textbullet\ Leçon \DOCLECON}{white}{popink}}
\lfoot{\popxbold\fontsize{7.6}{7.6}\selectfont\color{popmuted}\MakeUppercase{Cours signé OnBuch+}\ \textbullet\ {\popsemi\DOCTITRE}}
\rfoot{\tikz[baseline=-0.6ex]\node[rounded corners=8.5pt,fill=popink,minimum height=17pt,minimum width=17pt,inner xsep=5pt,inner sep=0]{\popxbold\fontsize{8}{8}\selectfont\color{popcream}\thepage\,/\,\pageref*{LastPage}};}

% ============================================================
% Titres
% ============================================================
\newcommand{\chaptitre}[2]{%
  \begin{tcolorbox}[enhanced,colback=poporange,colframe=popink,boxrule=2.6pt,arc=11pt,
      drop shadow=popink,left=15pt,right=15pt,top=12pt,bottom=11pt,before skip=3pt,after skip=18pt]
    \poppill{#2}{popink}{popcream}\par\vspace{9pt}
    {\raggedright\hyphenpenalty=10000\exhyphenpenalty=10000\popdisplay\fontsize{22}{24}\selectfont\color{popcream}\MakeUppercase{#1}\par}
  \end{tcolorbox}
}
% Section numérotée : pastille orange avec le numéro + titre Archivo
\newcounter{popsec}
\newcommand{\coursec}[1]{%
  \stepcounter{popsec}\setcounter{popsub}{0}%
  \par\vspace{16pt}\noindent
  \begin{minipage}{\linewidth}
  \tikz[baseline=-0.7ex]\node[draw=popink,line width=1.6pt,fill=poporange,rounded corners=4pt,minimum size=22pt,inner sep=0]{\popdisplay\fontsize{12}{12}\selectfont\color{popcream}\thepopsec};%
  \hspace{8pt}{\popdisplay\fontsize{15}{17}\selectfont\color{popink}\MakeUppercase{#1}}\par
  \vspace{4pt}\noindent{\color{poporange}\rule{36pt}{3.6pt}}
  \end{minipage}\par\nopagebreak\vspace{8pt}
}
\newcounter{popsub}
\newcommand{\courssub}[1]{%
  \stepcounter{popsub}%
  \par\vspace{9pt}\noindent{\popxbold\fontsize{11.5}{13}\selectfont\color{popink}{\color{poporange}\thepopsec.\thepopsub}\hspace{6pt}#1}\par\nopagebreak\vspace{4pt}
}

% ============================================================
% Boîtes pédagogiques — même recette : fond blanc, bordure épaisse colorée,
% ombre dure encre, badge plein. Titre optionnel [#1].
% ============================================================
\tcbset{popbase/.style={breakable,enhanced,colback=white,boxrule=2.3pt,arc=9pt,
  shadow={3pt}{-3pt}{0pt}{popink},left=14pt,right=14pt,top=11pt,bottom=11pt,before skip=11pt,after skip=15pt,
  fontupper=\popsemi\fontsize{10.6}{14.5}\selectfont\color{popink},
  fontlower=\popsemi\fontsize{10.6}{14.5}\selectfont\color{popink}}}
% Coche dessinée (le glyphe ✓ n'existe pas dans les polices du document)
\newcommand{\popcheck}[1]{\tikz[baseline=-0.5ex]\draw[#1,line width=1.6pt,line cap=round,line join=round](-0.13,0.0)--(-0.04,-0.09)--(0.13,0.1);}
\providecommand{\coche}{\popcheck{popgreen}}
\providecommand{\resultat}[1]{\par\smallskip\noindent\fcolorbox{popgreen}{popgreenL}{\parbox{\dimexpr\linewidth-2\fboxsep-2\fboxrule\relax}{\textbf{Résultat.} #1}}\par\smallskip}
\newcommand{\popboxhead}[3]{\poppill{#1}{#2}{white}\ifx\relax#3\relax\else\hspace{6pt}{\popxbold\fontsize{10.6}{12}\selectfont\color{popink}#3}\fi\par\vspace{7pt}}

\newtcolorbox{aretenir}[1][]{popbase,colframe=popgreen,before upper={\popboxhead{À retenir}{popgreen}{#1}}}
% Texte en allemand (document de lecture, dialogue, extrait) : cadre
% d'étude. Un titre optionnel (source, auteur) s'affiche dans l'en-tête.
\newtcolorbox{textebox}[1][]{popbase,colframe=popblue,before upper={\popboxhead{Text}{popblue}{#1}\begin{otherlanguage}{ngerman}},after upper={\end{otherlanguage}}}
% Marques ✓ / ✗ (forme correcte / fautive) souvent utilisées par les modèles
\providecommand{\cross}{\textcolor{poppink}{\ensuremath{\times}}}
\providecommand{\xmark}{\cross}\providecommand{\crossmark}{\cross}\providecommand{\cmark}{\ensuremath{\checkmark}}
% Ligne de réponse / justification à compléter (inventée par les modèles)
\providecommand{\justification}[1]{\par\noindent\textit{Justification~:} \dotfill\par}
% \textipa (package tipa non chargé) : la police ne couvre pas l'API -> simple passage
\providecommand{\textipa}[1]{#1}
% Soulignement (ulem non chargé) : \uline, \ul
\providecommand{\uline}[1]{\underline{#1}}\providecommand{\ul}[1]{\underline{#1}}
% \alert (beamer) inventée par les modèles : texte mis en évidence
\providecommand{\alert}[1]{\textbf{\textcolor{poppink}{#1}}}
% variantes ASCII d'ordinaux écrites en macro
\providecommand{\iere}{\textsuperscript{re}}\providecommand{\ieme}{\textsuperscript{e}}\providecommand{\ere}{\textsuperscript{re}}\providecommand{\eme}{\textsuperscript{e}}\providecommand{\er}{\textsuperscript{er}}\providecommand{\ie}{\textsuperscript{e}}
% Ordinaux français écrits en macro par les modèles : 1\ière, 3\ième
\newcommand{\ière}{\textsuperscript{re}}\newcommand{\ième}{\textsuperscript{e}}
% \exemple (inventée par les modèles) : étiquette « Exemple : »
\providecommand{\exemple}{\textbf{Exemple~:}\ }
% \square utilisable aussi hors mode math (cases à cocher)
\AtBeginDocument{\let\squareMath\square\renewcommand{\square}{\ensuremath{\squareMath}}}
% Mot ou phrase en allemand à l'intérieur d'un paragraphe français
\newcommand{\all}[1]{\ifmmode\text{\textit{#1}}\else\begingroup\selectlanguage{ngerman}\itshape #1\endgroup\fi}
\let\esp\all % alias toléré
\newtcolorbox{exemplebox}[1][]{popbase,colframe=poporange,before upper={\popboxhead{Exemple}{poporange}{#1}}}
\newtcolorbox{definition}[1][]{popbase,colframe=popblue,before upper={\popboxhead{Définition}{popblue}{#1}}}
\newtcolorbox{propriete}[1][]{popbase,colframe=poppurple,before upper={\popboxhead{Propriété}{poppurple}{#1}}}
\newtcolorbox{methode}[1][]{popbase,colframe=popgold,before upper={\popboxhead{Méthode}{popgold}{#1}}}
\newtcolorbox{attention}[1][]{popbase,colframe=poppink,before upper={\popboxhead{Attention}{poppink}{#1}}}
\newtcolorbox{experience}[1][]{popbase,colframe=popblue,colback=popblueL,before upper={\popboxhead{Expérience}{popblue}{#1}}}
% « Le savais-tu ? » : carte violette pleine (contexte, culture, Cameroun)
\newtcolorbox{savaistu}[1][]{popbase,colback=poppurple,colframe=popink,
  fontupper=\popsemi\fontsize{10.4}{14.2}\selectfont\color{white},
  before upper={\poppill{Le savais-tu\,?}{white}{poppurple}\ifx\relax#1\relax\else\hspace{6pt}{\popxbold\color{white}#1}\fi\par\vspace{7pt}}}
% Exercice résolu : énoncé en haut, \tcblower puis la solution sur fond crème
\newcounter{popexo}\newcounter{popexr}
\newtcolorbox{exoresolu}[1][]{popbase,colframe=poporange,colbacklower=popcream,
  segmentation style={popink,line width=1.2pt,dashed},
  before upper={\stepcounter{popexr}\popboxhead{Exercice résolu \thepopexr}{poporange}{#1}},
  before lower={\poppill{Solution}{popink}{popcream}\par\vspace{6pt}}}
% Exercice (non résolu en place) — la correction est dans la partie Corrigés
\newtcolorbox{exercice}[1][]{popbase,colframe=popink,
  before upper={\stepcounter{popexo}\popboxhead{Exercice \thepopexo}{popink}{#1}}}
% Corrigé
\newtcolorbox{corrige}[1][]{popbase,colframe=popgreen,colback=popgreenL,
  before upper={\popboxhead{Corrigé}{popgreen}{#1}}}
% Objectifs / prérequis / bilan
\newtcolorbox{objectifs}{popbase,colframe=popink,colback=poporangeL,
  before upper={\popboxhead{Objectifs de la leçon}{popink}{}}}
\newtcolorbox{prerequis}{popbase,colframe=popink,colback=popblueL,
  before upper={\popboxhead{Prérequis}{popblue}{}}}
\newtcolorbox{bilan}{popbase,colframe=popink,colback=white,boxrule=2.8pt,
  before upper={\popboxhead{Fiche bilan}{poporange}{L'essentiel à connaître pour l'examen}}}
% Figure encadrée avec légende
\newtcolorbox{popfigure}[1][]{enhanced,breakable=false,colback=white,colframe=popline,boxrule=1.4pt,arc=8pt,
  left=10pt,right=10pt,top=10pt,bottom=8pt,before skip=10pt,after skip=12pt,halign=center,
  lower separated=false,halign lower=center,
  fontlower=\popsemi\fontsize{9}{11.5}\selectfont\color{popmuted},
  #1}
\newcommand{\legende}[1]{\par\vspace{6pt}{\popsemi\fontsize{9}{11.5}\selectfont\color{popmuted}#1}}

% ============================================================
% Signature OnBuch+
% ============================================================
\newcommand{\poplogoinline}[1]{{\poplogo\fontsize{#1}{#1}\selectfont\color{popink}OnBuch\color{poporange}+}}
\newcommand{\signatureonbuch}{%
  \begin{tcolorbox}[enhanced,colback=popink,colframe=popink,boxrule=2.2pt,arc=10pt,
      drop shadow=poporange,left=14pt,right=14pt,top=10pt,bottom=10pt,before skip=0pt,after skip=0pt]
    \tikz[baseline=-0.7ex]\node[circle,fill=popgreen,draw=popcream,line width=1.3pt,minimum size=22pt,inner sep=0]{\popcheck{white}};%
    \hspace{9pt}\begin{minipage}[c]{0.62\linewidth}
      {\popxbold\fontsize{12}{13}\selectfont\color{popcream}Signé\ \textbullet\ {\poplogo OnBuch\color{poporange}+}}\par\vspace{2pt}
      {\raggedright\popsemi\fontsize{8}{10.5}\selectfont\color{popline}\MakeUppercase{Équipe pédagogique OnBuch+}\\\MakeUppercase{Conforme au programme officiel MINESEC}\par}
    \end{minipage}\hfill
    {\poplogo\fontsize{22}{22}\selectfont\color{popcream}OnBuch\color{poporange}+}
  \end{tcolorbox}%
}

% ============================================================
% Page de garde
% #1 titre · #2 matière · #3 niveau · #4 module · #5 n° leçon · #6 durée ·
% #7 description / objectifs (texte ou itemize)
% ============================================================
\newcommand{\pagegarde}[7]{%
  \thispagestyle{empty}
  \newgeometry{a4paper,margin=1.7cm,top=1.6cm,bottom=1.4cm}%
  \begin{tikzpicture}[remember picture,overlay]
    \fill[poporange!18] ([xshift=-3.2cm,yshift=-3.6cm]current page.north east) circle (4.6cm);
    \fill[poppurple!14] ([xshift=2.4cm,yshift=7.6cm]current page.south west) circle (3.8cm);
  \end{tikzpicture}%
  {\poplogo\fontsize{30}{30}\selectfont\color{popink}OnBuch\color{poporange}+}\hfill
  \raisebox{0.6ex}{\poppill{Cours signé \textbullet\ Premium}{popink}{popcream}}\par
  \vspace{1pt}\popsquiggle{poppurple}\par
  \vspace{8pt}
  \begin{tcolorbox}[enhanced,colback=poporange,colframe=popink,boxrule=2.8pt,arc=13pt,
      drop shadow=popink,left=18pt,right=18pt,top=12pt,bottom=13pt,after skip=10pt]
    \poppill{#2 \textbullet\ #3}{popink}{popcream}\hspace{5pt}\poppill{#4}{white}{popink}\par\vspace{8pt}
    {\popdisplay\fontsize{14}{14}\selectfont\color{popink}LEÇON #5}\par\vspace{4pt}
    {\raggedright\hyphenpenalty=10000\exhyphenpenalty=10000\popdisplay\fontsize{25}{27}\selectfont\color{popcream}\MakeUppercase{#1}\par}
    \vspace{8pt}
    {\popxbold\fontsize{9.5}{9.5}\selectfont\color{popink}\MakeUppercase{Durée conseillée : #6}}
  \end{tcolorbox}
  \begin{tcolorbox}[enhanced,colback=white,colframe=popink,boxrule=2.2pt,arc=10pt,
      drop shadow=popink,left=16pt,right=16pt,top=10pt,bottom=10pt,after skip=8pt]
    \poppill{Dans cette leçon}{poppurple}{white}\par\vspace{5pt}
    \popsemi\fontsize{9.3}{12.6}\selectfont\color{popink}#7
  \end{tcolorbox}
  \noindent\begin{minipage}[t]{0.487\linewidth}
    \begin{tcolorbox}[enhanced,colback=popgreen,colframe=popink,boxrule=2.2pt,arc=10pt,
        drop shadow=popink,left=11pt,right=11pt,top=10pt,bottom=10pt,height=3.1cm,valign=center]
      \tcbox[colback=white,colframe=popink,boxrule=1.3pt,arc=5pt,boxsep=0pt,nobeforeafter,
        left=3pt,right=3pt,top=3pt,bottom=3pt,valign=center]{\qrcode[height=1.9cm]{https://play.google.com/store/apps/details?id=com.luvvix.onbuch&hl=fr}}%
      \hspace{8pt}\begin{minipage}[c]{0.5\linewidth}
        {\popdisplay\fontsize{11}{12.5}\selectfont\color{white}\MakeUppercase{Télécharge OnBuch}}\par\vspace{4pt}
        {\raggedright\popsemi\fontsize{8.4}{11}\selectfont\color{white}Gratuit \textbullet\ Google Play\\Tuteur IA, annales, concours, résultats\par}
      \end{minipage}
    \end{tcolorbox}
  \end{minipage}\hfill
  \begin{minipage}[t]{0.487\linewidth}
    \begin{tcolorbox}[enhanced,colback=poppurple,colframe=popink,boxrule=2.2pt,arc=10pt,
        drop shadow=popink,left=11pt,right=11pt,top=10pt,bottom=10pt,height=3.1cm,valign=center]
      \tikz[baseline=-0.7ex]\node[circle,fill=white,draw=popink,line width=1.3pt,minimum size=1.6cm,inner sep=0]{\popdisplay\fontsize{24}{24}\selectfont\color{poppurple}+};%
      \hspace{8pt}\begin{minipage}[c]{0.6\linewidth}
        {\popdisplay\fontsize{11}{12.5}\selectfont\color{white}\MakeUppercase{Passe à OnBuch+}}\par\vspace{4pt}
        {\raggedright\popsemi\fontsize{8.4}{11}\selectfont\color{white}Tous les cours signés, Tuteur IA illimité, fiches PDF — onglet OnBuch+ de l'app\par}
      \end{minipage}
    \end{tcolorbox}
  \end{minipage}
  \vspace{8pt}
  \popcontenu
  \vfill
  \signatureonbuch
  \restoregeometry
  \newpage
}

% Rangée « Ce cours contient » (page de garde)
\newcommand{\poptile}[3]{% #1 couleur #2 gros texte #3 légende
  \begin{tcolorbox}[enhanced,colback=#1,colframe=popink,boxrule=2pt,arc=9pt,shadow={2.5pt}{-2.5pt}{0pt}{popink},
      left=6pt,right=6pt,top=9pt,bottom=9pt,halign=center,valign=center,height=1.75cm,width=\linewidth,nobeforeafter]
    {\popdisplay\fontsize{12.5}{13}\selectfont\color{white}\MakeUppercase{#2}}\par\vspace{3pt}
    {\centering\popsemi\fontsize{7.8}{9.5}\selectfont\color{white}#3\par}
  \end{tcolorbox}}
\newcommand{\popcontenu}{%
  \par\noindent{\popxboldls\fontsize{7.5}{7.5}\selectfont\color{popmuted}CE COURS SIGNÉ CONTIENT}\par\vspace{7pt}
  \noindent\begin{minipage}[t]{0.235\linewidth}\poptile{popblue}{Cours}{complet, illustré, pas à pas}\end{minipage}\hfill
  \begin{minipage}[t]{0.235\linewidth}\poptile{poporange}{Méthodes}{et exercices résolus}\end{minipage}\hfill
  \begin{minipage}[t]{0.235\linewidth}\poptile{poppink}{Exercices}{type examen + corrigés}\end{minipage}\hfill
  \begin{minipage}[t]{0.235\linewidth}\poptile{popgreen}{Fiche bilan}{l'essentiel à retenir}\end{minipage}\par}

% Sommaire
\newtcolorbox{sommaire}{enhanced,colback=white,colframe=popink,boxrule=2.2pt,arc=10pt,
  drop shadow=popink,left=18pt,right=18pt,top=13pt,bottom=13pt,before skip=6pt,after skip=20pt,
  before upper={\poppill{Sommaire}{poppurple}{white}\par\vspace{9pt}\popsemi\fontsize{10.6}{14.5}\selectfont\color{popink}}}

% Signature de fin de cours — poussée tout en bas de la dernière page
\newcommand{\findecours}{%
  \par\vfill\nopagebreak
  \begin{center}\popsquiggle{poporange}\end{center}\vspace{4pt}
  \signatureonbuch
}
\AtBeginDocument{\def\llbracket{\mathopen{[\mkern-3.5mu[}}\def\rrbracket{\mathclose{]\mkern-3.5mu]}}} % intervalles d'entiers (glyphe ⟦ absent de la police)
% Glyphes absents de Fira Math : redéfinis à partir de glyphes disponibles
\AtBeginDocument{%
  \def\setminus{\mathbin{\backslash}}\let\smallsetminus\setminus
  \def\vdots{\mathord{\vbox{\baselineskip3.2pt\lineskiplimit0pt\kern4pt\hbox{.}\hbox{.}\hbox{.}}}}%
  \def\ddots{\mathinner{\mkern1mu\raise6.5pt\hbox{.}\mkern2mu\raise3.6pt\hbox{.}\mkern2mu\raise0.7pt\hbox{.}\mkern1mu}}%
  \def\triangle{\mathord{\scalebox{0.9}{$\symup{\Delta}$}}}%
  \def\nabla{\mathord{\raisebox{\depth}{\scalebox{1}[-1]{$\symup{\Delta}$}}}}%
  \def\checkmark{\popcheck{popdark}}%
  \def\star{\mathbin{\ast}}\def\bigstar{\mathord{\ast}}%
  \def\diamond{\mathbin{\scalebox{0.6}{$\Diamond$}}}%
  \def\bigcup{\mathop{\mathchoice{\raisebox{-0.3em}{\scalebox{1.7}{$\cup$}}}{\scalebox{1.25}{$\cup$}}{\cup}{\cup}}\displaylimits}%
  \def\bigcap{\mathop{\mathchoice{\raisebox{-0.3em}{\scalebox{1.7}{$\cap$}}}{\scalebox{1.25}{$\cap$}}{\cap}{\cap}}\displaylimits}%
  \def\lesseqqgtr{\mathrel{\lessgtr}}\def\bigoplus{\mathop{\oplus}\displaylimits}%
}
\providecommand{\ssi}{si et seulement si} % abréviation fréquente
\AtBeginDocument{\providecommand{\wideparen}{\overparen}} % arc de cercle

%% FIGFIX-BEGIN v3 — correctifs globaux des figures (docs/onbuch-generation/tools/figfix)
\usepackage{adjustbox}\usepackage{etoolbox}
% toute figure TikZ/pgfplots plus large que la ligne est réduite à la largeur disponible
\BeforeBeginEnvironment{tikzpicture}{\begin{adjustbox}{max width=\linewidth}}
\AfterEndEnvironment{tikzpicture}{\end{adjustbox}}
% légendes pgfplots lisibles par-dessus les courbes
\pgfplotsset{every axis/.append style={legend style={fill=white,fill opacity=0.92,text opacity=1,draw=black!15,inner sep=3pt}}}
% étiquettes d'arêtes (syntaxe edge["..."]) avec fond blanc
\tikzset{every edge quotes/.append style={fill=white,inner sep=1.2pt}}
% bulles de cartes mentales : texte à la ligne dans la bulle, bulle un peu plus grande
\tikzset{every concept/.append style={text width=2.0cm,align=center,minimum size=2.3cm,inner sep=2pt}}
%% FIGFIX-END
\begin{document}

\pagegarde{La gestion des déchets (le tri des ordures)}{Allemand}{Tle A}{Chapitre 3 : Environnement, santé et bien-être}{AL11}{2 h}{Cette leçon de type « texte » propose la lecture et le commentaire d'un texte allemand sur la gestion des déchets et le tri des ordures. Elle permet d'acquérir le vocabulaire thématique de l'environnement et de maîtriser les verbes à particule séparable, le passif avec verbes de modalité et l'impératif.
\vspace{4pt}
\begin{multicols}{2}\small
\begin{itemize}
\item À la fin de cette leçon, je sais comprendre globalement puis en détail un texte allemand sur la gestion des déchets et le tri des ordures.
\item À la fin de cette leçon, je sais employer correctement le vocabulaire du thème (der Müll, die Mülltrennung, das Recycling, entsorgen, sauber...) avec articles et pluriels.
\item À la fin de cette leçon, je sais reconnaître et conjuguer les verbes à particule séparable au présent et au parfait.
\item À la fin de cette leçon, je sais former et utiliser le passif avec les verbes de modalité (muss getrennt werden, kann recycelt werden...).
\item À la fin de cette leçon, je sais former l'impératif (2e personne du singulier, du pluriel, forme de politesse) pour donner des consignes écologiques.
\item À la fin de cette leçon, je sais répondre à des questions de compréhension en phrases complètes.
\end{itemize}\end{multicols}}

\begin{sommaire}
\begin{multicols}{2}
\begin{enumerate}
\item Texte support et lecture globale
\item Compréhension détaillée
\item Lexique thématique : le tri des déchets
\item Grammaire 1 : les verbes à particule séparable
\item Grammaire 2 : le passif avec verbes de modalité et l'impératif
\item Méthode du commentaire et production guidée
\item Activité d'intégration
\item Méthodes et exercices résolus
\item Exercices
\item Corrigés des exercices
\item Fiche bilan
\end{enumerate}
\end{multicols}
\end{sommaire}

\begin{objectifs}
\begin{itemize}
\item À la fin de cette leçon, je sais comprendre globalement puis en détail un texte allemand sur la gestion des déchets et le tri des ordures.
\item À la fin de cette leçon, je sais employer correctement le vocabulaire du thème (der Müll, die Mülltrennung, das Recycling, entsorgen, sauber...) avec articles et pluriels.
\item À la fin de cette leçon, je sais reconnaître et conjuguer les verbes à particule séparable au présent et au parfait.
\item À la fin de cette leçon, je sais former et utiliser le passif avec les verbes de modalité (muss getrennt werden, kann recycelt werden...).
\item À la fin de cette leçon, je sais former l'impératif (2e personne du singulier, du pluriel, forme de politesse) pour donner des consignes écologiques.
\item À la fin de cette leçon, je sais répondre à des questions de compréhension en phrases complètes.
\item À la fin de cette leçon, je sais rédiger un résumé et un court paragraphe de production guidée sur le tri des déchets au Cameroun.
\item À la fin de cette leçon, je sais comparer les pratiques de gestion des déchets au Cameroun et dans les pays germanophones.
\end{itemize}
\end{objectifs}

\begin{prerequis}
\begin{itemize}
\item Le présent et le parfait (Perfekt) des verbes faibles et forts (Première)
\item Le passif simple au présent (Das Haus wird gebaut) vu en Première
\item Les verbes de modalité au présent (müssen, können, sollen, dürfen, wollen)
\item Le vocabulaire de base de l'environnement (die Umwelt, die Natur, die Verschmutzung)
\item La place du verbe dans la phrase déclarative et interrogative
\end{itemize}
\end{prerequis}

\newpage

\coursec{Texte support et lecture globale}

\courssub{Présentation du texte et situation d'accroche}

À Yaoundé comme à Douala, les ordures s'accumulent parfois dans les rues, obstruant les caniveaux et polluant l'air, alors qu'en Allemagne, en Autriche ou en Suisse, chaque foyer trie ses déchets dans des poubelles de couleurs différentes : \cle{gelbe Tonne} (jaune) pour le plastique et les emballages, \cle{blaue Tonne} (bleue) pour le papier, \cle{braune Tonne} (brune) pour les biodéchets. Que se passerait-il si le Cameroun adoptait le système allemand de \cle{Mülltrennung} ? Comment les pays germanophones sont-ils devenus des champions du recyclage ? Dans cette leçon, nous lisons un texte authentique sur la gestion des déchets en Allemagne et nous apprenons à donner des consignes (impératif) et à décrire des obligations ou des processus (passif avec verbes de modalité).

Le texte support s'intitule \textbf{„Mülltrennung in Deutschland – ein Vorbild für die Welt?“}. C'est un \cle{texte informatif et argumentatif} d'environ 180 mots, adapté du niveau B1/B2, destiné à un public de lycéens germanophones ou apprenants avancés. Il décrit le système de tri, son fonctionnement juridique (Verpackungsverordnung) et ouvre une perspective comparative avec le Cameroun.

\courssub{Lecture globale et hypothèses de sens}

\begin{methode}[Comment aborder un texte allemand inconnu — 5 étapes]
\begin{enumerate}
    \item \textbf{Lire le titre et anticiper :} Repérer les mots transparents (\all{Recycling}, \all{System}, \all{Problem}, \all{Plastik}, \all{Modell}) et formuler des hypothèses sur le thème.
    \item \textbf{Repérer la structure visuelle :} Paragraphes, mots en gras, connecteurs logiques (\all{außerdem}, \all{jedoch}, \all{deshalb}).
    \item \textbf{Lecture silencieuse globale :} Lire une première fois sans dictionnaire pour saisir le \cle{sens global} (qui ? quoi ? où ? quand ?).
    \item \textbf{Consulter le glossaire de marge :} Clarifier les 8--10 mots-clés bloquants fournis par l'enseignant.
    \item \textbf{Vérifier les hypothèses :} Répondre aux questions globales (W-Fragen) et reformuler oralement le contenu en 2--3 phrases en français ou en allemand.
\end{enumerate}
\end{methode}

\courssub{Texte support : « Mülltrennung in Deutschland »}

\begin{textebox}[Mülltrennung in Deutschland – ein Vorbild für die Welt?]
In Deutschland ist die Mülltrennung gesetzliche Pflicht. Jeder Haushalt besitzt mindestens vier verschiedene Behälter: die gelbe Tonne für Verpackungen aus Plastik, Metall und Verbundstoffen, die blaue Tonne für Papier und Pappe, die braune Tonne für Biomüll wie Küchen- und Gartenabfälle sowie die graue Tonne für den Restmüll. Glasflaschen und Einweckgläser werden nicht in die Tonnen geworfen, sondern in öffentliche Glascontainer nach Farben sortiert: Weißglas, Grünglas und Braunglas.

Das System funktioniert nur, wenn alle Bürger mitmachen. Falsch sortierter Müll wird oft nicht geleert und kann Bußgelder nach sich ziehen. Durch das konsequente Recycling werden wertvolle Rohstoffe gespart und die Umwelt geschützt. Seit 1991 regelt die Verpackungsverordnung die Rücknahme von Verkaufsverpackungen. Der „Grüne Punkt“ auf vielen Produkten zeigt an, dass der Hersteller für die Entsorgung bezahlt hat.

Auch in Österreich und in der Schweiz gibt es ähnliche Trennsysteme mit hohen Recyclingquoten. In Kamerun hingegen landet ein großer Teil des Mülls noch auf wilden Deponien oder wird unkontrolliert verbrannt. Ein Umdenken ist dringend nötig: Umweltbildung in den Schulen, klare Gesetze und eine funktionierende Infrastruktur könnten die Situation nachhaltig verbessern.
\end{textebox}

\courssub{Glossaire de marge — Mots-clés du texte}

\begin{center}
\begin{tabularx}{\linewidth}{|X|X|X|}
\hline
\rowcolor{popblueL}
\textbf{Deutsch (Artikel, Plural)} & \textbf{Français} & \textbf{Remarque / Famille de mots} \\
\hline
der Abfall, die Abfälle & le déchet, l'ordure & \all{abfallen} (tomber), \all{der Abfalleimer} (poubelle) \\
\hline
die Wertstofftonne, die Wertstofftonnen & la poubelle de tri (matériaux valorisables) & \all{der Wertstoff} (matière recyclable) \\
\hline
der Restmüll (kein Plural) & les ordures résiduelles (non recyclables) & \all{der Rest} (le reste) \\
\hline
wegwerfen (wirft weg, warf weg, hat weggeworfen) & jeter & verbe à particule \cle{séparable} : \all{Ich werfe den Müll weg.} \\
\hline
wiederverwerten & revaloriser, recycler & \all{wieder} (de nouveau) + \all{verwerten} (valoriser) \\
\hline
die Deponie, die Deponien & la décharge & \all{deponieren} (déposer, mettre en décharge) \\
\hline
schaden (schadet, schadete, hat geschadet) + Datif & nuire à, endommager & \all{Das schadet der Umwelt.} (Cela nuit à l'env.) \\
\hline
der Umweltschutz & la protection de l'environnement & \all{die Umwelt} (l'environnement) + \all{der Schutz} (la protection) \\
\hline
die Verpackungsverordnung, die -en & le règlement sur les emballages & \all{die Verpackung} (l'emballage), \all{verordnen} (ordonner) \\
\hline
der Rohstoff, die Rohstoffe & la matière première & \all{roh} (brut) + \all{der Stoff} (la matière) \\
\hline
\end{tabularx}
\end{center}

\begin{attention}[Piège de la lecture mot à mot]
Ne traduisez pas \all{„Mülltrennung ist Pflicht“} par « La poubelle séparation est devoir ».
\begin{itemize}
    \item \cle{Mülltrennung} = le tri des déchets (nom composé, genre féminin).
    \item \cle{ist Pflicht} = est obligatoire / est une obligation (\all{Pflicht} est un nom, pas un verbe).
    \item \textbf{Stratégie :} Cherchez le verbe principal (\all{ist}), identifiez le sujet (\all{Mülltrennung}), puis le prédicat (\all{Pflicht}). Reformulez en français naturel : « Le tri des déchets est obligatoire. »
\end{itemize}
\end{attention}

\courssub{Vérification des hypothèses — Questions globales (W-Fragen)}

Répondez en allemand (phrases courtes) ou en français après une lecture silencieuse de 3 minutes.

\begin{enumerate}
    \item \all{Worum geht es im Text?} (Quel est le sujet du texte ?)
    \item \all{Wie viele Tonnen hat ein deutscher Haushalt mindestens?} (Combien de poubelles minimum ?)
    \item \all{Was kommt in die gelbe / blaue / braune Tonne?} (Qu'y met-on ?)
    \item \all{Wohin bringt man Glas?} (Où apporte-t-on le verre ?)
    \item \item \all{Was passiert bei falscher Sortierung?} (Qu'arrive-t-il en cas de mauvais tri ?)
    \item \all{Was bedeutet der „Grüne Punkt“?} (Que signifie le « Point vert » ?)
    \item \all{Wie ist die Situation in Kamerun im Vergleich?} (Situation au Cameroun ?)
\end{enumerate}

\courssub{Réformulation orale du contenu général}

\begin{exemplebox}[Modèle de reformulation (niveau B1)]
\textbf{En français :} « Le texte explique que l'Allemagne a un système obligatoire de tri des déchets avec plusieurs poubelles de couleurs (jaune, bleue, brune, grise) et des conteneurs à verre. Ce système permet le recyclage et protège l'environnement. L'auteur suggère que le Cameroun devrait s'en inspirer en améliorant l'éducation environnementale et les infrastructures. »

\textbf{En allemand (Zusammenfassung) :}
\all{Der Text beschreibt das deutsche Mülltrennungssystem. Es gibt vier Tonnen und Glascontainer. Recycling schont Rohstoffe und die Umwelt. Der „Grüne Punkt“ finanziert die Entsorgung. Für Kamerun wäre ein ähnliches System mit Bildung und Gesetzen sinnvoll.}
\end{exemplebox}

\courssub{Le savais-tu ? — L'Allemagne, championne du recyclage}

\begin{savaistu}[Le système du « Grüner Punkt » et les chiffres clés]
L'Allemagne recycle environ \textbf{67 \%} de ses déchets ménagers (\all{Hausmüll}), l'un des meilleurs taux au monde (moyenne UE : env. 48 \%). Depuis \textbf{1991}, la \all{Verpackungsverordnung} (ordonnance sur les emballages) oblige les fabricants à reprendre leurs emballages. Le symbole \textbf{„Grüner Punkt“} (point vert) sur un emballage prouve que l'entreprise a payé une licence à un organisme de récupération (aujourd'hui \all{Der Grüne Punkt – Duales System Deutschland GmbH}). En Autriche et en Suisse, les systèmes sont similaires mais gérés différemment (ex. \all{ARA} en Autriche). Au Cameroun, des initiatives locales (ex. \all{HYSACAM} à Yaoundé/Douala) existent pour la collecte, mais le tri à la source (\all{Quelletrennung}) reste peu développé.
\end{savaistu}

\courssub{Exercice résolu : Analyse globale du texte}

\begin{exoresolu}[Analyse du type et de la structure du texte]
\textbf{Énoncé :} Identifiez le type de texte, le public visé, l'intention de l'auteur et la structure en 3 parties du texte ci-dessus.

\tcblower
\textbf{Solution détaillée :}
\begin{enumerate}
    \item \textbf{Type de texte :} Texte informatif et argumentatif (\all{informierender und appellativer Text}). Il informe sur le système allemand et plaide (\all{appellieren}) pour son adoption au Cameroun.
    \item \textbf{Public visé :} Lycéens / grand public germanophone ou apprenants d'allemand (langue standard, phrases complètes, vocabulaire technique expliqué par le contexte).
    \item \textbf{Intention :} Informer sur le modèle allemand (\all{Vorbildfunktion}) et sensibiliser à la nécessité d'un tri au Cameroun (\all{Handlungsaufforderung} implicite).
    \item \textbf{Structure en 3 parties :}
    \begin{itemize}
        \item \textbf{Partie 1 (L. 1--4) : Description du système.} Énumération des 4 poubelles + conteneurs verre. Vocabulaire technique (\all{Wertstofftonne}, \all{Biomüll}, \all{Restmüll}).
        \item \textbf{Partie 2 (L. 5--9) : Fonctionnement et cadre légal.} Condition de réussite (\all{mitmachen}), sanctions (\all{Bußgelder}), bénéfices écologiques (\all{Rohstoffe}, \all{Umweltschutz}), base légale (\all{Verpackungsverordnung}, \all{Grüner Punkt}).
        \item \textbf{Partie 3 (L. 10--fin) : Comparaison et perspectives.} Autriche/Suisse (modèles similaires) vs Cameroun (problèmes : \all{wilde Deponien}, \all{verbrannt}). Solutions proposées : \all{Umweltbildung}, \all{Gesetze}, \all{Infrastruktur}.
    \end{itemize}
\end{enumerate}
\end{exoresolu}

\courssub{Carte mentale lexicale : Le système de tri allemand}

\begin{popfigure}
\begin{tikzpicture}[
    mindmap,
    grow cyclic,
    every node/.style={concept, execute at begin node=\hskip0pt},
    root concept/.append style={concept color=popblue, minimum size=3.5cm, text width=3cm, font=\bfseries\large},
    level 1 concept/.append style={concept color=popgreen, sibling angle=72, minimum size=2.5cm, text width=2.2cm, font=\bfseries\normalsize, level distance=4.5cm},
    level 2 concept/.append style={concept color=poporange, sibling angle=40, minimum size=1.8cm, text width=1.6cm, font=\small, level distance=3.5cm},
    text=white
]
\node[root concept] {Mülltrennung}
    child[concept color=poporange] {node[level 1 concept, align=center]{Gelbe Tonne\\ (Verpackungen)}
        child {node[level 2 concept] {Plastik}}
        child {node[level 2 concept] {Metall}}
        child {node[level 2 concept] {Verbundstoffe}}
    }
    child[concept color=popblue] {node[level 1 concept, align=center]{Blaue Tonne\\ (Papier/Pappe)}
        child {node[level 2 concept] {Zeitungen}}
        child {node[level 2 concept] {Kartons}}
        child {node[level 2 concept] {Prospekte}}
    }
    child[concept color=popgreen] {node[level 1 concept, align=center]{Braune Tonne\\ (Biomüll)}
        child {node[level 2 concept] {Küchenabfälle}}
        child {node[level 2 concept] {Gartenabfälle}}
        child {node[level 2 concept] {Kompost}}
    }
    child[concept color=poppurple] {node[level 1 concept, align=center]{Graue Tonne\\ (Restmüll)}
        child {node[level 2 concept] {Hygieneartikel}}
        child {node[level 2 concept] {Kehrricht}}
        child {node[level 2 concept] {Deponie}}
    }
    child[concept color=poppink] {node[level 1 concept] {Glascontainer}
        child {node[level 2 concept] {Weißglas}}
        child {node[level 2 concept] {Grünglas}}
        child {node[level 2 concept] {Braunglas}}
    };
\end{tikzpicture}
\legende{Carte mentale du vocabulaire de la Mülltrennung (niveau Tle A).}
\end{popfigure}

\courssub{Synthèse de la lecture globale}

Avant de passer à la compréhension détaillée (Section 2), assurez-vous de maîtriser :
\begin{itemize}
    \item Le \cle{champ lexical} du tri (conteneurs, matériaux, verbes d'action : \all{trennen}, \all{sortieren}, \all{entsorgen}, \all{recyceln}).
    \item La \cle{structure logique} du texte : Description $\rightarrow$ Réglementation/Avantages $\rightarrow$ Transfert/Comparaison.
    \item Les \cle{mots transparents} qui facilitent l'entrée dans le texte : \all{System}, \all{Recycling}, \all{Problem}, \all{Modell}, \all{Container}, \all{Infrastruktur}, \all{Information}, \all{Motivation}.
\end{itemize}

Dans la section suivante, nous exploiterons ce texte pour travailler la \cle{compréhension détaillée} (questions fermées/ouvertes, vrai/faux justifié, recherche d'équivalents) et nous introduirons la première notion de grammaire : \textbf{les verbes à particule séparable} (comme \all{wegwerfen}, \all{mitmachen}, \all{abholen}) au présent et au parfait.

\coursec{Compréhension détaillée}

Dans la section précédente, nous avons lu le texte support \emph{„Mülltrennung – ein Beispiel für Deutschland und Kamerun“} de manière globale : nous avons identifié le thème (la gestion des déchets et le tri des ordures), le type de texte (un article informatif à visée argumentative) et les grandes étapes du raisonnement de l'auteur. Nous passons maintenant à la \cle{lecture détaillée} : il s'agit de vérifier chaque information, de répondre à des questions précises et de repérer les procédés que l'auteur utilise pour convaincre son lecteur.

\courssub{Rappel du texte support}

Pour travailler efficacement, relisons le texte. Les numéros de lignes nous serviront à justifier nos réponses.

\begin{textebox}[„Mülltrennung – ein Beispiel für Deutschland und Kamerun“ (extrait structuré)]
(Zeile 1--3) In Deutschland wird der Müll getrennt: Papier, Glas, Biomüll und Restmüll kommen in verschiedene Tonnen. Für Papier benutzt man die blaue Tonne, für Verpackungen die gelbe Tonne. Glas muss in Container gebracht werden.\\
(Zeile 4--6) Das Recycling ist wichtig, denn aus altem Papier wird neues Papier, und aus Plastikflaschen können sogar Jacken produziert werden. So werden Rohstoffe gespart und die Umwelt wird geschützt. Etwa 65 \% des Mülls werden in Deutschland recycelt.\\
(Zeile 7--10) In Kamerun, zum Beispiel in Yaoundé und Douala, gibt es noch wenige Mülltonnen für die Trennung. Zuerst müssen die Bürger informiert werden. Dann müssen mehr Tonnen aufgestellt werden, und die Abfälle müssen regelmäßig abgeholt werden. Wenn alle mitmachen, wird die Stadt sauberer und gesünder.
\end{textebox}

\textbf{Glossaire :} die Tonne, -n (la poubelle, la benne) ; der Biomüll (les déchets organiques) ; der Restmüll (les déchets résiduels) ; der Container, - (le conteneur) ; das Recycling (le recyclage) ; die Plastikflasche, -n (la bouteille en plastique) ; der Rohstoff, -e (la matière première) ; sparen (économiser) ; schützen (protéger) ; informieren (informer) ; aufstellen (installer) ; abholen (ramasser, venir chercher) ; mitmachen (participer).

\courssub{Méthode : comment répondre à une question de compréhension}

Avant de répondre aux questions, adoptons une démarche rigoureuse. En allemand comme en français, une bonne réponse de compréhension repose sur trois gestes précis.

\begin{methode}[Répondre à une question de compréhension écrite]
\begin{enumerate}
\item \textbf{Repérer} : lis la question et souligne le mot-clé (par exemple \all{Papier}, \all{Recycling}, \all{Kamerun}). Retrouve dans le texte la ligne qui contient ce mot-clé.
\item \textbf{Citer} : recopie entre guillemets allemands „…“ l'élément du texte qui justifie ta réponse (ex. : Zeile 2 : „Für Papier benutzt man die blaue Tonne.“).
\item \textbf{Reformuler} : rédige une phrase complète, avec sujet et verbe conjugué, en reprenant le vocabulaire du texte : \all{Für Papier benutzt man die blaue Tonne.}
\end{enumerate}
\end{methode}

\begin{attention}[Jamais de réponse en un seul mot]
À l'écrit, on ne répond \textbf{jamais} par un mot isolé. À la question \all{Welche Tonne benutzt man für Papier?}, la réponse \all{Blau.} est \emph{fausse} sur le plan de l'expression écrite, même si l'information est juste. Il faut écrire : \all{Man benutzt die blaue Tonne für Papier.} « On utilise la poubelle bleue pour le papier. » De même, pour un vrai/faux, la justification par une citation est obligatoire : écrire seulement \all{Falsch.} ne suffit pas.
\end{attention}

\courssub{Questions de compréhension : vrai ou faux justifié}

Pour chaque affirmation, indique \all{Richtig} ou \all{Falsch} et justifie par une citation du texte avec le numéro de ligne.

\begin{exoresolu}[Vrai / faux justifié]
\textbf{Énoncé.} Entscheiden Sie : richtig oder falsch ? Begründen Sie mit einem Zitat.
\begin{enumerate}
\item Der Müll wird in Deutschland nicht getrennt.
\item Glas kommt in die gelbe Tonne.
\item Aus Plastikflaschen können Jacken produziert werden.
\item In Kamerun gibt es schon viele Tonnen für die Mülltrennung.
\end{enumerate}
\tcblower
\textbf{Solution détaillée.}
\begin{enumerate}
\item \textbf{Falsch.} Zeile 1 : „In Deutschland wird der Müll getrennt.“ — Le texte dit exactement le contraire : les déchets \emph{sont} triés en Allemagne.
\item \textbf{Falsch.} Zeile 2--3 : „Glas muss in Container gebracht werden.“ — Le verre ne va pas dans la poubelle jaune (réservée aux emballages), il doit être apporté dans des conteneurs.
\item \textbf{Richtig.} Zeile 4--5 : „aus Plastikflaschen können sogar Jacken produziert werden.“ — L'information figure mot pour mot dans le texte.
\item \textbf{Falsch.} Zeile 7--8 : „In Kamerun … gibt es noch wenige Mülltonnen für die Trennung.“ — Il y a \emph{encore peu} de poubelles de tri, et non beaucoup.
\end{enumerate}
\end{exoresolu}

Remarque la stratégie du correcteur : les affirmations fausses jouent sur une \textbf{négation ajoutée} (\all{nicht getrennt}), un \textbf{détail substitué} (\all{gelbe Tonne} au lieu de \all{Container}) ou un \textbf{quantificateur inversé} (\all{viele} au lieu de \all{wenige}). Lis donc chaque phrase en la comparant mot à mot avec le texte.

\courssub{Questions ouvertes : réponses rédigées}

Passons aux questions ouvertes. La règle d'or : phrase complète, reprise du vocabulaire du texte, justification si possible.

\textbf{Question 1.} \all{Welche Tonne benutzt man für Papier?} « Quelle poubelle utilise-t-on pour le papier ? »

Réponse rédigée : \all{Für Papier benutzt man die blaue Tonne.} « Pour le papier, on utilise la poubelle bleue. » (Zeile 2)

\textbf{Question 2.} \all{Warum ist Recycling wichtig?} « Pourquoi le recyclage est-il important ? »

Réponse rédigée : \all{Recycling ist wichtig, weil Rohstoffe gespart werden und die Umwelt geschützt wird.} « Le recyclage est important parce que les matières premières sont économisées et que l'environnement est protégé. » (Zeile 5--6)

Observe la transformation : le texte utilise \all{denn} (coordination), la réponse utilise \all{weil} (subordination, verbe à la fin). Les deux sont correctes ; \all{weil} est plus naturel dans une réponse à \all{Warum…?}.

\textbf{Question 3.} \all{Was muss in Kamerun zuerst gemacht werden?} « Que faut-il faire d'abord au Cameroun ? »

Réponse rédigée : \all{Zuerst müssen die Bürger informiert werden.} « D'abord, les citoyens doivent être informés. » (Zeile 8--9)

\begin{exemplebox}[Deux phrases modèles du texte, traduites]
\begin{itemize}
\item \all{Glas muss in Container gebracht werden.} « Le verre doit être apporté dans des conteneurs. »
\item \all{Die Bürger müssen informiert werden.} « Les citoyens doivent être informés. »
\end{itemize}
Ces deux phrases illustrent le passif avec verbe de modalité (\all{muss/müssen + Partizip II + werden}), que nous étudierons en grammaire dans la section 5.
\end{exemplebox}

\courssub{Les procédés de l'auteur pour convaincre}

Le texte n'est pas seulement informatif : il veut \textbf{convaincre} le lecteur que le tri des déchets est nécessaire et possible, en Allemagne comme au Cameroun. Relevons ses procédés.

\begin{center}
\begin{tabularx}{\linewidth}{|X|X|X|}
\hline
\rowcolor{popblueL} Procédé & Exemple dans le texte & Effet sur le lecteur \\
\hline
Exemples concrets & „aus Plastikflaschen können sogar Jacken produziert werden“ & Rendre le recyclage tangible, surprenant \\
\hline
Chiffre & „Etwa 65 \% des Mülls werden recycelt“ & Donner de la crédibilité, prouver l'efficacité \\
\hline
Comparaison Allemagne / Kamerun & „In Kamerun … gibt es noch wenige Mülltonnen“ & Montrer le chemin à parcourir, adapter le message au lecteur camerounais \\
\hline
Étapes ordonnées & „Zuerst … Dann …“ & Présenter un plan d'action clair et réaliste \\
\hline
Appel au lecteur & „Wenn alle mitmachen, wird die Stadt sauberer“ & Impliquer chacun, finir sur une note positive \\
\hline
\end{tabularx}
\end{center}

\begin{popfigure}
\begin{tikzpicture}
\draw[fill=popblueL, draw=popblue, rounded corners] (0,0) rectangle (5.4,8.6);
\draw[fill=poporangeL, draw=poporange, rounded corners] (5.8,0) rectangle (11.2,8.6);
\node[font=\bfseries, text=popblue] at (2.7,8.1) {Question};
\node[font=\bfseries, text=poporange] at (8.5,8.1) {Élément de réponse dans le texte};
\draw[popline] (0,7.7) -- (5.4,7.7);
\draw[popline] (5.8,7.7) -- (11.2,7.7);
\node[align=left, text width=4.8cm, anchor=north west, font=\small] at (0.2,7.5)
 {1. Welche Tonne für Papier ?\\[0.35cm]
  2. Warum ist Recycling wichtig ?\\[0.35cm]
  3. Was muss in Kamerun zuerst gemacht werden ?\\[0.35cm]
  4. Wird der Müll in Deutschland getrennt ?};
\node[align=left, text width=4.8cm, anchor=north west, font=\small] at (6.0,7.5)
 {„Für Papier benutzt man die blaue Tonne.“ (Z. 2)\\[0.35cm]
  „So werden Rohstoffe gespart und die Umwelt wird geschützt.“ (Z. 5--6)\\[0.35cm]
  „Zuerst müssen die Bürger informiert werden.“ (Z. 8--9)\\[0.35cm]
  „In Deutschland wird der Müll getrennt.“ (Z. 1)};
\draw[popmuted, dashed] (0,5.9) -- (5.4,5.9);
\draw[popmuted, dashed] (5.8,5.9) -- (11.2,5.9);
\draw[popmuted, dashed] (0,4.3) -- (5.4,4.3);
\draw[popmuted, dashed] (5.8,4.3) -- (11.2,4.3);
\draw[popmuted, dashed] (0,2.7) -- (5.4,2.7);
\draw[popmuted, dashed] (5.8,2.7) -- (11.2,2.7);
\end{tikzpicture}
\legende{Tableau de correspondance : chaque question renvoie à une citation précise du texte.}
\end{popfigure}

\courssub{Relevé lexical : la famille de \all{Müll}}

Le mot \all{der Müll} (les déchets, les ordures) forme de nombreux composés. Relevons dans le texte tous les mots de cette famille — c'est un excellent exercice de lecture détaillée et d'enrichissement du vocabulaire.

\begin{center}
\begin{tabularx}{\linewidth}{|X|X|X|}
\hline
\rowcolor{popgreenL} Mot composé & Traduction & Ligne du texte \\
\hline
die Mülltrennung & le tri des déchets & titre, Z. 8 \\
\hline
der Biomüll & les déchets organiques & Z. 1 \\
\hline
der Restmüll & les déchets résiduels & Z. 1 \\
\hline
die Mülltonne, -n & la poubelle (de tri) & Z. 8 \\
\hline
der Müll & les déchets, les ordures & Z. 1, 6 \\
\hline
\end{tabularx}
\end{center}

\begin{aretenir}[Le dernier élément commande]
En allemand, dans un nom composé, c'est le \textbf{dernier élément} qui donne le genre et le sens principal : \all{die Tonne} donne \all{die Mülltonne} ; \all{die Trennung} donne \all{die Mülltrennung}. Astuce pour deviner le sens : lis le composé de droite à gauche — \all{Mülltrennung} = « séparation (Trennung) des déchets (Müll) ».
\end{aretenir}

\begin{popfigure}
\begin{tikzpicture}[mindmap, grow cyclic, every node/.style={concept, font=\small},
root concept/.append style={concept color=popgreen, font=\bfseries},
level 1/.append style={level distance=3.4cm, sibling angle=72}]
\node[root concept, text=white]{der Müll}
 child[concept color=popblue] { node[concept, align=center]{die Mülltrennung\\ \tiny le tri} }
 child[concept color=poporange] { node[concept, align=center]{der Biomüll\\ \tiny le compost} }
 child[concept color=poppurple] { node[concept, align=center]{der Restmüll\\ \tiny le résidu} }
 child[concept color=poppink] { node[concept, align=center]{die Mülltonne\\ \tiny la poubelle} }
 child[concept color=popgold] { node[concept, align=center]{der Abfalleimer\\ \tiny la corbeille} };
\end{tikzpicture}
\legende{Carte mentale : la famille du mot \all{Müll}.}
\end{popfigure}

\courssub{Correction collective : la justification par la citation}

Lors de la correction au tableau, chaque élève doit pouvoir défendre sa réponse en trois temps : « Je dis \all{richtig/falsch}, parce que le texte dit à la ligne … : „…“, donc … ». Cette habitude — \textbf{affirmer, citer, conclure} — est exactement celle du commentaire de texte que tu pratiqueras au baccalauréat.

\begin{exercice}[À toi de jouer]
Réponds par des phrases complètes, avec citation justificative :
\begin{enumerate}
\item Wohin kommt der Biomüll ? (cherche une réponse logique à partir du texte)
\item Was wird aus altem Papier ?
\item Nenne zwei Maßnahmen für Kamerun.
\end{enumerate}
\end{exercice}

\begin{corrige}[Exercice : réponses attendues]
\begin{enumerate}
\item \all{Der Biomüll kommt in eine eigene Tonne, denn der Müll wird in verschiedene Tonnen getrennt.} « Les déchets organiques vont dans une poubelle spéciale, car les déchets sont triés dans différentes poubelles. » (Z. 1--2)
\item \all{Aus altem Papier wird neues Papier.} « Avec le vieux papier, on fait du papier neuf. » (Z. 4)
\item \all{Zuerst müssen die Bürger informiert werden, dann müssen mehr Tonnen aufgestellt werden.} « D'abord les citoyens doivent être informés, ensuite davantage de poubelles doivent être installées. » (Z. 8--10)
\end{enumerate}
\end{corrige}

\begin{attention}[Pièges de traduction dans les réponses]
\begin{itemize}
\item \all{die Tonne} n'est pas « la tonne » (poids) mais « la poubelle, la benne » : faux ami classique.
\item \all{zuerst} signifie « d'abord » (première étape), \all{zunächst} est son synonyme ; ne confonds pas avec \all{zuletzt} « enfin, en dernier ».
\item Dans la citation, conserve l'orthographe exacte du texte : majuscules aux noms (\all{Müll}, \all{Tonne}, \all{Bürger}) et ß dans \all{müssen}.
\end{itemize}
\end{attention}

La compréhension détaillée est maintenant acquise : tu sais vérifier une information, rédiger une réponse complète et justifiée, et repérer les procédés de conviction de l'auteur. La section suivante approfondira le lexique thématique du tri des déchets.

\coursec{Lexique thématique : le tri des déchets}

Après notre lecture détaillée du texte support, nous avons relevé de nombreux mots nouveaux : \all{der Müll}, \all{die gelbe Tonne}, \all{trennen}, \all{entsorgen}... Cette section a pour but de les organiser, de les expliquer et de les mémoriser durablement. Un bon lexique ne s'apprend pas mot par mot, mais par \cle{champs lexicaux} et par \cle{familles de mots} : c'est la méthode que nous allons suivre ici.

\courssub{Observer d'abord : un mini-texte de réactivation}

\begin{textebox}[Mülltrennung in Yaoundé — ein neues Projekt]
Seit einem Jahr gibt es in einigen Quartieren von Yaoundé ein neues Projekt: Die Einwohner lernen, ihren Müll zu trennen. In der Schule stehen jetzt drei bunte Tonnen: eine gelbe Tonne für Plastikmüll, eine blaue Tonne für Papier und eine braune Tonne für Biomüll, zum Beispiel Bananenschalen und Gemüsereste. „Früher haben wir alles in einen Abfalleimer geworfen“, erzählt die Schülerin Sandrine. „Heute bringen wir den Restmüll in die schwarze Tonne und das Altglas zum Container neben dem Markt.“ Der Müllwagen kommt jeden Montag und Donnerstag und holt die vollen Tonnen ab. Ein Teil des Mülls wird recycelt: Aus alten Plastikflaschen macht eine kleine Firma in Douala neue Gegenstände. „Wir müssen die Umwelt schützen“, sagt der Direktor der Schule. „Eine saubere Stadt ist gesünder für alle.“ Die Schüler finden das Projekt umweltfreundlich und sinnvoll. Manche bringen sogar von zu Hause getrennten Müll mit. Nur ein Problem bleibt: Nicht alle Familien haben zu Hause genug Eimer, und manche werfen ihren Müll immer noch auf die Straße. Aber die Gewohnheiten ändern sich langsam. „Schmutzig war gestern“, lacht Sandrine, „sauber ist heute!“

\emph{Glossaire :} der Einwohner (–) = l'habitant ; bunt = coloré ; die Bananenschale (-n) = la peau de banane ; der Gemüserest (-e) = le reste de légumes ; früher = autrefois ; abholen = ramasser, venir chercher ; der Gegenstand (die Gegenstände) = l'objet ; sinnvoll = utile, sensé ; die Gewohnheit (-en) = l'habitude.
\end{textebox}

\textbf{Questions de réactivation (réponses orales) :}
\begin{enumerate}
\item \all{Welche Tonnen stehen in der Schule?} (Quelles poubelles se trouvent dans l'école ?)
\item \all{Wohin bringt Sandrine das Altglas?} (Où Sandrine apporte-t-elle le verre usagé ?)
\item \all{Was macht die Firma in Douala?} (Que fait l'entreprise de Douala ?)
\end{enumerate}

Réponses attendues : 1. \all{In der Schule stehen eine gelbe, eine blaue und eine braune Tonne.} 2. \all{Sie bringt das Altglas zum Container neben dem Markt.} 3. \all{Sie macht aus alten Plastikflaschen neue Gegenstände.}

\courssub{Le champ lexical des déchets}

\begin{definition}[die Mülltrennung]
\all{Die Mülltrennung} (–, sans pluriel) désigne l'action de \cle{séparer les déchets selon leur matière} (plastique, papier, verre, matière organique) afin de permettre leur \cle{recyclage}. C'est un mot composé : \all{der Müll} (les déchets) + \all{die Trennung} (la séparation). En allemand, c'est toujours le \textbf{dernier élément} d'un mot composé qui donne le genre : \all{die Trennung} est féminin, donc \all{die Mülltrennung}.
\end{definition}

Observons les mots du texte et classons-les. En allemand, on apprend \textbf{toujours} un nom avec son \cle{article} et son \cle{pluriel} : c'est indispensable pour construire des phrases correctes.

\begin{center}
\begin{tabularx}{\linewidth}{|l|l|X|}
\hline
\rowcolor{popblueL} \textbf{Allemand} & \textbf{Pluriel} & \textbf{Français} \\
\hline
der Müll & (pas de pluriel) & les déchets, les ordures \\
\hline
der Abfall & die Abfälle & le déchet, les ordures \\
\hline
der Restmüll & (pas de pluriel) & les déchets résiduels (non recyclables) \\
\hline
der Biomüll & (pas de pluriel) & les déchets organiques, le biodéchet \\
\hline
der Plastikmüll & (pas de pluriel) & les déchets plastiques \\
\hline
das Altglas & (pas de pluriel) & le verre usagé (à recycler) \\
\hline
das Papier & die Papiere & le papier \\
\hline
die Plastikflasche & die Plastikflaschen & la bouteille en plastique \\
\hline
\end{tabularx}
\end{center}

\begin{attention}[der Müll : un nom sans pluriel]
\all{Der Müll} est un nom \cle{singulier tantum} : il n'a \textbf{pas de pluriel}, comme son équivalent français « les ordures » employé globalement. On ne dit jamais \all{die Mülle}. Pour parler de plusieurs sortes de déchets, on utilise \all{die Abfälle} (pluriel de \all{der Abfall}) ou des composés : \all{der Plastikmüll}, \all{der Biomüll}, \all{der Restmüll}. Même logique pour \all{das Altglas} et \all{der Restmüll}, qui restent au singulier.
\end{attention}

\courssub{Le champ lexical des contenants}

Le texte oppose deux réalités qu'il faut bien distinguer : la petite poubelle d'intérieur et la grande poubelle extérieure.

\begin{center}
\begin{tabularx}{\linewidth}{|l|l|X|}
\hline
\rowcolor{popblueL} \textbf{Allemand} & \textbf{Pluriel} & \textbf{Français} \\
\hline
die Mülltonne & die Mülltonnen & la (grande) poubelle extérieure, la benne \\
\hline
der Abfalleimer & die Abfalleimer & la petite poubelle, la corbeille (intérieur) \\
\hline
der Mülleimer & die Mülleimer & la poubelle, la corbeille (intérieur) \\
\hline
der Container & die Container & le conteneur (verre, vêtements) \\
\hline
die gelbe Tonne & die gelben Tonnen & la poubelle jaune (plastique, emballages) \\
\hline
die blaue Tonne & die blauen Tonnen & la poubelle bleue (papier, carton) \\
\hline
die braune Tonne & die braunen Tonnen & la poubelle brune (biodéchets) \\
\hline
die schwarze Tonne & die schwarzen Tonnen & la poubelle noire/grise (déchets résiduels) \\
\hline
\end{tabularx}
\end{center}

\begin{attention}[Abfalleimer ou Mülltonne ?]
Ne confonds pas \all{der Abfalleimer} (la \textbf{petite} poubelle d'intérieur, dans la cuisine ou la salle de classe) et \all{die Mülltonne} (la \textbf{grande} poubelle extérieure que le camion vient vider). Erreur typique de francophone : traduire « jette cela à la poubelle » (en classe) par \all{Wirf das in die Mülltonne!} — la bonne phrase est \all{Wirf das in den Abfalleimer!} Note aussi le cas : après \all{in} + mouvement, on met l'\textbf{accusatif} : \all{in die Tonne}, \all{in den Eimer} ; après \all{in} + position, le \textbf{datif} : \all{Das Papier ist in der blauen Tonne.}
\end{attention}

\begin{popfigure}
\begin{tikzpicture}[scale=1]
% Poubelle jaune
\draw[fill=popgold!80, draw=popdark] (0,0) rectangle (1.6,2.2);
\draw[fill=popgold, draw=popdark] (-0.15,2.2) rectangle (1.75,2.5);
\node[align=center] at (0.8,1.4) {\textbf{gelb}};
\node[align=center, font=\small] at (0.8,0.6) {Plastik\\Verpackungen};
\node[align=center, font=\footnotesize] at (0.8,-0.4) {plastique,\\emballages};
% Poubelle bleue
\draw[fill=popblue!70, draw=popdark] (2.4,0) rectangle (4.0,2.2);
\draw[fill=popblue, draw=popdark] (2.25,2.2) rectangle (4.15,2.5);
\node[align=center] at (3.2,1.4) {\textbf{blau}};
\node[align=center, font=\small] at (3.2,0.6) {Papier\\Karton};
\node[align=center, font=\footnotesize] at (3.2,-0.4) {papier,\\carton};
% Poubelle brune
\draw[fill=poporange!60!popdark!40, draw=popdark] (4.8,0) rectangle (6.4,2.2);
\draw[fill=poporange!70!popdark!50, draw=popdark] (4.65,2.2) rectangle (6.55,2.5);
\node[align=center] at (5.6,1.4) {\textbf{braun}};
\node[align=center, font=\small] at (5.6,0.6) {Biomüll\\Essensreste};
\node[align=center, font=\footnotesize] at (5.6,-0.4) {biodéchets,\\restes de repas};
% Poubelle noire
\draw[fill=popdark!75, draw=popdark] (7.2,0) rectangle (8.8,2.2);
\draw[fill=popdark, draw=popdark] (7.05,2.2) rectangle (8.95,2.5);
\node[align=center, text=white] at (8.0,1.4) {\textbf{schwarz}};
\node[align=center, font=\small, text=white] at (8.0,0.6) {Restmüll};
\node[align=center, font=\footnotesize] at (8.0,-0.4) {déchets\\résiduels};
\end{tikzpicture}
\legende{Les quatre poubelles du tri sélectif : chaque couleur correspond à une matière.}
\end{popfigure}

\begin{savaistu}[Der Gelbe Sack]
En Allemagne, les emballages en plastique ne vont pas toujours dans une poubelle jaune : dans beaucoup de régions, on utilise le \all{Gelbe Sack} (« le sac jaune »), un sac en plastique jaune distribué gratuitement, que l'on dépose devant la maison le jour de la collecte. Le système s'appelle \all{das Duale System} : les emballages portent souvent un logo appelé \all{der Grüne Punkt} (« le point vert »), qui signifie que l'entreprise paie pour le recyclage de l'emballage.
\end{savaistu}

\courssub{Le champ lexical des actions}

Les verbes sont le moteur du thème. Voici les sept verbes essentiels, avec leurs formes principales (à mémoriser comme pour tout verbe fort) :

\begin{center}
\begin{tabularx}{\linewidth}{|l|l|X|}
\hline
\rowcolor{popblueL} \textbf{Infinitif} & \textbf{Präsens / Perfekt} & \textbf{Français} \\
\hline
trennen & trennt, hat getrennt & trier, séparer \\
\hline
wegwerfen & wirft weg, hat weggeworfen & jeter (au rebut) \\
\hline
entsorgen & entsorgt, hat entsorgt & éliminer, traiter (les déchets) \\
\hline
recyceln & recycelt, hat recycelt & recycler \\
\hline
wiederverwerten & verwertet wieder, hat wiederverwertet & valoriser, réutiliser (recycler) \\
\hline
sammeln & sammelt, hat gesammelt & ramasser, collecter \\
\hline
rausbringen & bringt raus, hat rausgebracht & sortir (la poubelle) \\
\hline
\end{tabularx}
\end{center}

Remarque : \all{rausbringen} est la forme courante et parlée de \all{herausbringen} ; à l'écrit formel, préfère \all{herausbringen} (\all{bringt heraus, hat herausgebracht}). \all{wegwerfen} est un verbe fort irrégulier : \all{ich werfe weg, du wirfst weg, er wirft weg}.

\begin{exemplebox}[Les verbes en contexte]
\begin{itemize}
\item \all{Der Plastikmüll wird in die gelbe Tonne geworfen.} « Les déchets plastiques sont jetés dans la poubelle jaune. »
\item \all{Man muss den Müll entsorgen.} « Il faut éliminer les déchets. »
\item \all{Wir trennen jeden Tag unseren Müll.} « Nous trions nos déchets chaque jour. »
\item \all{Sandrine hat das Altglas zum Container gebracht.} « Sandrine a apporté le verre usagé au conteneur. »
\item \all{Alte Flaschen werden recycelt.} « Les vieilles bouteilles sont recyclées. »
\item \all{Ich bringe heute Abend die Mülltonne raus.} « Je sors la poubelle ce soir. »
\end{itemize}
\end{exemplebox}

\courssub{Le champ lexical des qualités}

Pour décrire une ville, une poubelle ou un comportement, quatre adjectifs sont indispensables :

\begin{center}
\begin{tabularx}{\linewidth}{|X|X|X|}
\hline
\rowcolor{popblueL} \textbf{Français} & \textbf{Deutsch} & \textbf{Remarque} \\
\hline
propre & sauber & contraire : schmutzig \\
\hline
sale & schmutzig & \all{eine schmutzige Straße} \\
\hline
écologique, respectueux de l'environnement & umweltfreundlich & mot composé : Umwelt + freundlich \\
\hline
recyclable, valorisable & wiederverwertbar & suffixe \all{-bar} = « -able » \\
\hline
\end{tabularx}
\end{center}

Exemples : \all{Eine saubere Stadt ist gesünder.} « Une ville propre est plus saine. » ; \all{Dieses Projekt ist sehr umweltfreundlich.} « Ce projet est très écologique. » ; \all{Papier ist wiederverwertbar.} « Le papier est recyclable. »

\courssub{Familles de mots : apprendre par réseaux}

L'allemand adore les \cle{mots composés}. En connaissant un mot de base, tu peux deviner toute une famille :

\begin{itemize}
\item \all{der Müll} $\rightarrow$ \all{die Mülltrennung} (le tri des déchets), \all{die Mülltonne} (la poubelle), \all{der Müllwagen} (le camion poubelle, le camion à ordures), \all{der Mülleimer} (la corbeille).
\item \all{werfen} (lancer) $\rightarrow$ \all{wegwerfen} (jeter au rebut), \all{der Abfall} (ce qui est « jeté de côté » = le déchet).
\item \all{verwerten} (exploiter, valoriser) $\rightarrow$ \all{wiederverwerten} (ré-exploiter = recycler), \all{wiederverwertbar} (recyclable), \all{die Wiederverwertung} (la valorisation, le recyclage).
\item \all{die Umwelt} (l'environnement) $\rightarrow$ \all{umweltfreundlich} (écologique), \all{der Umweltschutz} (la protection de l'environnement), \all{die Umwelt schützen} (protéger l'environnement).
\end{itemize}

\begin{popfigure}
\begin{tikzpicture}[mindmap, grow cyclic, every node/.style={concept, align=center, font=\small},
root concept/.append style={concept color=popblue!60, font=\bfseries},
level 1/.append style={level distance=3.2cm, sibling angle=90},
level 2/.append style={level distance=2.2cm, sibling angle=50}]
\node[root concept]{der Müll}
  child[concept color=popgold!70] { node {die Müll\-tonne} child { node[align=center]{die gelbe\\Tonne} } }
  child[concept color=popgreen!60] { node {die Müll\-trennung} child { node[align=center]{trennen\\recyceln} } }
  child[concept color=poporange!70] { node {der Müll\-wagen} }
  child[concept color=poppurple!50] { node {der Müll\-eimer} child { node {wegwerfen} } };
\end{tikzpicture}
\legende{Carte mentale : la famille du mot \all{der Müll} et les verbes associés.}
\end{popfigure}

\courssub{Expressions utiles à mémoriser par paires}

Retiens ces expressions comme des blocs, avec leur verbe et leur cas :

\begin{center}
\begin{tabularx}{\linewidth}{|X|X|}
\hline
\rowcolor{popblueL} \textbf{Deutsch} & \textbf{Français} \\
\hline
den Müll trennen & trier les déchets \\
\hline
den Müll rausbringen & sortir la poubelle \\
\hline
etwas in die Tonne werfen & jeter quelque chose dans la poubelle \\
\hline
die Umwelt schützen & protéger l'environnement \\
\hline
den Müll entsorgen & éliminer les déchets \\
\hline
Plastik sammeln & collecter le plastique \\
\hline
die Stadt sauber halten & garder la ville propre \\
\hline
\end{tabularx}
\end{center}

\begin{methode}[Mémoriser par paires de mots associés]
\begin{enumerate}
\item Associe toujours un \textbf{verbe} à son \textbf{nom} typique : \all{Müll + trennen}, \all{Tonne + rausbringen}, \all{Umwelt + schützen}.
\item Apprends le nom avec article et pluriel : \all{die Tonne, die Tonnen}.
\item Fabrique une phrase personnelle sur ton quartier : \all{In meinem Quartier bringen wir den Müll am Montag raus.}
\item Récite les paires à voix haute, allemand puis français, puis français puis allemand.
\end{enumerate}
\end{methode}

\begin{exoresolu}[Traduis en allemand]
Énoncé : 1. Les déchets plastiques sont jetés dans la poubelle jaune. 2. Nous devons protéger l'environnement. 3. Jette la peau de banane dans la petite poubelle ! 4. Le camion poubelle vient le lundi.
\tcblower
Solution détaillée :
\begin{enumerate}
\item \all{Der Plastikmüll wird in die gelbe Tonne geworfen.} — passif avec \all{werden} + participe II ; \all{in die Tonne} = accusatif (mouvement).
\item \all{Wir müssen die Umwelt schützen.} — \all{müssen} + infinitif en fin de phrase ; \all{die Umwelt} = accusatif.
\item \all{Wirf die Bananenschale in den Abfalleimer!} — impératif de \all{werfen} (verbe fort : \all{wirf}, sans -st, avec changement de voyelle e $\rightarrow$ i) ; \all{Abfalleimer} car il s'agit d'une petite poubelle d'intérieur ; \all{in den} = accusatif masculin.
\item \all{Der Müllwagen kommt am Montag.} — \all{am Montag} = le lundi (jour habituel).
\end{enumerate}
\end{exoresolu}

\begin{attention}[Pièges fréquents des francophones]
\begin{itemize}
\item Ne dis pas \all{die Mülle} : \all{der Müll} n'a pas de pluriel.
\item \all{entsorgen} n'est pas séparable (préfixe inséparable \all{ent-}) : \all{Man muss den Müll entsorgen} (jamais \all{sorgt ... ent}), et son participe II est \all{entsorgt}, sans \all{ge-}.
\item « Recycler » se dit \all{recyceln} (et non \all{recyclieren} dans la langue courante).
\item Attention au genre : \all{das Altglas} (neutre), \all{der Container} (masculin), \all{die Tonne} (féminin).
\item \all{sauber} se prononce « za-ou-beur » et \all{schmutzig} « chmou-tsik ».
\item Faux ami : \all{die Tonne} n'est pas « la tonne » (poids = \all{die Tonne} aussi, mais le contexte décide) ; dans le thème des déchets, c'est toujours la poubelle.
\end{itemize}
\end{attention}

\begin{savaistu}[Le sac poubelle payant en Suisse]
En Suisse, on ne peut pas jeter ses déchets dans n'importe quel sac : il faut acheter des sacs poubelles officiels, le \all{Gebührensack} (« sac à taxe »), qui sont assez chers. Plus tu produis de déchets, plus tu paies ! Résultat : les habitants trient soigneusement, compostent le \all{Biomüll} et réduisent leurs ordures. C'est un exemple célèbre de politique \all{umweltfreundlich} : le portefeuille aide à protéger \all{die Umwelt} !
\end{savaistu}

\begin{exercice}[À toi de jouer]
1. Donne l'article et le pluriel (s'il existe) de : \all{Tonne}, \all{Abfalleimer}, \all{Müll}, \all{Container}, \all{Altglas}, \all{Plastikflasche}.
2. Complète avec le bon contenant : \all{Die Bananenschale kommt in die \_\_\_ Tonne. Das alte Heft kommt in die \_\_\_ Tonne. Die Plastikflasche kommt in die \_\_\_ Tonne.}
3. Traduis : « Nous trions nos déchets et nous sortons la poubelle le soir. »
\end{exercice}

\begin{corrige}[Exercice « À toi de jouer »]
1. \all{die Tonne, die Tonnen} ; \all{der Abfalleimer, die Abfalleimer} ; \all{der Müll} (sans pluriel) ; \all{der Container, die Container} ; \all{das Altglas} (sans pluriel) ; \all{die Plastikflasche, die Plastikflaschen}.
2. \all{Die Bananenschale kommt in die braune Tonne. Das alte Heft kommt in die blaue Tonne. Die Plastikflasche kommt in die gelbe Tonne.}
3. \all{Wir trennen unseren Müll und bringen die Mülltonne am Abend raus.} (Attention : \all{rausbringen} est séparable — la particule \all{raus} va en fin de phrase.)
\end{corrige}

\begin{aretenir}[L'essentiel du lexique]
Trois familles à connaître par cœur : les \textbf{déchets} (\all{der Müll}, \all{der Abfall, die Abfälle}, \all{das Altglas}), les \textbf{contenants} (\all{die Mülltonne} dehors, \all{der Abfalleimer} dedans, \all{der Container} pour le verre) et les \textbf{actions} (\all{trennen}, \all{wegwerfen}, \all{entsorgen}, \all{recyceln}, \all{rausbringen}). Chaque nom s'apprend avec article et pluriel ; chaque verbe s'apprend avec son complément typique. Ce lexique sera immédiatement réemployé dans la section suivante, consacrée aux verbes à particule séparable comme \all{wegwerfen} et \all{rausbringen}.
\end{aretenir}

\coursec{Grammaire 1 : les verbes à particule séparable}

Dans la section précédente, nous avons constitué le lexique du tri des déchets : \all{der Müll}, \all{die Mülltrennung}, \all{wegwerfen}, \all{rausbringen}, \all{einsammeln}. Vous avez peut-être remarqué que plusieurs de ces verbes se comportent de façon étrange dans les phrases du texte support : tantôt ils apparaissent en un seul mot, tantôt ils sont coupés en deux morceaux placés aux deux extrémités de la phrase. Ce phénomène, typique de l'allemand et redouté des francophones, porte un nom : les \cle{verbes à particule séparable} (\all{trennbare Verben}). C'est l'une des règles les plus importantes de la syntaxe allemande, et elle est au cœur de cette leçon.

\courssub{Observation : un verbe coupé en deux}

Reprenons trois phrases tirées de notre texte sur la gestion des déchets et observons attentivement la place des verbes.

\begin{exemplebox}[Observation dans le texte]
\begin{itemize}
\item \all{Man wirft den Müll weg.} « On jette les déchets. »
\item \all{Sie bringt den Sack zur Sammelstelle mit.} « Elle apporte le sac au point de collecte. »
\item \all{Die Schüler sammeln den Plastikmüll auf dem Schulhof ein.} « Les élèves ramassent les déchets plastiques dans la cour de l'école. »
\item \all{Man muss den Müll wegwerfen.} « Il faut jeter les déchets. »
\item \all{Er hat den Sack weggeworfen.} « Il a jeté le sac. »
\end{itemize}
\end{exemplebox}

Que constate-t-on ?

\begin{enumerate}
\item Dans les trois premières phrases, le verbe est \textbf{coupé en deux} : le radical conjugué (\all{wirft}, \all{bringt}, \all{sammeln}) occupe la deuxième position de la phrase, et la particule (\all{weg}, \all{mit}, \all{ein}) est rejetée \textbf{tout à la fin}.
\item Dans la quatrième phrase, le verbe reste \textbf{en un seul mot} à l'infinitif, en fin de phrase, parce qu'il dépend du verbe modal \all{muss}.
\item Dans la cinquième phrase, au parfait, le verbe est à nouveau \textbf{en un seul mot} : \all{weggeworfen}, mais le \all{-ge-} du participe passé s'est glissé \textbf{entre la particule et le radical}.
\end{enumerate}

La particule n'est donc pas un simple préfixe : c'est un élément mobile, qui obéit à des règles de position précises. C'est elle qui porte le sens : \all{werfen} signifie « lancer », mais \all{wegwerfen} signifie « jeter (pour se débarrasser) » ; \all{bringen} signifie « apporter », \all{mitbringen} signifie « apporter avec soi », \all{rausbringen} signifie « sortir (la poubelle) ».

\begin{propriete}[La règle fondamentale des verbes séparables]
Un verbe à particule séparable se compose d'un \textbf{radical} (verbe de base) et d'une \textbf{particule} (petit mot : \all{ab-, an-, auf-, aus-, ein-, mit-, weg-, raus-}...).
\begin{itemize}
\item Aux \textbf{temps simples} (présent, imparfait) : la particule se \textbf{détache} et va en \textbf{fin de phrase} ; le radical conjugué reste en deuxième position. \all{Ich werfe den Müll weg.}
\item Au \textbf{parfait} (et plus-que-parfait) : le participe passé s'écrit en \textbf{un seul mot}, le \all{-ge-} s'insérant entre la particule et le radical : \all{weg + ge + worfen} $\rightarrow$ \all{weggeworfen}. \all{Ich habe den Müll weggeworfen.}
\item À l'\textbf{infinitif} (avec un verbe modal, \all{zu}, etc.) : le verbe reste \textbf{en un seul mot}, en fin de phrase : \all{Man muss den Müll wegwerfen.}
\item La particule séparable porte toujours \textbf{l'accent tonique} : on prononce « \all{WEGwerfen} », « \all{RAUSbringen} ».
\end{itemize}
\end{propriete}

\begin{popfigure}
\begin{tikzpicture}[
  box/.style={draw=popblue, fill=popblueL, rounded corners=3pt, minimum height=0.9cm, align=center, font=\small},
  partbox/.style={draw=poporange, fill=poporangeL, rounded corners=3pt, minimum height=0.9cm, align=center, font=\small\bfseries}
]
\node[box] (s) at (0,0) {Ich};
\node[box] (v) at (1.8,0) {werfe};
\node[box] (c) at (4.6,0) {den Müll heute};
\node[partbox] (p) at (7.4,0) {weg.};
\draw[-{Stealth[length=3mm]}, very thick, poporange] (v.south) .. controls (3.2,-1.3) and (5.8,-1.3) .. (p.south)
  node[midway, below, font=\small\itshape, text=popdark] {la particule part en fin de phrase};
\node[font=\small, text=popmuted, align=center] at (3.7,1.0) {Position 1 : sujet \quad Position 2 : radical conjugué \quad Fin : particule};
\end{tikzpicture}
\legende{Le verbe \all{wegwerfen} au présent : la particule \all{weg} est renvoyée en fin de phrase.}
\end{popfigure}

\courssub{Particules séparables et particules inséparables}

Comment savoir si un verbe est séparable ? Deux indices : l'accent tonique et la nature de la particule.

\begin{center}
\begin{tabularx}{\linewidth}{|l|X|X|}
\hline
\rowcolor{popblueL} \textbf{Type} & \textbf{Particules} & \textbf{Exemples} \\
\hline
Séparables (accentuées) & \all{ab-, an-, auf-, aus-, ein-, mit-, weg-, raus-, rein-, vor-, zurück-} & \all{wegwerfen, aufstehen, mitmachen, einsammeln, rausbringen} \\
\hline
Inséparables (non accentuées) & \all{be-, ent-, er-, ge-, ver-, zer-} & \all{bekommen, entsorgen, erklären, verstehen, zerstören} \\
\hline
\end{tabularx}
\end{center}

\begin{attention}[Ne pas confondre séparable et inséparable]
Les verbes inséparables ne se coupent \textbf{jamais} et leur participe passé se forme \textbf{sans} \all{-ge-} :
\begin{itemize}
\item \all{verstehen} $\rightarrow$ \all{ich verstehe} $\rightarrow$ \all{ich habe verstanden} (jamais « ich stehe ver », jamais « vergestanden »).
\item \all{entsorgen} $\rightarrow$ \all{ich entsorge den Müll} $\rightarrow$ \all{ich habe den Müll entsorgt}.
\end{itemize}
Astuce : si la particule porte l'accent tonique (« \all{WEGwerfen} »), le verbe est séparable ; si l'accent est sur le radical (« \all{verSTEhen} », « \all{entSORgen} »), il est inséparable. Attention aussi : \all{trennen} (« trier, séparer ») est un verbe de base, sans particule : \all{ich trenne}, \all{ich habe getrennt}.
\end{attention}

\courssub{Conjugaison au présent et au parfait}

Voici la conjugaison complète du verbe séparable \all{wegwerfen}. Remarquez que la particule, en fin de phrase, ne change jamais de forme.

\begin{center}
\begin{tabular}{|l|l|l|}
\hline
\rowcolor{popblueL} \textbf{Personne} & \textbf{Présent} & \textbf{Parfait} \\
\hline
ich & werfe den Müll \textbf{weg} & habe den Müll \textbf{weggeworfen} \\
\hline
du & wirfst den Müll \textbf{weg} & hast den Müll \textbf{weggeworfen} \\
\hline
er/sie/es & wirft den Müll \textbf{weg} & hat den Müll \textbf{weggeworfen} \\
\hline
wir & werfen den Müll \textbf{weg} & haben den Müll \textbf{weggeworfen} \\
\hline
ihr & werft den Müll \textbf{weg} & habt den Müll \textbf{weggeworfen} \\
\hline
sie/Sie & werfen den Müll \textbf{weg} & haben den Müll \textbf{weggeworfen} \\
\hline
\end{tabular}
\end{center}

\begin{center}
\begin{tabular}{|l|l|l|}
\hline
\rowcolor{popgreenL} \textbf{Verbe} & \textbf{Présent (3\textsuperscript{e} pers. du sg.)} & \textbf{Parfait} \\
\hline
\all{rausbringen} (sortir) & er bringt den Müll \textbf{raus} & er hat den Müll \textbf{rausgebracht} \\
\hline
\all{einsammeln} (ramasser) & sie sammelt Papier \textbf{ein} & sie hat Papier \textbf{eingesammelt} \\
\hline
\all{mitmachen} (participer) & wir machen \textbf{mit} & wir haben \textbf{mitgemacht} \\
\hline
\all{mitbringen} (apporter) & er bringt den Sack \textbf{mit} & er hat den Sack \textbf{mitgebracht} \\
\hline
\end{tabular}
\end{center}

\begin{propriete}[Participes passés des verbes séparables]
Le participe passé d'un verbe séparable s'écrit en un seul mot : \textbf{particule + ge + radical (fort) ou particule + ge + radical + t (faible)}.
\begin{center}
\begin{tabular}{|l|l|l|}
\hline
\rowcolor{popblueL} \textbf{Infinitif} & \textbf{Participe passé} & \textbf{Français} \\
\hline
\all{wegwerfen} (fort) & \all{weggeworfen} & jeté \\
\hline
\all{mitbringen} (irrégulier) & \all{mitgebracht} & apporté \\
\hline
\all{rausbringen} (irrégulier) & \all{rausgebracht} & sorti \\
\hline
\all{einsammeln} (faible) & \all{eingesammelt} & ramassé \\
\hline
\all{mitmachen} (faible) & \all{mitgemacht} & participé \\
\hline
\all{aufmachen} (faible) & \all{aufgemacht} & ouvert \\
\hline
\end{tabular}
\end{center}
\end{propriete}

\begin{attention}[L'erreur n° 1 des francophones au parfait]
Le participe passé est \textbf{un seul mot}. N'écrivez jamais « Ich habe geworfen weg » ni « Ich habe weg geworfen » : la seule forme correcte est \all{Ich habe den Sack weggeworfen.} Le \all{-ge-} se colle entre la particule et le radical, sans espace.
\end{attention}

\courssub{L'infinitif avec un verbe modal}

Quand le verbe séparable dépend d'un verbe modal (\all{müssen, können, sollen, wollen, dürfen}), il reste \textbf{en un seul mot} à l'infinitif, placé en fin de phrase. C'est le modal, conjugué, qui occupe la deuxième position.

\begin{exemplebox}[Infinitif après modal]
\begin{itemize}
\item \all{Man muss den Müll wegwerfen.} « Il faut jeter les déchets. »
\item \all{Wir sollen das Glas zur Sammelstelle mitbringen.} « Nous devons apporter le verre au point de collecte. »
\item \all{Alle Schüler können bei der Aktion mitmachen.} « Tous les élèves peuvent participer à l'action. »
\item \all{Ich will heute Abend die Tonne rausbringen.} « Je veux sortir la poubelle ce soir. »
\end{itemize}
\end{exemplebox}

\begin{aretenir}[Et dans une subordonnée ?]
Dans une subordonnée introduite par \all{weil}, \all{dass}, \all{wenn}..., le verbe conjugué va en fin de phrase et le verbe séparable \textbf{se recolle en un seul mot} : \all{Der Lehrer sagt, dass man den Müll wegwerfen muss.} « Le professeur dit qu'il faut jeter les déchets. » ; \all{Sie ist zufrieden, weil alle bei der Aktion mitmachen.} « Elle est contente parce que tout le monde participe à l'action. »
\end{aretenir}

\begin{exemplebox}[Exemples traduits à retenir]
\begin{itemize}
\item \all{Sie bringt den Müll raus.} « Elle sort la poubelle. »
\item \all{Wir haben den Müll getrennt.} « Nous avons trié les déchets. »
\item \all{Er hat den Sack weggeworfen.} « Il a jeté le sac. »
\item \all{Die Kinder sammeln die Flaschen ein.} « Les enfants ramassent les bouteilles. »
\item \all{Machst du bei der Müllaktion mit?} « Est-ce que tu participes à l'action de ramassage ? »
\end{itemize}
\end{exemplebox}

\courssub{Texte d'application : le grand nettoyage du quartier}

\begin{textebox}[„Sauberes Yaoundé“ — eine Aktion im Quartier]
Am Samstagmorgen steht Frau Mbarga früh auf. Sie macht das Fenster auf und sieht die Nachbarn auf der Straße. Heute räumt das ganze Quartier auf! Die Aktion fängt um acht Uhr an. Frau Mbarga zieht ihre Handschuhe an und geht raus.

Zuerst sammeln die Kinder die Plastikflaschen ein. Ein Junge bringt einen großen Sack mit. „Wirf den Müll nicht auf die Straße!“, ruft er seinem kleinen Bruder zu. „Wirf ihn in den Sack!“ Die Frauen trennen den Müll: Papier kommt in die blaue Tonne, Glas in die grüne Tonne und Restmüll in die graue Tonne.

„Man muss den Müll trennen“, erklärt Herr Etoundi, der Lehrer der Nachbarschule. „Dann kann man Papier und Glas recyceln. Der Restmüll wird abgeholt und entsorgt.“ Um elf Uhr bringt ein Lastwagen die vollen Säcke zur Mülldeponie.

Am Nachmittag macht Frau Mbarga die Tür zu und ruht sich aus. Sie ist müde, aber zufrieden. „Wir haben heute viel Müll weggeworfen“, sagt sie, „und das Quartier ist jetzt sauber. Nächsten Monat machen wir wieder mit!“

\textbf{Glossaire :} \all{aufstehen} (se lever) ; \all{aufräumen} (ranger, nettoyer) ; \all{anfangen} (commencer) ; \all{anziehen} (mettre, enfiler) ; \all{die Handschuhe} (les gants) ; \all{zurufen} (crier à qqn) ; \all{die Tonne, -n} (la poubelle) ; \all{der Restmüll} (les déchets résiduels) ; \all{abholen} (ramasser, venir chercher) ; \all{die Mülldeponie, -n} (la décharge) ; \all{sich ausruhen} (se reposer).
\end{textebox}

\textbf{Questions de compréhension :}
\begin{enumerate}
\item \all{Wann fängt die Aktion an?}
\item \all{Was sammeln die Kinder ein?}
\item \all{Warum muss man den Müll trennen?}
\item \all{Wer bringt die Säcke zur Mülldeponie?}
\end{enumerate}

Soulignez dans le texte tous les verbes séparables : vous devriez en trouver plus de dix (\all{steht... auf, macht... auf, räumt... auf, fängt... an, zieht... an, sammeln... ein, bringt... mit, wirf... weg, wird... abgeholt, macht... zu, machen... mit}).

\courssub{Exercice résolu et entraînement}

\begin{exoresolu}[Mettez les verbes séparables à la bonne forme]
Énoncé — Complétez avec le verbe entre parenthèses au présent ou au parfait :
\begin{enumerate}
\item \all{Frau Mbarga \_\_\_\_ (aufstehen) um sechs Uhr \_\_\_\_.}
\item \all{Die Kinder \_\_\_\_ (einsammeln) die Flaschen gestern \_\_\_\_.}
\item \all{\_\_\_\_ (mitmachen) du bei der Aktion \_\_\_\_?}
\item \all{Man \_\_\_\_ (müssen / wegwerfen) den Restmüll \_\_\_\_.}
\end{enumerate}
\tcblower
Solution détaillée :
\begin{enumerate}
\item Présent, temps simple : la particule va en fin de phrase. \all{Frau Mbarga \textbf{steht} um sechs Uhr \textbf{auf}.} « Mme Mbarga se lève à six heures. »
\item Parfait : participe en un seul mot, \all{-ge-} inséré. \all{Die Kinder \textbf{haben} die Flaschen gestern \textbf{eingesammelt}.} « Les enfants ont ramassé les bouteilles hier. »
\item Présent, phrase interrogative sans mot interrogatif : le radical en première position, la particule en fin. \all{\textbf{Machst} du bei der Aktion \textbf{mit}?} « Participes-tu à l'action ? »
\item Avec un modal, l'infinitif reste en un seul mot en fin de phrase. \all{Man \textbf{muss} den Restmüll \textbf{wegwerfen}.} « Il faut jeter les déchets résiduels. »
\end{enumerate}
\end{exoresolu}

\begin{exercice}[À vous de jouer]
\begin{enumerate}
\item Traduisez en allemand : « Elle sort la poubelle le soir. » / « Nous avons participé à l'action. » / « Il faut trier les déchets. »
\item Conjuguez \all{rausbringen} au présent et au parfait à toutes les personnes.
\item Transformez au parfait : \all{Ich werfe den Sack weg. / Sie bringt das Glas mit. / Wir räumen das Quartier auf.}
\end{enumerate}
\end{exercice}

\begin{corrige}[Exercice « À vous de jouer »]
\begin{enumerate}
\item \all{Sie bringt die Tonne abends raus. / Wir haben bei der Aktion mitgemacht. / Man muss den Müll trennen.}
\item Présent : \all{ich bringe raus, du bringst raus, er bringt raus, wir bringen raus, ihr bringt raus, sie bringen raus.} Parfait : \all{ich habe rausgebracht, du hast rausgebracht, er hat rausgebracht, wir haben rausgebracht, ihr habt rausgebracht, sie haben rausgebracht.}
\item \all{Ich habe den Sack weggeworfen. / Sie hat das Glas mitgebracht. / Wir haben das Quartier aufgeräumt.}
\end{enumerate}
\end{corrige}

\begin{aretenir}[L'essentiel sur les verbes séparables]
\begin{itemize}
\item Temps simples : particule \textbf{détachée en fin de phrase} : \all{Ich werfe den Müll weg.}
\item Parfait : \textbf{un seul mot}, \all{-ge-} au milieu : \all{Ich habe den Müll weggeworfen.}
\item Avec modal : infinitif \textbf{en un seul mot} en fin de phrase : \all{Man muss den Müll wegwerfen.}
\item Subordonnée : le verbe conjugué se recolle en fin de phrase : \all{..., weil alle mitmachen.}
\item Séparables (accentuées) : \all{ab-, an-, auf-, aus-, ein-, mit-, weg-, raus-} ; inséparables (sans \all{-ge-}) : \all{be-, ent-, er-, ver-, zer-}.
\end{itemize}
\end{aretenir}

Ces mécanismes maîtrisés, nous pourrons aborder dans la prochaine section une construction qui s'y combine souvent dans les textes écologiques : le passif avec les verbes de modalité, comme dans \all{Der Müll muss getrennt werden}, ainsi que l'impératif des consignes de tri.

\coursec{Grammaire 2 : le passif avec verbes de modalité et l'impératif}

Dans la section précédente, nous avons étudié les verbes à particule séparable (\all{trennen}, \all{wegwerfen}, \all{rausbringen}), indispensables pour parler du tri des déchets. Mais les textes écologiques — affiches, règlements municipaux, campagnes de sensibilisation — emploient deux autres structures grammaticales très fréquentes : le \cle{passif avec verbe de modalité} (\all{Der Müll muss getrennt werden.}) et l'\cle{impératif} (\all{Trennt den Müll!}). C'est l'objet de cette section.

\courssub{Observation : le passif avec modalité dans un texte authentique}

\begin{textebox}[Règlement municipal — „Sauberes Bamenda“ (texte original)]
Liebe Bürgerinnen und Bürger!\\
Ab nächstem Monat wird in unserer Stadt der Müll getrennt. Der Restmüll muss in die schwarze Tonne geworfen werden, Papier und Karton müssen in die blaue Tonne gegeben werden. Altglas kann in den Containern am Marktplatz entsorgt werden. Plastikflaschen sollen gespült und dann in den gelben Sack gelegt werden. Batterien dürfen nicht in den Hausmüll geworfen werden — sie müssen zu den Sammelstellen gebracht werden, denn sie können wiederverwertet werden. Viele Materialien können recycelt werden, aber die Bürger müssen zuerst informiert werden. Deshalb \textbf{sollen Plakate an allen Schulen aufgehängt werden}. Jeder Haushalt soll die neuen Abfalleimer benutzen. Wer die Regeln nicht beachtet, muss eine Strafe zahlen. Helfen Sie mit: Eine saubere Stadt kann nur gemeinsam geschaffen werden!
\end{textebox}

\textbf{Glossaire :} die Tonne, -n (poubelle) ; der Restmüll (déchets résiduels) ; das Altglas (verre usagé) ; der Container, - (conteneur) ; der Marktplatz (place du marché) ; spülen (rincer) ; der gelbe Sack (sac jaune) ; die Batterie, -n (pile) ; die Sammelstelle, -n (point de collecte) ; wiederverwerten (valoriser, réutiliser) ; das Plakat, -e (affiche) ; aufhängen (accrocher) ; der Haushalt, -e (ménage) ; die Strafe, -n (amende) ; beachten (respecter) ; gemeinsam (ensemble).

\begin{exemplebox}[Observation grammaticale]
Repérons les structures en gras dans le texte :

\all{Der Müll \textbf{muss getrennt werden}.} « Les déchets doivent être triés. »

\all{Viele Materialien \textbf{können wiederverwertet werden}.} « Beaucoup de matériaux peuvent être valorisés. »

\all{Batterien \textbf{dürfen nicht} in den Hausmüll \textbf{geworfen werden}.} « Les piles ne doivent pas être jetées dans les ordures ménagères. »

\all{Deshalb \textbf{sollen Plakate ... aufgehängt werden}.} « C'est pourquoi des affiches doivent être accrochées. »

On observe toujours le même schéma : \textbf{sujet + modal conjugué + (compléments) + participe passé + werden}. Le verbe \all{werden} n'est jamais conjugué : il reste à l'infinitif, en toute fin de phrase.
\end{exemplebox}

\courssub{La règle : formation du passif avec verbe de modalité}

\begin{propriete}[Structure du passif avec modalité]
Au passif avec un verbe de modalité, c'est le \textbf{modal} qui porte la conjugaison (personne, temps). Le verbe principal passe au \textbf{participe passé}, et \all{werden} reste à l'\textbf{infinitif en dernière position}, immédiatement précédé du participe passé :

\begin{center}
\textbf{Sujet + modal conjugué + compléments + participe passé + werden}
\end{center}

\all{Der Müll muss heute getrennt werden.} — \all{Die Bürger müssen informiert werden.}
\end{propriete}

\begin{center}
\begin{tabularx}{\linewidth}{|l|X|X|}
\hline
\rowcolor{popblueL} Modal & Phrase allemande & Traduction française \\
\hline
müssen & Der Müll \textbf{muss} getrennt \textbf{werden}. & Les déchets \textbf{doivent être} triés. \\
\hline
können & Das Altglas \textbf{kann} recycelt \textbf{werden}. & Le verre usagé \textbf{peut être} recyclé. \\
\hline
sollen & Die Flaschen \textbf{sollen} gespült \textbf{werden}. & Les bouteilles \textbf{devraient être} rincées. \\
\hline
dürfen & Batterien \textbf{dürfen nicht} weggeworfen \textbf{werden}. & Les piles ne \textbf{doivent pas être} jetées (interdiction). \\
\hline
\end{tabularx}
\end{center}

\begin{aretenir}[Comparaison français / allemand]
Le français dit « doit \textbf{être} trié », « peut \textbf{être} recyclé » : l'auxiliaire « être » est à l'infinitif. L'allemand fait exactement pareil avec \all{werden} : \all{muss getrennt \textbf{werden}}, \all{kann recycelt \textbf{werden}}. Le français utilise aussi souvent « il faut + infinitif » (\all{Man muss den Müll trennen}), mais dès qu'il y a un passif, l'allemand exige \textbf{toujours} l'infinitif \all{werden} en fin de phrase, jamais \all{sein}.
\end{aretenir}

\courssub{Du actif au passif avec modalité}

\begin{methode}[Transformer une phrase active en phrase passive avec modalité]
Prenons : \all{Man muss den Müll trennen.} (« On doit trier les déchets. »)
\begin{enumerate}
\item Le \textbf{complément d'objet direct} de l'actif devient le \textbf{sujet} du passif : \all{den Müll} (accusatif) $\rightarrow$ \all{der Müll} (nominatif).
\item Le \textbf{modal reste conjugué} et s'accorde avec le nouveau sujet : \all{muss}.
\item Le \textbf{verbe principal} devient \textbf{participe passé} : \all{trennen} $\rightarrow$ \all{getrennt}.
\item On ajoute \all{werden} \textbf{à l'infinitif, en toute fin} : \all{Der Müll muss getrennt werden.}
\item Le sujet de l'actif (\all{man}) disparaît généralement ; s'il faut le garder, on utilise \all{von + datif} : \all{Der Müll muss von den Bürgern getrennt werden.}
\end{enumerate}
\end{methode}

\begin{popfigure}
\begin{center}
\begin{tikzpicture}[node distance=1.6cm and 2.2cm]
\node[draw=popblue, fill=popblueL, rounded corners, thick, align=center, text width=5.2cm] (actif) {\textbf{ACTIF}\\ \all{Man muss den Müll trennen.}\\ \footnotesize sujet + modal + COD + infinitif};
\node[draw=popgreen, fill=popgreenL, rounded corners, thick, align=center, text width=5.2cm, right=of actif] (passif) {\textbf{PASSIF}\\ \all{Der Müll muss getrennt werden.}\\ \footnotesize sujet + modal + participe + \all{werden}};
\draw[-{Stealth[length=3mm]}, very thick, poporange] (actif) -- node[above, align=center, font=\footnotesize\bfseries, text=popdark]{Passivierung\\ COD $\rightarrow$ Sujet} (passif);
\node[below=0.9cm of actif, align=center, font=\footnotesize, text=popmuted] {\all{den Müll} (Akk.)};
\node[below=0.9cm of passif, align=center, font=\footnotesize, text=popmuted] {\all{der Müll} (Nom.)};
\draw[-{Stealth[length=2mm]}, dashed, poppurple] ([yshift=-0.35cm]actif.south) to[bend right=15] ([yshift=-0.35cm]passif.south);
\end{tikzpicture}
\legende{Schéma de transformation : de l'actif au passif avec verbe de modalité.}
\end{center}
\end{popfigure}

\begin{exemplebox}[Exemples traduits]
\all{Das Altglas kann recycelt werden.} « Le verre usagé peut être recyclé. »

\all{Die Bürger müssen informiert werden.} « Les citoyens doivent être informés. »

\all{Plastik soll getrennt werden.} « Le plastique devrait être trié. »

\all{Alte Batterien dürfen nicht in die Tonne geworfen werden.} « Les vieilles piles ne doivent pas être jetées dans la poubelle. »

\all{In Douala muss mehr Müll gesammelt werden.} « À Douala, davantage de déchets doivent être ramassés. »
\end{exemplebox}

\begin{attention}[Ne jamais conjuguer \all{werden} après un modal]
Erreur typique du francophone : \all{Der Müll muss getrennt \textbf{wird}.} — FAUX. Après un modal, \all{werden} reste \textbf{toujours à l'infinitif} : \all{Der Müll muss getrennt \textbf{werden}.} De même, ne jamais écrire \all{muss sein getrennt} (calque du français « doit être trié ») : l'allemand n'emploie jamais \all{sein} dans le passif d'action.
\end{attention}

\begin{exoresolu}[Mise au passif avec modalité]
Énoncé : transformez les phrases actives suivantes en phrases passives avec modalité.
a) \all{Man muss die Bürger informieren.} b) \all{Man kann Altglas recyceln.} c) \all{Man soll die Flaschen spülen.}
\tcblower
Solution détaillée :
a) COD \all{die Bürger} $\rightarrow$ sujet nominatif pluriel ; \all{informieren} $\rightarrow$ \all{informiert} ; ajout de \all{werden} en fin : \all{Die Bürger müssen informiert werden.} « Les citoyens doivent être informés. »
b) \all{Altglas} devient sujet (singulier), donc \all{kann} : \all{Altglas kann recycelt werden.} « Le verre usagé peut être recyclé. »
c) \all{Die Flaschen} sujet pluriel, \all{sollen} ; \all{spülen} $\rightarrow$ \all{gespült} : \all{Die Flaschen sollen gespült werden.} « Les bouteilles devraient être rincées. »
\end{exoresolu}

\courssub{L'impératif : donner des consignes écologiques}

Les affiches et campagnes de sensibilisation s'adressent directement au public : elles utilisent l'\cle{impératif}.

\begin{propriete}[Formation de l'impératif]
L'allemand possède trois formes d'impératif :
\begin{itemize}
\item \textbf{2\textsuperscript{e} pers. du singulier (du)} : radical du présent (+ \all{-e} facultatif) : \all{Trenn(e) den Müll!} « Trie les déchets ! »
\item \textbf{2\textsuperscript{e} pers. du pluriel (ihr)} : identique à la 2\textsuperscript{e} pers. plur. du présent, sans pronom : \all{Trennt den Müll!} « Triez les déchets ! »
\item \textbf{Forme de politesse (Sie)} : infinitif + \all{Sie} après le verbe : \all{Trennen Sie den Müll!} « Triez les déchets ! »
\end{itemize}
Le verbe est toujours en \textbf{première position}, suivi du point d'exclamation.
\end{propriete}

\begin{propriete}[Cas particuliers de l'impératif singulier]
\begin{itemize}
\item Verbes à \textbf{changement vocalique e $\rightarrow$ i / ie} : le changement est conservé et le \all{-e} final disparaît : \all{wegwerfen} $\rightarrow$ \all{Wirf weg!} ; \all{nehmen} $\rightarrow$ \all{Nimm!} ; \all{lesen} $\rightarrow$ \all{Lies!}
\item Verbes dont le radical se termine en \textbf{-d, -t, -ig, -tm} : le \all{-e} est \textbf{obligatoire} : \all{warten} $\rightarrow$ \all{Warte!} ; \all{öffnen} $\rightarrow$ \all{Öffne!}
\item Verbes en \textbf{-eln / -ern} : \all{recyceln} $\rightarrow$ \all{Recycl(e) den Müll!} ; \all{sammeln} $\rightarrow$ \all{Sammle das Papier!}
\item \all{sein} est irrégulier : \all{Sei vorsichtig!} « Sois prudent ! » ; \all{Seid leise!} ; \all{Seien Sie freundlich!}
\item \textbf{Verbes séparables} : la particule va en fin de phrase : \all{Bring den Müll raus!} « Sors les poubelles ! » ; \all{Wirf nichts weg!}
\end{itemize}
\end{propriete}

\begin{popfigure}
\begin{center}
\begin{tikzpicture}
\node[draw=popdark, fill=popcream, rounded corners, thick, font=\bfseries] (titre) at (0,0) {L'impératif : trois formes};
\foreach \i/\pers/\tr/\ww/\sn in {0/du/Trenn(e) den Müll!/Wirf weg!/Sei sauber!,
1/ihr/Trennt den Müll!/Werft weg!/Seid sauber!,
2/Sie/Trennen Sie den Müll!/Werfen Sie weg!/Seien Sie sauber!}{
\node[draw=popblue, fill=popblueL, rounded corners, minimum width=1.4cm, font=\bfseries] at (-4.6,-1.1-\i*1.1) {\pers};
\node[draw=popgreen, fill=popgreenL, rounded corners, align=center, minimum width=4.2cm] at (-1.6,-1.1-\i*1.1) {\footnotesize\all{\tr}};
\node[draw=poporange, fill=poporangeL, rounded corners, align=center, minimum width=3.2cm] at (2.3,-1.1-\i*1.1) {\footnotesize\all{\ww}};
\node[draw=poppurple, fill=poppurpleL, rounded corners, align=center, minimum width=3.2cm] at (6.0,-1.1-\i*1.1) {\footnotesize\all{\sn}};
}
\node[font=\footnotesize\bfseries, text=popmuted] at (-1.6,-0.55) {\all{trennen}};
\node[font=\footnotesize\bfseries, text=popmuted] at (2.3,-0.55) {\all{wegwerfen}};
\node[font=\footnotesize\bfseries, text=popmuted] at (6.0,-0.55) {\all{sein}};
\end{tikzpicture}
\legende{Tableau de l'impératif : \all{trennen}, \all{wegwerfen} (séparable, e$\rightarrow$i) et \all{sein} (irrégulier).}
\end{center}
\end{popfigure}

\begin{exemplebox}[Consignes écologiques à l'impératif]
\all{Trenn deinen Müll!} « Trie tes déchets ! »

\all{Werfen Sie nichts auf die Straße!} « Ne jetez rien dans la rue ! »

\all{Schützt die Umwelt!} « Protégez l'environnement ! »

\all{Bring den Müll raus!} « Sors les poubelles ! »

\all{Recycelt das Altglas!} « Recyclez le verre usagé ! »

\all{Spart Wasser und Energie!} « Économisez l'eau et l'énergie ! »
\end{exemplebox}

\begin{attention}[Pièges de l'impératif pour les francophones]
\begin{itemize}
\item Verbes à changement vocalique e $\rightarrow$ i : on dit \all{Wirf!} et non \all{Wirfe!} ; \all{Nimm!} et non \all{Nehme!}. En revanche, les verbes en \all{-d/-t/-ig} gardent le \all{-e} : \all{Warte!}, \all{Öffne!}
\item Ne pas garder le pronom : \all{Trennst du den Müll!} est faux ; correct : \all{Trenn den Müll!} (sauf \all{Sie}, obligatoire : \all{Trennen Sie...}).
\item Verbes séparables : la particule ne reste pas collée au verbe : \all{Bring den Müll raus!} et non \all{Rausbring den Müll!}
\end{itemize}
\end{attention}

\begin{exoresolu}[Affiche de sensibilisation]
Énoncé : rédigez trois slogans à l'impératif (forme \all{ihr}) avec : \all{die Umwelt schützen}, \all{den Müll trennen}, \all{das Papier recyceln}.
\tcblower
Solution : à la forme \all{ihr}, on reprend la 2\textsuperscript{e} personne du pluriel du présent sans le pronom ; la particule des verbes séparables va en fin :
\all{Schützt die Umwelt!} « Protégez l'environnement ! »
\all{Trennt den Müll!} « Triez les déchets ! »
\all{Recycelt das Papier!} « Recyclez le papier ! »
\end{exoresolu}

\begin{exercice}[À vous de jouer]
1. Mettez au passif avec modalité : a) \all{Man muss die Bürger informieren.} b) \all{Man kann Altglas recyceln.} c) \all{Man soll die Flaschen spülen.}
2. Mettez à l'impératif (\all{du}, puis \all{Sie}) : a) \all{wegwerfen / nichts} b) \all{rausbringen / den Müll} c) \all{sein / vorsichtig}.
\end{exercice}

\begin{corrige}[Exercice « À vous de jouer »]
1. a) \all{Die Bürger müssen informiert werden.} b) \all{Altglas kann recycelt werden.} c) \all{Die Flaschen sollen gespült werden.}
2. a) \all{Wirf nichts weg!} / \all{Werfen Sie nichts weg!} b) \all{Bring den Müll raus!} / \all{Bringen Sie den Müll raus!} c) \all{Sei vorsichtig!} / \all{Seien Sie vorsichtig!}
\end{corrige}

\begin{aretenir}[L'essentiel de la section]
Passif avec modalité : \textbf{sujet + modal conjugué + participe passé + werden} (infinitif final, jamais conjugué). Impératif : verbe en première position ; \all{Trenn!} / \all{Trennt!} / \all{Trennen Sie!} ; changement vocalique conservé au singulier (\all{Wirf!}) ; particule séparable en fin (\all{Bring raus!}).
\end{aretenir}

\coursec{Méthode du commentaire et production guidée}

Dans les sections précédentes, nous avons découvert le texte support sur le tri des déchets, enrichi notre vocabulaire (\all{der Müll, die Mülltrennung, das Recycling, entsorgen, sauber}) et appris à construire le passif avec les verbes de modalité (\all{Der Müll muss getrennt werden.}) ainsi qu'à donner des consignes à l'impératif (\all{Trenne deinen Müll!}). Ces outils grammaticaux ne sont pas une fin en soi : ils servent à \emph{parler} et à \emph{écrire} sur l'environnement. C'est exactement ce que l'on vous demandera au baccalauréat : lire un texte, le comprendre, le commenter et produire quelques phrases personnelles. Cette dernière section vous donne la méthode complète, un modèle rédigé et des productions guidées pour vous entraîner.

\courssub{La méthode du commentaire de texte en quatre étapes}

Commenter un texte allemand, ce n'est ni le traduire ni le résumer mot à mot : c'est le \cle{présenter}, le \cle{résumer}, l'\cle{analyser} et donner son \cle{avis}. La méthode se déroule en quatre étapes obligatoires.

\begin{methode}[Rédiger un commentaire de texte (der Kommentar)]
\begin{enumerate}
\item \textbf{Introduction (die Einleitung)} : présentez le texte en deux ou trois phrases — de quoi parle-t-il (thème), d'où vient-il (type de texte : article, lettre, dialogue) — et posez une problématique sous forme de question. Annoncez toujours le thème dès la première phrase : \all{Der Text handelt von der Mülltrennung in Deutschland.}
\item \textbf{Résumé structuré (die Zusammenfassung)} : reformulez les idées principales avec vos propres mots, au présent, sans citer de phrases entières. Suivez l'ordre du texte : \all{Zuerst beschreibt der Autor\dots{} Dann erklärt er\dots}
\item \textbf{Analyse (die Analyse)} : examinez les arguments de l'auteur, le vocabulaire du thème (\all{der Müll, das Recycling, entsorgen}) et les procédés (questions rhétoriques, chiffres, impératifs pour convaincre). Citez de courts extraits du texte pour appuyer chaque remarque.
\item \textbf{Conclusion avec avis personnel (der Schluss)} : terminez toujours par votre opinion justifiée avec \all{weil} : \all{Ich finde die Mülltrennung wichtig, weil sie die Umwelt schützt.}
\end{enumerate}
\end{methode}

\begin{popfigure}
\begin{tikzpicture}[
  bloc/.style={rectangle, rounded corners=4pt, draw=popblue, fill=popblueL, text=popdark, align=center, minimum width=2.6cm, minimum height=1.1cm, font=\small},
  fleche/.style={-{Stealth[length=2.5mm]}, thick, poporange}
]
\node[bloc, align=center] (a) at (0,0){\textbf{1. Introduction}\\Thème + problématique};
\node[bloc, align=center] (b) at (3.6,0){\textbf{2. Résumé}\\Idées principales};
\node[bloc, align=center] (c) at (7.2,0){\textbf{3. Analyse}\\Arguments + citations};
\node[bloc, fill=popgreenL, draw=popgreen, align=center] (d) at (10.8,0){\textbf{4. Conclusion}\\Avis personnel};
\draw[fleche] (a) -- (b);
\draw[fleche] (b) -- (c);
\draw[fleche] (c) -- (d);
\end{tikzpicture}
\legende{Les quatre étapes du commentaire de texte : de la présentation à l'avis personnel.}
\end{popfigure}

\begin{attention}[La subordonnée avec weil et dass : le verbe à la fin]
Dans votre conclusion, vous utiliserez forcément \all{weil} ou \all{dass}. Attention : dans la subordonnée allemande, le verbe conjugué va \textbf{en dernière position}, contrairement au français. Avec un passif + modalité, l'ordre final est : participe + auxiliaire \all{werden} + modalité.

\all{\dots, weil der Müll getrennt werden muss.} « \dots parce que les déchets doivent être triés. »

Ne dites jamais \all{weil der Müll muss getrennt werden} — c'est la faute la plus fréquente des francophones, qui calquent l'ordre français.
\end{attention}

\courssub{Les connecteurs utiles : le squelette du commentaire}

Un bon commentaire est un texte \cle{articulé}. Les connecteurs montrent au correcteur que votre pensée est organisée. Mémorisez cette boîte à outils :

\begin{center}
\begin{tabularx}{\linewidth}{|l|X|X|}\hline
\rowcolor{popblueL} \textbf{Connecteur} & \textbf{Français} & \textbf{Fonction} \\ \hline
\all{zuerst} & d'abord, tout d'abord & ouvrir le résumé \\ \hline
\all{dann} & ensuite, puis & enchaîner les idées \\ \hline
\all{außerdem} & de plus, en outre & ajouter un argument \\ \hline
\all{deshalb} & c'est pourquoi, donc & conséquence \\ \hline
\all{meiner Meinung nach} & à mon avis & introduire l'avis personnel \\ \hline
\all{zum Schluss} & pour conclure, enfin & ouvrir la conclusion \\ \hline
\end{tabularx}
\end{center}

\begin{attention}[Place du verbe après les connecteurs]
\all{Zuerst}, \all{dann}, \all{außerdem}, \all{deshalb} et \all{meiner Meinung nach} occupent la \textbf{première place} de la phrase : le verbe conjugué vient donc \textbf{juste après}, avant le sujet (inversion) : \all{Deshalb müssen wir den Müll trennen.} (et non \all{Deshalb wir müssen\dots}). En revanche, \all{zum Schluss kann man sagen, dass\dots} introduit une subordonnée avec verbe final.
\end{attention}

\begin{exemplebox}[Phrases-modèles à réutiliser]
\all{Meiner Meinung nach sollte die Mülltrennung auch in Kamerun eingeführt werden.} « À mon avis, le tri des déchets devrait aussi être introduit au Cameroun. » (\all{sollte} est le Präteritum de \all{sollen}, employé ici avec une valeur de conditionnel pour nuancer : « devrait ».)

\all{Zum Schluss kann man sagen, dass Recycling die Umwelt schützt.} « Pour conclure, on peut dire que le recyclage protège l'environnement. »
\end{exemplebox}

\courssub{Commentaire modèle rédigé}

Voici un commentaire complet (10-12 lignes) sur le texte support de la leçon, avec la problématique imposée. Observez comment les quatre étapes et les connecteurs s'enchaînent, et comment la grammaire de la leçon (verbes séparables, passif + modalité, impératif) est réutilisée.

\begin{textebox}[Kommentar : Ist die Mülltrennung auch in Kamerun möglich ?]
Der Text handelt von der Mülltrennung in Deutschland und stellt die Frage : Ist die Mülltrennung auch in Kamerun möglich ? Zuerst beschreibt der Autor, wie die Deutschen ihren Müll in verschiedene Tonnen werfen : Papier, Glas, Plastik und Biomüll. Dann erklärt er, dass der Müll in Deutschland recycelt werden muss und dass das System viel Geld kostet. Außerdem zeigt er, dass viele Bürger am Anfang keine Lust hatten, aber heute fast alle mitmachen. Der Autor benutzt viele Imperative wie „Trennen Sie Ihren Müll!“, um die Leser zu überzeugen. Deshalb wirkt der Text wie ein Appell an die Bevölkerung. Meiner Meinung nach sollte die Mülltrennung auch in Kamerun eingeführt werden, weil unsere Städte wie Yaoundé und Douala immer schmutziger werden. Zum Schluss kann man sagen, dass Recycling die Umwelt schützt und dass jeder Bürger anfangen muss, seinen Müll zu trennen.
\end{textebox}

\textbf{Glossaire} : die Tonne, -n (la poubelle, le conteneur) ; der Biomüll (les déchets organiques) ; Lust haben (avoir envie) ; mitmachen (participer, séparable) ; der Appell, -e (l'appel) ; überzeugen (convaincre) ; schmutzig (sale) ; anfangen (commencer, séparable).

\begin{aretenir}[Ce qui fait la qualité de ce commentaire]
\begin{itemize}
\item La problématique est posée dès l'introduction sous forme de question.
\item Les connecteurs \all{zuerst}, \all{dann}, \all{außerdem}, \all{deshalb}, \all{meiner Meinung nach}, \all{zum Schluss} structurent le texte.
\item L'analyse cite le texte („Trennen Sie Ihren Müll!“) et identifie un procédé (l'impératif pour convaincre).
\item La conclusion donne un avis personnel justifié par \all{weil} + verbe final, avec un passif + modalité (\all{eingeführt werden sollte}).
\end{itemize}
\end{aretenir}

\courssub{Production guidée 1 : le résumé en 4-5 phrases}

Résumer, c'est garder les idées et changer les mots. Interdiction de recopier des phrases entières du texte : utilisez des synonymes et des verbes de la leçon. Le résumé se rédige au \textbf{présent} et ne contient \textbf{aucune opinion personnelle}.

\begin{exoresolu}[Résumez le texte support en 4-5 phrases avec vos propres mots]
Consigne : \all{Fassen Sie den Text mit eigenen Worten zusammen (4-5 Sätze).} Rédigez le résumé d'un texte qui explique que les Allemands trient leurs déchets dans des poubelles de couleurs, que le recyclage est obligatoire et coûteux, mais que les citoyens participent aujourd'hui massivement.
\tcblower
\textbf{Solution détaillée} — modèle de résumé :

\all{In Deutschland wirft man den Müll in verschiedene Abfalleimer : gelb für Plastik, blau für Papier und braun für Biomüll. Das Recycling ist Pflicht, aber es ist teuer. Am Anfang wollten viele Leute nicht mitmachen. Heute trennen fast alle Bürger ihren Müll, weil sie die Umwelt schützen wollen.}

Méthode : chaque phrase reprend une idée du texte avec des mots différents (\all{Pflicht} au lieu de \all{muss recycelt werden} ; \all{fast alle Bürger} au lieu de \all{die Mehrheit}). Le résumé reste au présent et ne contient aucune opinion personnelle — l'avis est réservé à la conclusion du commentaire.
\end{exoresolu}

\courssub{Production guidée 2 : « Der Müll in meiner Stadt »}

Vous devez maintenant rédiger un court paragraphe (8-10 lignes) sur les déchets dans votre ville (Yaoundé, Douala, Bafoussam\dots). La consigne impose la grammaire de la leçon : \textbf{2 verbes séparables}, \textbf{1 passif avec modalité} et \textbf{1 impératif}.

Avant d'écrire, rappelez-vous les trois formes de l'impératif avec un verbe séparable — la particule va en fin de phrase :

\begin{center}
\begin{tabular}{|l|l|l|}\hline
\rowcolor{popblueL} \textbf{Personne} & \textbf{Forme} & \textbf{Exemple} \\ \hline
\all{du} & radical (+ e) & \all{Trenne deinen Müll!} \\ \hline
\all{ihr} & radical + t & \all{Trennt euren Müll!} \\ \hline
\all{Sie} & infinitif + Sie & \all{Trennen Sie Ihren Müll!} \\ \hline
\end{tabular}
\end{center}

\begin{methode}[Construire le paragraphe]
\begin{enumerate}
\item Décrivez la situation (2-3 phrases) : \all{In meiner Stadt liegt viel Müll auf den Straßen.}
\item Insérez deux verbes séparables, par exemple \all{anfangen}, \all{mitmachen} ou \all{wegwerfen} : \all{Viele Leute werfen den Müll einfach weg.}
\item Ajoutez un passif + modalité : \all{Der Plastikmüll muss endlich recycelt werden.}
\item Terminez par un impératif adressé au lecteur : \all{Hilf mit und trenne deinen Müll!}
\end{enumerate}
\end{methode}

\begin{exemplebox}[Modèle de paragraphe]
\all{In Douala gibt es viel Müll in den Straßen, besonders nach dem Regen. Viele Menschen werfen ihre Plastikflaschen einfach weg, und das Wasser kann nicht abfließen. Meiner Meinung nach fängt die Lösung in der Schule an : die Kinder müssen lernen, den Müll zu trennen. Außerdem sollte in jedem Quartier ein Abfalleimer aufgestellt werden. Wenn alle mitmachen, wird unsere Stadt sauberer. Deshalb : Trenne deinen Müll und schütze deine Umwelt!}
\end{exemplebox}

Vérifiez : \all{wegwerfen} et \all{mitmachen} (2 verbes séparables, particule en fin de phrase) ; \all{aufgestellt werden sollte} (passif + modalité) ; \all{Trenne\dots{} schütze\dots} (impératif). Le contrat est rempli.

\begin{exercice}[À vous d'écrire]
Rédigez votre propre paragraphe \all{Der Müll in meiner Stadt} (8-10 lignes) en respectant le contrat : 2 verbes séparables soulignés, 1 passif avec modalité, 1 impératif, et au moins 3 connecteurs de la liste (\all{zuerst, außerdem, deshalb, zum Schluss}).
\end{exercice}

\begin{corrige}[Exercice : Der Müll in meiner Stadt]
Éléments de correction attendus. Exemple de production acceptable :

\all{In Yaoundé sehe ich jeden Tag Müll vor den Häusern und an den Märkten. Viele Bürger werfen Papier und Plastik einfach weg. Zuerst muss die Stadt mehr Abfalleimer kaufen, dann muss der Müll regelmäßig abgeholt werden. Außerdem sollte Plastik recycelt werden, weil es die Flüsse verschmutzt. Wenn die Schüler früh anfangen, den Müll zu trennen, wird das Problem kleiner. Zum Schluss : Mach mit und halte deine Stadt sauber!}

Barème indicatif : 2 verbes séparables correctement placés (\all{werfen\dots weg}, \all{abgeholt werden}, \all{anfangen}, \all{Mach mit}) ; passif + modalité correct (\all{recycelt werden sollte} / \all{abgeholt werden muss}) ; impératif correct (\all{Mach}, \all{halte}) ; connecteurs présents ; vocabulaire du thème (\all{der Müll, der Abfalleimer, recyceln, verschmutzen, sauber}).
\end{corrige}

\courssub{Production orale : mini-exposé par binômes}

\begin{experience}[Yaoundé ou Douala contre une ville allemande]
Par binômes, préparez un mini-exposé de 3 minutes comparant la gestion des déchets dans votre ville et dans une ville allemande (par exemple Freiburg ou Berlin). Répartition : l'élève A présente la ville camerounaise, l'élève B la ville allemande, puis les deux concluent ensemble.

Plan imposé :
\begin{enumerate}
\item \all{In Douala sieht man\dots{} / In Freiburg gibt es\dots}
\item Comparaison : \all{In Deutschland wird der Müll getrennt, aber bei uns\dots}
\item Proposition : \all{Meiner Meinung nach sollte\dots{} eingeführt werden, weil\dots}
\end{enumerate}

Le reste de la classe écoute et pose une question à chaque binôme : \all{Was kann man in deiner Stadt besser machen ?}
\end{experience}

\courssub{Grille d'auto-évaluation}

Avant de rendre votre production écrite, évaluez-vous honnêtement avec cette grille — c'est exactement celle du correcteur.

\begin{center}
\begin{tabularx}{\linewidth}{|X|c|c|}\hline
\rowcolor{popblueL} \textbf{Critère} & \textbf{Oui} & \textbf{À revoir} \\ \hline
J'ai utilisé au moins 5 mots du thème (\all{der Müll, die Mülltrennung, das Recycling, entsorgen, sauber}) & $\square$ & $\square$ \\ \hline
Mes verbes séparables sont correctement placés (particule en fin) & $\square$ & $\square$ \\ \hline
Mon passif + modalité est correct (\all{muss getrennt werden}) & $\square$ & $\square$ \\ \hline
Mon impératif est bien formé (\all{Trenne!} / \all{Trennen Sie!}) & $\square$ & $\square$ \\ \hline
Après \all{weil/dass}, mon verbe est à la fin de la phrase & $\square$ & $\square$ \\ \hline
Mon texte a une structure claire (introduction, développement, conclusion) & $\square$ & $\square$ \\ \hline
J'ai donné un avis personnel justifié (\all{Ich finde\dots, weil\dots}) & $\square$ & $\square$ \\ \hline
Tous les noms ont une majuscule et un article correct & $\square$ & $\square$ \\ \hline
\end{tabularx}
\end{center}

\begin{savaistu}[La production écrite au baccalauréat]
Au Bac, la production écrite est évaluée sur trois critères principaux : la richesse du vocabulaire, la correction grammaticale et la cohérence du texte. Bonne nouvelle : réutiliser les structures de la leçon rapporte directement des points. Un passif avec modalité bien construit (\all{Der Müll muss getrennt werden.}), un verbe séparable bien placé et un avis personnel avec \all{weil} + verbe final montrent au correcteur que vous maîtrisez le programme de Terminale. N'inventez pas de structures compliquées : recyclez celles du cours — c'est aussi cela, le recyclage !
\end{savaistu}

\begin{aretenir}[L'essentiel de la section]
\begin{itemize}
\item Un commentaire suit toujours quatre étapes : introduction (thème + problématique), résumé, analyse avec citations, conclusion avec avis personnel justifié.
\item Les connecteurs \all{zuerst, dann, außerdem, deshalb, meiner Meinung nach, zum Schluss} structurent le texte et provoquent l'inversion du sujet.
\item Après \all{weil} et \all{dass}, le verbe conjugué va en fin de phrase : \all{\dots, weil der Müll getrennt werden muss.}
\item L'impératif d'un verbe séparable rejette la particule en fin de phrase : \all{Trenne deinen Müll!} / \all{Trennen Sie Ihren Müll!}
\item Dans toute production guidée, respectez le contrat grammatical imposé et vérifiez-vous avec la grille d'auto-évaluation.
\end{itemize}
\end{aretenir}

\newpage

\coursec{Activité d'intégration}

\courssub{Objectifs de l'activité}

À l'issue de cette activité d'intégration, l'élève sera capable de :
\begin{itemize}
\item mobiliser le \cle{vocabulaire thématique} du tri des déchets (\all{der Müll, die Mülltrennung, das Recycling, die Tonne, entsorgen, sauber}\ldots) dans une tâche de communication authentique ;
\item employer correctement les \cle{verbes à particule séparable} au présent et au parfait (\all{wegwerfen, aufsammeln, mitmachen, aufräumen}\ldots) ;
\item construire des phrases au \cle{passif avec verbes de modalité} (\all{Plastik muss recycelt werden.}) ;
\item utiliser l'\cle{impératif} pour formuler des slogans et des consignes (\all{Trennt euren Müll !}) ;
\item produire un court texte cohérent de 6 à 8 lignes et le présenter à l'oral en 2 minutes ;
\item travailler en groupe, évaluer la production des autres et justifier un choix.
\end{itemize}

\courssub{Situation}

Le quartier Nkol-Eton, à Yaoundé, souffre d'un problème croissant : les ordures s'accumulent près du marché et des caniveaux, surtout pendant la saison des pluies. L'association des jeunes du quartier, en partenariat avec le club d'allemand du lycée, lance une campagne de sensibilisation intitulée \all{„Sauberes Viertel“} (« Un quartier propre »). Votre groupe a été choisi pour concevoir l'affiche officielle de la campagne, en allemand, car une délégation d'étudiants allemands de la ville jumelée visitera le lycée le mois prochain.

Voici le mot d'information publié par l'association :

\begin{textebox}[Appel à participation de l'association des jeunes]
Liebe Schülerinnen und Schüler,

unser Viertel hat ein Problem: Überall liegt Müll auf den Straßen. Plastikflaschen werden einfach weggeworfen, Papier und Reste werden nicht getrennt. Das muss sich ändern! Deshalb starten wir die Kampagne „Sauberes Viertel“. Wir suchen die besten Plakate auf Deutsch. Euer Plakat braucht: einen Slogan, drei klare Regeln und einen kurzen Text über unser Projekt. Die beste Gruppe präsentiert ihr Plakat vor der deutschen Delegation.

Macht mit — für ein sauberes Viertel!

\emph{Traduction : Chers élèves, notre quartier a un problème : il y a des ordures partout dans les rues. Les bouteilles en plastique sont simplement jetées, le papier et les restes ne sont pas triés. Cela doit changer ! C'est pourquoi nous lançons la campagne « Un quartier propre ». Nous cherchons les meilleures affiches en allemand. Votre affiche doit contenir : un slogan, trois règles claires et un court texte sur notre projet. Le meilleur groupe présentera son affiche devant la délégation allemande. Participez — pour un quartier propre !}
\end{textebox}

\begin{definition}[Lexique du texte à retenir]
\begin{itemize}
\item \all{das Viertel, -} : le quartier ; \all{die Straße, -n} : la rue.
\item \all{die Plastikflasche, -n} : la bouteille en plastique ; \all{der Rest, -e} : le reste, les déchets alimentaires.
\item \all{wegwerfen (wirft weg, warf weg, hat weggeworfen)} : jeter (verbe fort séparable).
\item \all{trennen (trennte, hat getrennt)} : trier, séparer — attention, verbe \textbf{non séparable}.
\item \all{das Plakat, -e} : l'affiche ; \all{der Slogan, -s} : le slogan.
\item \all{mitmachen (machte mit, hat mitgemacht)} : participer (verbe séparable).
\item \all{sich ändern} : changer (verbe pronominal) ; \all{die Delegation, -en} : la délégation.
\end{itemize}
\end{definition}

\courssub{Consignes}

Par groupes de 4, réalisez les tâches suivantes, dans l'ordre :

\begin{enumerate}
\item \textbf{Compréhension.} Lisez le mot de l'association et répondez par écrit, en allemand et en phrases complètes : \all{Was ist das Problem im Viertel ? Was sucht der Verein ? Wer besucht bald die Schule ?} (« Quel est le problème dans le quartier ? Que cherche l'association ? Qui visitera bientôt l'établissement ? »)
\item \textbf{Lexique.} Complétez en groupe une liste de dix mots utiles pour votre affiche : cinq noms avec article et pluriel (ex. \all{die Flasche, -n}), cinq verbes (dont au moins trois verbes séparables).
\item \textbf{Slogan.} Rédigez un slogan court et frappant à l'\cle{impératif} (forme \all{ihr}, car l'affiche s'adresse à tous les habitants). Exemple : \all{Trennt euren Müll !}
\item \textbf{Règles.} Formulez trois consignes au \cle{passif avec verbe de modalité} : une avec \all{müssen}, une avec \all{sollen}, une avec \all{können}. Exemple : \all{Plastik muss in die gelbe Tonne geworfen werden.}
\item \textbf{Texte de présentation.} Rédigez ensemble un texte de 6 à 8 lignes qui présente le projet : le problème, la solution, l'appel à participer. Employez \textbf{au moins trois verbes séparables} (au présent ou au parfait) et deux noms du lexique du tri.
\item \textbf{Affiche.} Disposez sur une feuille : le slogan en grand, les trois règles avec un symbole ou un dessin simple, le texte en bas.
\item \textbf{Présentation orale.} Chaque groupe présente son affiche en 2 minutes : un élève lit le slogan et les règles, un autre lit le texte, les deux autres répondent aux questions de la classe.
\item \textbf{Vote et évaluation.} La classe vote pour la meilleure campagne selon la grille : vocabulaire (5 points), grammaire (5 points), prononciation (5 points), créativité (5 points).
\end{enumerate}

\begin{attention}[Pièges à éviter dans votre production]
\begin{itemize}
\item Verbe séparable : la particule va \textbf{en fin de phrase} au présent : \all{Wir werfen den Müll nicht weg.} et non \all{Wir wegwerfen den Müll.}
\item Au parfait : \all{ge-} s'insère entre particule et radical : \all{Wir haben den Müll weggeworfen.}
\item Passif avec modalité : l'infinitif \all{werden} se place \textbf{tout à la fin} : \all{Papier muss getrennt werden.}
\item Impératif \all{ihr} : radical seul, sans pronom : \all{Trennt den Müll !} (et non \all{Trennt ihr den Müll !}).
\item Attention : \all{trennen} n'est \textbf{pas} un verbe séparable ; à l'impératif, rien ne se détache : \all{Trennt den Müll !} En revanche, \all{mitmachen} est séparable : \all{Macht mit !}
\item N'oubliez pas la majuscule à tous les noms : \all{der Müll, die Tonne, das Viertel}.
\end{itemize}
\end{attention}

\courssub{Ressources à mobiliser}

\begin{aretenir}[Boîte à outils linguistique]
\textbf{Lexique :} \all{der Müll (Sg.), der Abfall, die Abfälle, die Mülltrennung (Sg.), das Recycling (Sg.), die Tonne, -n, der Abfalleimer, -, die Flasche, -n, das Papier (Sg.), der Rest, -e, der Kompost (Sg.), sauber, schmutzig, wegwerfen, aufsammeln, trennen, entsorgen, recyceln, mitmachen, aufräumen}.

\textbf{Verbe séparable au présent (\all{wegwerfen}) :} \all{ich werfe weg, du wirfst weg, er/sie/es wirft weg, wir werfen weg, ihr werft weg, sie/Sie werfen weg.}

\textbf{Passif avec modalité :} sujet + modal conjugué + participe II + \all{werden} : \all{Der Müll muss entsorgt werden.} / « Les déchets doivent être éliminés. »

\textbf{Impératif :} \all{Trenn den Müll !} (du) — \all{Trennt den Müll !} (ihr) — \all{Trennen Sie den Müll !} (Sie).

\textbf{Connecteurs utiles :} \all{deshalb, außerdem, zuerst, dann, endlich, für, ohne}.
\end{aretenir}

\courssub{Proposition de production et corrigé commenté}

\begin{exoresolu}[Modèle d'affiche et de texte — groupe type]
\textbf{Énoncé.} Produisez le contenu complet d'une affiche « Sauberes Viertel » : slogan, trois règles, texte de 6 à 8 lignes.
\tcblower
\textbf{Solution détaillée — production modèle :}

\emph{Slogan (impératif, forme ihr) :}

\all{Trennt euren Müll — für ein sauberes Viertel !}

\emph{Trois règles (passif avec modalité) :}

\all{1. Plastikflaschen müssen in die gelbe Tonne geworfen werden.}

\all{2. Papier und Glas sollen getrennt gesammelt werden.}

\all{3. Küchenabfälle können zu Kompost verarbeitet werden.}

\emph{Texte de présentation :}

\all{In unserem Viertel Nkol-Eton liegt viel Müll auf den Straßen. Das ist schmutzig und gefährlich für unsere Gesundheit. Deshalb räumen wir jeden Samstag gemeinsam auf. Wir sammeln Plastikflaschen auf und werfen sie nicht einfach weg. Papier wird recycelt, und aus Resten wird Kompost gemacht. Viele Nachbarn machen schon mit. Letzte Woche haben wir zwanzig Säcke Müll aufgesammelt. Macht auch mit — gemeinsam sind wir stark !}

\emph{Traduction : Dans notre quartier Nkol-Eton, il y a beaucoup d'ordures dans les rues. C'est sale et dangereux pour notre santé. C'est pourquoi nous nettoyons ensemble chaque samedi. Nous ramassons les bouteilles en plastique et ne les jetons pas simplement. Le papier est recyclé, et à partir des restes on fait du compost. Beaucoup de voisins participent déjà. La semaine dernière, nous avons ramassé vingt sacs d'ordures. Participez aussi — ensemble nous sommes forts !}

\textbf{Commentaire des choix de langue :}
\begin{itemize}
\item \all{Trennt} : impératif \all{ihr} du verbe \all{trennen} (verbe simple, non séparable : rien ne se détache). La forme \all{ihr} interpelle tous les habitants sans pronom exprimé.
\item \all{müssen\ldots geworfen werden}, \all{sollen\ldots gesammelt werden}, \all{können\ldots verarbeitet werden} : passif avec modalité ; le modal est conjugué en deuxième position, le participe II précède \all{werden}, qui ferme la phrase.
\item Verbes séparables au présent : \all{räumen\ldots auf}, \all{sammeln\ldots auf}, \all{werfen\ldots weg}, \all{machen\ldots mit} — la particule est renvoyée en fin de phrase.
\item Verbe séparable au parfait : \all{haben wir\ldots aufgesammelt} — \all{ge-} inséré entre la particule et le radical (\all{auf + ge + sammelt}).
\item Passif simple au présent : \all{Papier wird recycelt}, \all{aus Resten wird Kompost gemacht} — sujet patient au nominatif, \all{werden} conjugué + participe II en finale.
\item Lexique du thème : \all{der Müll, die Tonne, recyceln, der Kompost, sauber, schmutzig}.
\item Connecteurs : \all{deshalb, und, letzte Woche} assurent la cohérence du texte.
\end{itemize}
\end{exoresolu}

\courssub{Bilan de l'activité}

Cette activité d'intégration a permis de réinvestir l'ensemble des savoir-faire de la leçon dans une situation de communication réelle et motivante, ancrée dans le quotidien d'un quartier de Yaoundé. Sur le plan de la \cle{compréhension}, les élèves ont lu et exploité un document authentique (l'appel à participation). Sur le plan du \cle{lexique}, ils ont sélectionné et réemployé le vocabulaire du tri des déchets avec articles et pluriels. Sur le plan de la \cle{grammaire}, trois acquis majeurs ont été mobilisés de manière fonctionnelle : les verbes à particule séparable (au présent et au parfait), le passif avec verbes de modalité pour formuler des règles impersonnelles, et l'impératif pour interpeller le public. Enfin, la présentation orale et le vote ont développé l'\cle{expression orale}, l'écoute active et l'évaluation entre pairs.

\begin{methode}[Pour réussir ce type de tâche d'intégration]
\begin{enumerate}
\item Lire deux fois le document déclencheur et souligner les mots-clés.
\item Construire d'abord une banque de mots et de structures avant de rédiger.
\item Vérifier chaque phrase : place du verbe, particule séparable en finale, \all{werden} à la fin au passif.
\item Lire le texte à voix haute dans le groupe pour contrôler la prononciation.
\item Se répartir clairement les rôles pour la présentation orale.
\end{enumerate}
\end{methode}

L'enseignant pourra prolonger l'activité en affichant les meilleures productions dans la salle de classe et en proposant, en devoir, la traduction (version) de trois phrases de l'affiche vers le français, puis le thème de deux consignes simples vers l'allemand, afin de consolider le passage entre les deux langues et de préparer les épreuves du Baccalauréat.

\newpage

\coursec{Méthodes et exercices résolus}

\courssub{Méthode 1 : Comprendre un texte sur la gestion des déchets}

\begin{methode}[Lire et exploiter un texte sur \all{die Mülltrennung}]
\begin{enumerate}
\item \textbf{Première lecture (globale).} Lis le texte une première fois sans dictionnaire. Repère : le thème (\all{der Müll}, \all{die Mülltrennung}, \all{das Recycling}), le lieu (Allemagne ? Cameroun ?), le ton (informatif, polémique ?) et le type de texte (article, reportage, interview).
\item \textbf{Repérage des mots-clés du champ lexical.} Souligne les mots de la famille : \all{der Abfall, die Abfälle} ; \all{der Mülleimer, -} ; \all{die Tonne, -n} ; \all{entsorgen} ; \all{recyceln} ; \all{sauber} ; \all{die Umwelt}. Ces mots te guident vers les réponses.
\item \textbf{Lecture détaillée.} Pour chaque question, localise le passage du texte qui contient la réponse (les questions suivent presque toujours l'ordre du texte). Ne recopie jamais une phrase entière : reformule avec tes mots.
\item \textbf{Rédaction de la réponse.} Réponds par une phrase complète en allemand, verbe conjugué en deuxième position, en reprenant le temps de la question.
\item \textbf{Vérification.} Relis : majuscule aux noms, place du verbe, accord sujet-verbe, cas après les prépositions.
\end{enumerate}
\textbf{Structures à retenir :} \all{Der Text handelt von ...} (le texte traite de ...) ; \all{Dem Text zufolge ...} (selon le texte ...) ; \all{Man muss den Müll trennen.} (il faut trier les déchets).
\end{methode}

\begin{exoresolu}[Compréhension de texte — type Bac]
\textbf{Texte :} \all{In Deutschland trennen die meisten Familien ihren Müll: Papier kommt in die blaue Tonne, Plastik in die gelbe Tonne und Biomüll in die braune Tonne. So können viele Abfälle recycelt werden. In Yaoundé dagegen werfen viele Menschen alles in denselben Eimer. Die Stadt will jetzt das Recycling ausbauen und die Bürger informieren.}

\textbf{Questions :} 1. \all{Was trennen die deutschen Familien?} 2. \all{In welche Tonne kommt das Papier?} 3. \all{Was ist das Problem in Yaoundé?} 4. \all{Was will die Stadt Yaoundé machen?}
\tcblower
\textbf{Raisonnement :} les questions suivent l'ordre du texte. Q1 → première phrase ; Q2 → deuxième phrase ; Q3 → troisième phrase (\all{dagegen} = « en revanche », signal du contraste) ; Q4 → dernière phrase.

\textbf{Réponses rédigées :}
\begin{enumerate}
\item \all{Die deutschen Familien trennen ihren Müll.} (Les familles allemandes trient leurs déchets.) — Verbe séparable \all{trennen} : particule en fin de phrase.
\item \all{Das Papier kommt in die blaue Tonne.} (Le papier va dans la poubelle bleue.) — \all{kommen in} + accusatif (mouvement).
\item \all{Das Problem in Yaoundé ist, dass viele Menschen alles in denselben Eimer werfen.} (Le problème à Yaoundé est que beaucoup de gens jettent tout dans la même poubelle.) — Subordonnée avec \all{dass} : verbe conjugué à la fin.
\item \all{Die Stadt will das Recycling ausbauen und die Bürger informieren.} (La ville veut développer le recyclage et informer les citoyens.) — \all{wollen} + infinitif en fin de phrase ; \all{ausbauen} est séparable mais reste entier à l'infinitif.
\end{enumerate}
\textbf{Conclusion :} chaque réponse est une phrase complète, avec le verbe conjugué à la deuxième place et la particule séparable renvoyée en fin de phrase principale.
\end{exoresolu}

\courssub{Méthode 2 : Employer les verbes à particule séparable}

\begin{methode}[Maîtriser \all{wegwerfen}, \all{aufräumen}, \all{mitmachen}...]
\begin{enumerate}
\item \textbf{Identifier la particule.} Si l'infinitif s'écrit en un mot avec un préfixe \textbf{accentué} (\all{WEGwerfen}, \all{AUFräumen}, \all{ANfangen}), le verbe est séparable.
\item \textbf{Phrase principale au présent.} Le verbe conjugué prend la 2\textsuperscript{e} place, la particule va à la fin : \all{Ich werfe den Müll weg.} (Je jette les déchets.)
\item \textbf{Au Perfekt.} Le \all{ge-} se place entre la particule et le radical : \all{weggeworfen}, \all{aufgeräumt}. Auxiliaire \all{haben} (ou \all{sein} si mouvement) : \all{Wir haben den Hof aufgeräumt.}
\item \textbf{Avec un modal ou un infinitif.} Le verbe reste entier à la fin : \all{Man muss den Müll wegwerfen.}
\item \textbf{Subordonnée.} Le verbe se recolle à la fin : \all{..., weil ich den Müll wegwerfe.}
\end{enumerate}
\textbf{À retenir :} particule \textbf{accentuée} = séparable ; particule non accentuée (\all{be-, ver-, er-, ent-}) = inséparable (\all{entsorgen} → \all{entsorgt}, sans \all{ge-}).
\end{methode}

\begin{center}
\begin{tabularx}{\linewidth}{|X|X|X|}
\hline
\rowcolor{popblueL} Infinitif & Präsens (phrase principale) & Perfekt \\
\hline
\all{wegwerfen} (jeter) & \all{Ich werfe den Müll weg.} & \all{Ich habe den Müll weggeworfen.} \\
\hline
\all{aufräumen} (ranger, nettoyer) & \all{Wir räumen den Hof auf.} & \all{Wir haben den Hof aufgeräumt.} \\
\hline
\all{mitmachen} (participer) & \all{Sie macht bei der Aktion mit.} & \all{Sie hat bei der Aktion mitgemacht.} \\
\hline
\all{einsammeln} (ramasser) & \all{Er sammelt die Flaschen ein.} & \all{Er hat die Flaschen eingesammelt.} \\
\hline
\all{entsorgen} (inséparable !) & \all{Man entsorgt den Müll.} & \all{Man hat den Müll entsorgt.} \\
\hline
\end{tabularx}
\end{center}

\begin{exoresolu}[Verbes séparables — transformation de phrases]
\textbf{Énoncé :} Mets les phrases suivantes au Präsens, puis au Perfekt : 1. \all{wir / aufräumen / der Hof} ; 2. \all{sie / wegwerfen / die Flaschen} ; 3. \all{die Schüler / mitmachen / bei der Aktion} ; 4. Traduis : « Vous devez trier les déchets. » (\all{der Müll, trennen, müssen})
\tcblower
\textbf{Raisonnement :} \all{aufräumen}, \all{wegwerfen}, \all{mitmachen}, \all{trennen} sont séparables (particule accentuée). Au présent : verbe en 2\textsuperscript{e} position + particule finale. Au Perfekt : participe passé avec \all{-ge-} inséré entre particule et radical.

\textbf{Solutions :}
\begin{enumerate}
\item Präsens : \all{Wir räumen den Hof auf.} — Perfekt : \all{Wir haben den Hof aufgeräumt.} (Nous avons nettoyé la cour.) — \all{der Hof} → accusatif \all{den Hof}.
\item Präsens : \all{Sie wirft die Flaschen weg.} — Perfekt : \all{Sie hat die Flaschen weggeworfen.} (Elle a jeté les bouteilles.) — verbe fort : \all{werfen – warf – geworfen} ; 3\textsuperscript{e} pers. \all{wirft} (changement \all{e} → \all{i}).
\item Präsens : \all{Die Schüler machen bei der Aktion mit.} — Perfekt : \all{Die Schüler haben bei der Aktion mitgemacht.} (Les élèves ont participé à l'action.) — \all{bei} + datif : \all{bei der Aktion}.
\item \all{Sie müssen den Müll trennen.} — Avec \all{müssen}, l'infinitif reste entier (non séparé) à la fin de la phrase ; \all{der Müll} → accusatif \all{den Müll}.
\end{enumerate}
\textbf{Conclusion :} la règle d'or : séparation en phrase principale simple, verbe recollé à la fin avec modal ou dans une subordonnée.
\end{exoresolu}

\courssub{Méthode 3 : Le passif avec verbes de modalité}

\begin{methode}[Construire \all{Der Müll muss getrennt werden}]
\begin{enumerate}
\item \textbf{Structure.} Sujet (patient) + modal conjugué + participe passé + \all{werden} à l'infinitif, en fin de phrase : \all{Die Flaschen können recycelt werden.} (Les bouteilles peuvent être recyclées.)
\item \textbf{Passage actif → passif.} L'objet accusatif de l'actif devient le sujet nominatif du passif ; le sujet actif disparaît ou devient \all{von} + datif : \all{Man trennt den Müll.} → \all{Der Müll wird getrennt.}
\item \textbf{Accord.} Le modal s'accorde avec le nouveau sujet : \all{Die Abfälle müssen entsorgt werden.} (pluriel → \all{müssen}).
\item \textbf{Ordre des mots en fin de phrase.} Modal conjugué en 2\textsuperscript{e} position, puis : participe passé AVANT \all{werden} : \all{... recycelt werden}, jamais l'inverse.
\item \textbf{Emploi.} Ce passif est idéal pour les consignes, règlements et textes sur l'environnement : \all{Plastik darf nicht in die Natur geworfen werden.} (Le plastique ne doit pas être jeté dans la nature.)
\end{enumerate}
\end{methode}

\begin{exoresolu}[Passif avec modal — thème grammatical]
\textbf{Énoncé :} Traduis en allemand : 1. « Les déchets doivent être triés. » 2. « Le papier peut être recyclé. » 3. « Le plastique ne doit pas être jeté dans la rivière. » 4. Transforme au passif avec \all{müssen} : \all{Man sammelt die Flaschen ein.}
\tcblower
\textbf{Raisonnement :} « doivent être triés » = \all{müssen} + participe + \all{werden}. Attention : \all{der Müll} est singulier en allemand (pas de pluriel pour « les déchets » au sens général ; on peut aussi dire \all{die Abfälle}).

\textbf{Solutions :}
\begin{enumerate}
\item \all{Der Müll muss getrennt werden.} (ou \all{Die Abfälle müssen getrennt werden.}) — \all{der Müll} singulier → \all{muss} ; participe \all{getrennt} avant \all{werden}.
\item \all{Das Papier kann recycelt werden.} — \all{das Papier} neutre, singulier → \all{kann}.
\item \all{Das Plastik darf nicht in den Fluss geworfen werden.} — négation \all{nicht} après le modal ; \all{in} + accusatif (mouvement) : \all{in den Fluss} ; participe du verbe fort \all{werfen} : \all{geworfen}.
\item Actif : \all{Man sammelt die Flaschen ein.} → Passif : \all{Die Flaschen müssen eingesammelt werden.} (Les bouteilles doivent être ramassées.) — \all{die Flaschen} pluriel → \all{müssen} ; verbe séparable \all{einsammeln} → participe \all{eingesammelt}.
\end{enumerate}
\textbf{Conclusion :} la structure figée \all{modal + participe + werden} est un classique du Bac : vérifie l'accord du modal et la place de \all{werden} en toute fin.
\end{exoresolu}

\courssub{Méthode 4 : L'impératif pour les consignes écologiques}

\begin{methode}[Former l'impératif (tu, vous, nous)]
\begin{enumerate}
\item \textbf{2\textsuperscript{e} pers. du singulier (\all{du}).} Radical du présent, sans pronom : \all{Trenn(e) den Müll!} Les verbes en \all{e → i/ie} gardent le changement : \all{Wirf den Abfall weg!} (mais pas les verbes en \all{a → ä} : \all{Fahr langsam!}).
\item \textbf{2\textsuperscript{e} pers. du pluriel (\all{ihr}).} Comme le présent sans pronom : \all{Trennt den Müll!} ; \all{Werft die Flaschen weg!}
\item \textbf{Forme de politesse (\all{Sie}).} Comme le présent, pronom APRÈS le verbe : \all{Trennen Sie den Müll!} ; \all{Werfen Sie nichts auf den Boden!}
\item \textbf{1\textsuperscript{re} pers. du pluriel.} \all{Trennen wir den Müll!} (Trions les déchets !)
\item \textbf{Verbes séparables.} La particule va à la fin : \all{Räum den Hof auf!} ; \all{Macht bei der Aktion mit!}
\end{enumerate}
\textbf{Usage :} affiches, campagnes de sensibilisation, slogans : \all{Schützt die Umwelt!} (Protégez l'environnement !)
\end{methode}

\begin{center}
\begin{tabular}{|l|l|l|l|}
\hline
\rowcolor{popblueL} Infinitif & \all{du} & \all{ihr} & \all{Sie} \\
\hline
\all{trennen} & \all{Trenn(e)!} & \all{Trennt!} & \all{Trennen Sie!} \\
\hline
\all{wegwerfen} & \all{Wirf ... weg!} & \all{Werft ... weg!} & \all{Werfen Sie ... weg!} \\
\hline
\all{aufräumen} & \all{Räum ... auf!} & \all{Räumt ... auf!} & \all{Räumen Sie ... auf!} \\
\hline
\all{sauber halten} & \all{Halt(e) ... sauber!} & \all{Haltet ... sauber!} & \all{Halten Sie ... sauber!} \\
\hline
\end{tabular}
\end{center}

\begin{exoresolu}[Impératif — rédaction de slogans]
\textbf{Énoncé :} Tu rédiges une affiche pour la campagne « Une école propre » de ton lycée à Douala. Mets les verbes à l'impératif (\all{ihr}-Form) : 1. \all{wegwerfen / der Müll / in den Mülleimer} ; 2. \all{aufräumen / der Klassenraum / jeden Tag} ; 3. \all{mitmachen / bei der Müllaktion} ; 4. \all{sauber halten / die Schule} (forme \all{Sie}).
\tcblower
\textbf{Raisonnement :} forme \all{ihr} = présent sans pronom ; verbes forts avec \all{e → i} : \all{werfen → werft} ; \all{halten → haltet} (à la 2\textsuperscript{e} pers. pluriel, le radical est celui du présent : \all{ihr haltet}). Particules séparables en fin de phrase.

\textbf{Solutions :}
\begin{enumerate}
\item \all{Werft den Müll in den Mülleimer!} (Jetez les déchets dans la poubelle !) — \all{in} + accusatif (mouvement vers) : \all{in den Mülleimer}.
\item \all{Räumt den Klassenraum jeden Tag auf!} (Nettoyez la salle de classe chaque jour !) — particule \all{auf} renvoyée à la fin ; \all{jeden Tag} = accusatif de temps.
\item \all{Macht bei der Müllaktion mit!} (Participez à l'action de tri !) — \all{bei} + datif : \all{bei der Müllaktion}.
\item \all{Halten Sie die Schule sauber!} (Gardez l'école propre !) — forme de politesse : pronom \all{Sie} avec majuscule, placé après le verbe.
\end{enumerate}
\textbf{Conclusion :} à l'impératif, le verbe ouvre la phrase, la particule séparable la ferme : c'est la structure typique du slogan.
\end{exoresolu}

\courssub{Méthode 5 : Traduire (thème) un paragraphe sur l'environnement}

\begin{methode}[Thème : du français à l'allemand sans calque]
\begin{enumerate}
\item \textbf{Analyser la phrase française.} Repère sujet, verbe, compléments, temps. Demande-toi : y a-t-il un verbe séparable ? un modal ? un passif ?
\item \textbf{Poser le squelette allemand.} Verbe conjugué en 2\textsuperscript{e} position ; si deux verbes (modal, Perfekt, passif), le deuxième va en fin de phrase.
\item \textbf{Placer les compléments.} Ordre : temps avant lieu (\all{Ich bringe heute den Müll zur Tonne.}) ; complément d'objet direct à l'accusatif.
\item \textbf{Choisir le bon vocabulaire.} « jeter » = \all{wegwerfen} (à la poubelle) ; « trier » = \all{trennen} ; « propre » = \all{sauber} ; « sale » = \all{schmutzig}. Méfie-toi des faux amis : « ordures » ≠ \all{Ordnung} !
\item \textbf{Relire.} Majuscules aux noms, Umlaute, accord du verbe, particule séparable finale.
\end{enumerate}
\end{methode}

\begin{center}
\begin{tabularx}{\linewidth}{|X|X|X|}
\hline
\rowcolor{popblueL} Français & Deutsch & Remarque \\
\hline
les déchets, les ordures & \all{der Müll} (sing.), \all{die Abfälle} (pl.) & jamais \all{die Ordnung} (= l'ordre) ! \\
\hline
la poubelle & \all{die Tonne, -n} ; \all{der Mülleimer, -} & \all{in die Tonne} (acc., mouvement) \\
\hline
trier & \all{trennen} (séparable) & \all{die Mülltrennung} = le tri des déchets \\
\hline
jeter & \all{wegwerfen} (séparable, fort) & \all{wirft weg – warf weg – weggeworfen} \\
\hline
recycler & \all{recyceln} & \all{das Recycling} \\
\hline
propre / sale & \all{sauber} / \all{schmutzig} & \all{sauber halten} = garder propre \\
\hline
l'environnement & \all{die Umwelt} & toujours féminin \\
\hline
\end{tabularx}
\end{center}

\begin{exoresolu}[Thème — phrases à traduire]
\textbf{Énoncé :} Traduis en allemand : 1. « Chaque semaine, nous jetons les vieux journaux dans la poubelle bleue. » 2. « Dans notre quartier, les déchets doivent être ramassés le matin. » 3. « Hier, les élèves ont nettoyé la cour de l'école. »
\tcblower
\textbf{Raisonnement :} phrase 1 : complément de temps en tête → inversion sujet/verbe ; « jeter » = \all{wegwerfen} (séparable). Phrase 2 : passif + modal. Phrase 3 : Perfekt avec verbe séparable \all{aufräumen}.

\textbf{Solutions :}
\begin{enumerate}
\item \all{Jede Woche werfen wir die alten Zeitungen in die blaue Tonne.} — « chaque semaine » : \all{jede Woche} (accusatif de temps, \all{die Woche} féminin) ; le temps en tête provoque l'inversion : \all{werfen wir} ; \all{in} + accusatif (mouvement).
\item \all{In unserem Viertel müssen die Abfälle am Morgen eingesammelt werden.} — passif avec modal : \all{müssen ... eingesammelt werden} ; \all{unserem Viertel} : datif après \all{in} (situation) ; \all{das Viertel} = le quartier.
\item \all{Gestern haben die Schüler den Schulhof aufgeräumt.} — Perfekt : \all{haben ... aufgeräumt} ; « la cour de l'école » = mot composé \all{der Schulhof} (l'allemand préfère la composition au génitif).
\end{enumerate}
\textbf{Conclusion :} trois pièges évités : l'inversion après un complément initial, la double finale du passif avec modal, et la particule séparable au Perfekt.
\end{exoresolu}

\courssub{Méthode 6 : Produire un court paragraphe argumenté}

\begin{methode}[Rédiger 5 à 8 phrases sur le tri des déchets]
\begin{enumerate}
\item \textbf{Plan simple.} Introduction (le problème) → 2 ou 3 arguments (santé, environnement, économie) → proposition de solutions → conclusion personnelle.
\item \textbf{Connecteurs à mobiliser.} \all{zuerst} (d'abord), \all{dann} (ensuite), \all{außerdem} (de plus), \all{denn / weil} (car), \all{deshalb} (c'est pourquoi), \all{meiner Meinung nach} (à mon avis).
\item \textbf{Grammaire de la leçon à réutiliser.} Au moins un verbe séparable, un passif avec modal et un impératif : cela montre ta maîtrise du programme.
\item \textbf{Phrases modèles.} \all{In unserer Stadt gibt es viel Müll.} / \all{Der Müll muss getrennt werden.} / \all{Deshalb müssen wir alle mitmachen.}
\item \textbf{Relecture systématique.} Verbe en 2\textsuperscript{e} position, majuscules, cas après prépositions (\all{mit dem Bus}, \all{in der Stadt}).
\end{enumerate}
\end{methode}

\begin{exoresolu}[Production écrite guidée — type Bac]
\textbf{Énoncé :} \all{Schreiben Sie einen kurzen Text (6-8 Sätze): „Mülltrennung in unserer Schule — warum ist das wichtig?“} (Rédige un court texte de 6 à 8 phrases : « Le tri des déchets dans notre école — pourquoi est-ce important ? »)
\tcblower
\textbf{Raisonnement :} suivre le plan : problème → arguments → solutions → conclusion. Intégrer volontairement les points de grammaire de la leçon : verbes séparables (\all{wegwerfen}, \all{mitmachen}), passif avec modal (\all{muss getrennt werden}), impératif final.

\textbf{Production modèle :}

\all{In unserer Schule in Yaoundé gibt es jeden Tag viel Müll: Plastikflaschen, Papier und Essensreste. Viele Schüler werfen alles in denselben Eimer, und das ist ein Problem für die Umwelt und die Gesundheit. Dabei kann Papier recycelt werden, und Biomüll kann als Kompost benutzt werden. Deshalb muss der Müll in unserer Schule getrennt werden. Wir brauchen drei Tonnen: eine für Papier, eine für Plastik und eine für Biomüll. Außerdem müssen alle Schüler bei dieser Aktion mitmachen. Meiner Meinung nach ist eine saubere Schule wichtig für alle. Also: Trennt den Müll und haltet die Schule sauber!}

\textbf{Traduction :} « Dans notre école à Yaoundé, il y a chaque jour beaucoup de déchets : bouteilles en plastique, papier et restes de nourriture. Beaucoup d'élèves jettent tout dans la même poubelle, et c'est un problème pour l'environnement et la santé. Pourtant, le papier peut être recyclé et les biodéchets peuvent être utilisés comme compost. C'est pourquoi les déchets doivent être triés dans notre école. Nous avons besoin de trois poubelles : une pour le papier, une pour le plastique et une pour les biodéchets. De plus, tous les élèves doivent participer à cette action. À mon avis, une école propre est importante pour tous. Alors : triez les déchets et gardez l'école propre ! »

\textbf{Points grammaticaux exploités :} verbes séparables (\all{mitmachen}, \all{wegwerfen} dans \all{werfen alles ... weg} au sens implicite), passif avec modal (\all{kann recycelt werden}, \all{kann ... benutzt werden}, \all{muss ... getrennt werden}), impératif final (\all{Trennt}, \all{haltet}), connecteurs variés (\all{deshalb}, \all{außerdem}, \all{meiner Meinung nach}).
\end{exoresolu}

\begin{savaistu}[Le tri des déchets en Allemagne : une culture nationale]
En Allemagne, le tri des déchets fait partie de la vie quotidienne : chaque couleur de poubelle a sa fonction — la \all{blaue Tonne} pour le papier, la \all{gelbe Tonne} (ou le \all{Gelbe Sack}, le sac jaune) pour les emballages en plastique, la \all{braune Tonne} pour les biodéchets et la \all{schwarze Tonne} pour le reste. Les bouteilles en verre se trient par couleur dans des conteneurs (\all{der Glascontainer}) : blanc, vert, brun. Beaucoup de bouteilles en plastique sont consignées : c'est le système du \all{Pfand} — on paie une petite somme en plus à l'achat, remboursée quand on rapporte la bouteille au magasin. Au Cameroun, des initiatives de recyclage voient aussi le jour à Yaoundé et à Douala : le vocabulaire de cette leçon te permet d'en parler en allemand !
\end{savaistu}

\begin{attention}[Les 5 erreurs qui coûtent des points au Bac]
\begin{enumerate}
\item \textbf{Oublier de séparer la particule.} Faux : \all{Ich wegwerfe den Müll.} Juste : \all{Ich werfe den Müll weg.} En phrase principale, la particule va TOUJOURS à la fin.
\item \textbf{Inverser participe et \all{werden} au passif avec modal.} Faux : \all{Der Müll muss werden getrennt.} Juste : \all{Der Müll muss getrennt werden.} L'ordre figé est : modal + participe + \all{werden}.
\item \textbf{Calquer le pluriel français « les déchets ».} \all{Der Müll} est singulier : \all{Der Müll muss entsorgt werden.} Si tu veux un pluriel, utilise \all{die Abfälle}.
\item \textbf{Faux amis et confusion de mots.} « ordures » = \all{der Müll / die Abfälle}, jamais \all{die Ordnung} (qui signifie « l'ordre » !) ; « propre » = \all{sauber} ; « jeter » = \all{wegwerfen}, pas \all{schmeißen} dans un texte soigné (trop familier).
\item \textbf{Oublier l'inversion et les majuscules.} Après un complément en tête de phrase, le verbe reste en 2\textsuperscript{e} position : \all{Jeden Tag werfen wir den Müll weg} (et non \all{Jeden Tag wir werfen...}). Et majuscule à TOUS les noms : \all{die Tonne, das Papier, der Mülleimer}.
\end{enumerate}
\end{attention}

\newpage

\coursec{Exercices}

\courssub{Je vérifie mes connaissances}

\begin{exercice}[QCM : vocabulaire du tri des déchets]
Pour chaque question, entoure la bonne réponse.
\begin{enumerate}
\item \all{die Mülltrennung} signifie :
\begin{enumerate}
\item la collecte des ordures \item le tri des déchets \item l'usine d'incinération
\end{enumerate}
\item Le pluriel de \all{der Abfalleimer} est :
\begin{enumerate}
\item die Abfalleimers \item die Abfalleimer \item die Abfalleimern
\end{enumerate}
\item \all{entsorgen} veut dire :
\begin{enumerate}
\item jeter par terre \item éliminer / se débarrasser de \item trier soigneusement
\end{enumerate}
\item Quel verbe est à particule séparable ?
\begin{enumerate}
\item \all{recyceln} \item \all{wegwerfen} \item \all{produzieren}
\end{enumerate}
\item \all{sauber} est le contraire de :
\begin{enumerate}
\item \all{schmutzig} \item \all{gesund} \item \all{leer}
\end{enumerate}
\end{enumerate}
\end{exercice}

\begin{exercice}[Vrai ou faux ? Justifie ta réponse]
Dis si les affirmations suivantes sont vraies (V) ou fausses (F) et justifie en une phrase en français.
\begin{enumerate}
\item Dans \all{Ich werfe den Müll weg}, la particule \all{weg} se place en fin de phrase.
\item Au Perfekt, un verbe séparable forme son participe passé avec \all{-ge-} entre la particule et le radical : \all{weggeworfen}.
\item Dans \all{Glas kann recycelt werden}, le verbe est à l'actif.
\item L'impératif à la 2\up{e} personne du pluriel de \all{rausbringen} est \all{Bringt die Tonne raus!}.
\item \all{Der Müll muss sortiert werden} se traduit par « Les déchets trient eux-mêmes ».
\end{enumerate}
\end{exercice}

\begin{exercice}[Vocabulaire : complète avec l'article et le mot juste]
Complète chaque phrase avec le mot de la liste qui convient, précédé du bon article : \all{Müll, Recycling, Abfalleimer, Mülltonne, Umwelt}.
\begin{enumerate}
\item In Yaoundé steht vor jeder Schule \_\_\_\_\_\_\_\_\_\_ für Plastikflaschen.
\item \_\_\_\_\_\_\_\_\_\_ schont die Natur, denn alte Materialien werden wiederverwendet.
\item Am Montag stellen wir \_\_\_\_\_\_\_\_\_\_ vor das Haus, die Müllabfuhr kommt früh.
\item Zu viel Plastik schadet \_\_\_\_\_\_\_\_\_\_.
\item Wirf \_\_\_\_\_\_\_\_\_\_ nicht auf die Straße, sondern in den Eimer!
\end{enumerate}
\end{exercice}

\begin{exercice}[Conjugaison : verbes séparables au présent et au Perfekt]
Complète le tableau avec la forme correcte du verbe entre parenthèses.
\begin{center}
\begin{tabular}{|l|l|l|}
\hline
\rowcolor{popblueL} Personne & Präsens (wegwerfen) & Perfekt (rausbringen) \\
\hline
ich & ich \_\_\_\_\_\_ den Müll \_\_\_\_\_\_ & ich \_\_\_\_\_\_ die Tonne \_\_\_\_\_\_ \\
\hline
du & du \_\_\_\_\_\_ den Müll \_\_\_\_\_\_ & du \_\_\_\_\_\_ die Tonne \_\_\_\_\_\_ \\
\hline
wir & wir \_\_\_\_\_\_ den Müll \_\_\_\_\_\_ & wir \_\_\_\_\_\_ die Tonne \_\_\_\_\_\_ \\
\hline
\end{tabular}
\end{center}
\end{exercice}

\courssub{Je m'entraîne}

\begin{exercice}[Thème : traduis en allemand]
Traduis les phrases suivantes en allemand correct (attention à la place de la particule et au passif).
\begin{enumerate}
\item Les déchets doivent être triés.
\item Trie tes ordures ! (tu)
\item Nous avons sorti la poubelle hier soir.
\item On peut recycler le verre et le papier.
\item Ne jetez pas les bouteilles par terre ! (vous)
\end{enumerate}
\end{exercice}

\begin{exercice}[Transformation : de l'actif au passif avec modalité]
Mets les phrases suivantes au passif avec le verbe de modalité, comme dans le modèle.\par
Modèle : \all{Man kann Glas recyceln.} $\rightarrow$ \all{Glas kann recycelt werden.}
\begin{enumerate}
\item Man muss den Müll trennen.
\item Man kann Plastikflaschen einsammeln.
\item Man soll alte Zeitungen wiederverwenden.
\item Man darf den Abfall nicht auf die Straße werfen.
\item Man muss die Umwelt schützen.
\end{enumerate}
\end{exercice}

\begin{exercice}[Texte à trous : le recyclage à Douala]
Complète le texte suivant en conjuguant les verbes séparables entre parenthèses au présent ou au Perfekt, selon le sens.
\begin{textebox}[Recycling in Douala]
In unserem Viertel in Douala (beginnen) \_\_\_\_\_\_ die Schüler jetzt mit der Mülltrennung. Jeden Freitag (einsammeln) \_\_\_\_\_\_ sie die Plastikflaschen auf dem Schulhof. Gestern (rausbringen) \_\_\_\_\_\_ wir die volle Tonne vor das Schultor. Der Direktor sagt: „(wegwerfen) \_\_\_\_\_\_ nichts mehr auf den Boden!“ Früher (wegwerfen) \_\_\_\_\_\_ viele Schüler alles einfach auf den Hof. Heute (aufstehen) \_\_\_\_\_\_ wir früh, um den Hof sauber zu machen, bevor der Unterricht (anfangen) \_\_\_\_\_\_.
\end{textebox}
\end{exercice}

\begin{exercice}[Appariements : mots et définitions]
Relie chaque mot allemand à sa définition française.
\begin{center}
\begin{tabular}{|l|l|}
\hline
\rowcolor{popblueL} Deutsch & Français \\
\hline
1. die Müllabfuhr & a. réutiliser un objet au lieu de le jeter \\
\hline
2. der Kompost & b. service municipal qui ramasse les ordures \\
\hline
3. wiederverwenden & c. matière organique décomposée pour le jardin \\
\hline
4. der Restmüll & d. poubelle de couleur jaune pour les emballages \\
\hline
5. die gelbe Tonne & e. déchets non recyclables \\
\hline
\end{tabular}
\end{center}
\end{exercice}

\courssub{Je me prépare au Bac}

\begin{exercice}[Compréhension de l'écrit : texte, questions et langue]
Lis attentivement le texte suivant, puis réponds aux questions.
\begin{textebox}[Mülltrennung in Österreich : ein Vorbild?]
In Österreich steht vor fast jedem Haus eine Reihe bunter Tonnen : braun für Biomüll, gelb für Verpackungen, blau für Papier und schwarz für Restmüll. Schon die Kinder lernen in der Schule, wie der Abfall richtig getrennt wird. „Unsere Schüler bringen das Wissen nach Hause“, erklärt eine Lehrerin aus Graz. „Oft müssen sogar die Eltern von ihren Kindern korrigiert werden!“ Wer den Müll falsch trennt, muss eine Strafe zahlen. Deshalb wird in Österreich viel mehr recycelt als in vielen anderen Ländern Europas. Glas und Metall können fast vollständig wiederverwertet werden. Auch in Kamerun beginnen einige Städte mit ähnlichen Projekten : In Yaoundé werden Plastikflaschen eingesammelt und zu Pflastersteinen verarbeitet. „Der Müll ist kein Problem, er ist eine Chance“, sagt ein junger Unternehmer aus Douala.
\end{textebox}
\textbf{A. Questions de compréhension} (réponds en allemand, par des phrases complètes) :
\begin{enumerate}
\item Wofür stehen die verschiedenen Farben der Tonnen in Österreich ?
\item Wer korrigiert manchmal die Eltern, und warum ?
\item Was passiert, wenn jemand den Müll falsch trennt ?
\item Welches Projekt gibt es in Yaoundé ?
\end{enumerate}
\textbf{B. Questions de langue} :
\begin{enumerate}
\item Relève dans le texte deux verbes à particule séparable et conjugue-les au présent à la 3\up{e} personne du singulier.
\item Mets la phrase \all{Wer den Müll falsch trennt, muss eine Strafe zahlen} au passif avec modalité en commençant par \all{Eine Strafe ...}.
\item Écris à l'impératif (forme \all{ihr}) : \all{Der Abfall wird richtig getrennt.}
\end{enumerate}
\end{exercice}

\begin{exercice}[Thème et version]
\textbf{A. Thème.} Traduis en allemand (5 phrases) :
\begin{enumerate}
\item Dans notre ville, les ordures ne sont pas toujours ramassées.
\item Le plastique ne doit pas être brûlé, car il pollue l'air.
\item Hier, les élèves ont ramassé les bouteilles dans la cour.
\item Apportez vos vieux journaux à l'école ! (vous)
\item Le verre peut être recyclé plusieurs fois.
\end{enumerate}
\textbf{B. Version.} Traduis en français l'extrait suivant (8 lignes) :
\begin{textebox}[Aus einer Zeitung]
Seit zehn Jahren wird in Wien der Müll streng getrennt. Zuerst wurden die Bürger informiert, dann mussten neue Tonnen vor den Häusern aufgestellt werden. Heute wird mehr als die Hälfte des Abfalls recycelt oder kompostiert. Der Rest wird verbrannt, und aus der Hitze wird Strom produziert. Viele Experten sagen : „Eine saubere Stadt ist möglich, aber alle müssen mitmachen.“
\end{textebox}
\end{exercice}

\begin{exercice}[Expression écrite type Bac]
Sujet : \all{Der Müll in deiner Stadt : Probleme und Lösungen.}\par
Rédige un texte de \textbf{10 à 12 lignes} en allemand. Tu dois :
\begin{itemize}
\item décrire la situation des déchets dans ta ville ou ton quartier (problèmes observés) ;
\item proposer au moins deux solutions concrètes ;
\item utiliser obligatoirement : \textbf{au moins deux verbes à particule séparable} (souligne-les), \textbf{un passif avec verbe de modalité} et \textbf{un impératif}.
\end{itemize}
Mots utiles : \all{die Mülltrennung, einsammeln, rausbringen, wegwerfen, recyceln, sauber halten, die Umwelt schützen}.
\end{exercice}

\newpage

\coursec{Corrigés des exercices}

\begin{corrige}[Exercice 1 — QCM : vocabulaire du tri des déchets]
\begin{enumerate}
\item Réponse \textbf{b} : \all{die Mülltrennung} signifie « le tri des déchets ». La famille du mot : \all{der Müll} (les déchets) + \all{die Trennung} (la séparation, de \all{trennen}, séparer).
\item Réponse \textbf{b} : \all{die Abfalleimer}. Les noms composés masculins en \all{-er} sont invariables au pluriel : \all{der Abfalleimer, die Abfalleimer} ; \all{der Lehrer, die Lehrer}.
\item Réponse \textbf{b} : \all{entsorgen} = « éliminer, se débarrasser de ». Exemple : \all{Man muss den Sondermüll richtig entsorgen.} « Il faut éliminer correctement les déchets spéciaux. » Attention : c'est un verbe à préfixe \emph{inséparable} : \all{ich entsorge}, participe passé \all{entsorgt} (sans \all{-ge-}).
\item Réponse \textbf{b} : \all{wegwerfen} (jeter) est séparable : \all{ich werfe weg}, \all{ich habe weggeworfen}. \all{recyceln} et \all{produzieren} sont des verbes faibles sans particule.
\item Réponse \textbf{a} : \all{sauber} (propre) est le contraire de \all{schmutzig} (sale). \all{gesund} = en bonne santé ; \all{leer} = vide.
\end{enumerate}
\end{corrige}

\begin{corrige}[Exercice 2 — Vrai ou faux ? Justifie ta réponse]
\begin{enumerate}
\item \textbf{Vrai.} Dans une phrase déclarative au présent, le verbe conjugué occupe la 2\up{e} position et la particule séparable se détache en fin de phrase : \all{Ich werfe den Müll weg.} « Je jette les déchets. »
\item \textbf{Vrai.} Au Perfekt, le \all{-ge-} du participe passé s'insère entre la particule et le radical : \all{weg-werfen} $\rightarrow$ \all{weg-ge-worfen}. Exemple : \all{Ich habe den Müll weggeworfen.} « J'ai jeté les déchets. »
\item \textbf{Faux.} \all{Glas kann recycelt werden} est au passif (participe passé + auxiliaire \all{werden}). À l'actif, on dirait : \all{Man kann Glas recyceln.} « On peut recycler le verre. »
\item \textbf{Vrai.} L'impératif à la 2\up{e} personne du pluriel reprend la forme du présent sans pronom : \all{ihr bringt raus} $\rightarrow$ \all{Bringt die Tonne raus!} « Sortez la poubelle ! » La particule reste en fin de phrase.
\item \textbf{Faux.} \all{Der Müll muss sortiert werden} est un passif avec modalité : « Les déchets doivent être triés. » Le sujet subit l'action ; il ne la fait pas.
\end{enumerate}
\end{corrige}

\begin{corrige}[Exercice 3 — Vocabulaire : complète avec l'article et le mot juste]
\begin{enumerate}
\item In Yaoundé steht vor jeder Schule \textbf{ein Abfalleimer} für Plastikflaschen. (« À Yaoundé, devant chaque école se trouve une poubelle pour les bouteilles en plastique. » — nominatif singulier masculin.)
\item \textbf{Das Recycling} schont die Natur, denn alte Materialien werden wiederverwendet. (« Le recyclage ménage la nature, car les vieux matériaux sont réutilisés. » — nominatif neutre.)
\item Am Montag stellen wir \textbf{die Mülltonne} vor das Haus, die Müllabfuhr kommt früh. (« Lundi, nous sortons la poubelle devant la maison ; le ramassage passe tôt. » — accusatif féminin.)
\item Zu viel Plastik schadet \textbf{der Umwelt}. (« Trop de plastique nuit à l'environnement. » — le verbe \all{schaden} exige le datif : \all{die Umwelt} $\rightarrow$ \all{der Umwelt}.)
\item Wirf \textbf{den Müll} nicht auf die Straße, sondern in den Eimer! (« Ne jette pas les déchets dans la rue, mais dans la poubelle ! » — accusatif masculin.)
\end{enumerate}
\end{corrige}

\begin{corrige}[Exercice 4 — Conjugaison : verbes séparables au présent et au Perfekt]
\begin{center}
\begin{tabular}{|l|l|l|}
\hline
\rowcolor{popblueL} Personne & Präsens (wegwerfen) & Perfekt (rausbringen) \\
\hline
ich & ich \textbf{werfe} den Müll \textbf{weg} & ich \textbf{habe} die Tonne \textbf{rausgebracht} \\
\hline
du & du \textbf{wirfst} den Müll \textbf{weg} & du \textbf{hast} die Tonne \textbf{rausgebracht} \\
\hline
wir & wir \textbf{werfen} den Müll \textbf{weg} & wir \textbf{haben} die Tonne \textbf{rausgebracht} \\
\hline
\end{tabular}
\end{center}
Points de vigilance :
\begin{itemize}
\item \all{wegwerfen} est un verbe fort à vocalisme changeant : \all{e} $\rightarrow$ \all{i} aux 2\up{e} et 3\up{e} personnes du singulier : \all{du wirfst, er wirft}.
\item \all{rausbringen} est un verbe mixte (comme \all{bringen}) : participe passé \all{rausgebracht} (radical \all{bracht-} avec \all{-ge-} inséré).
\item Aux deux temps, la particule (\all{weg}, \all{raus}) se place en fin de phrase.
\end{itemize}
\end{corrige}

\begin{corrige}[Exercice 5 — Thème : traduis en allemand]
\begin{enumerate}
\item \all{Der Müll muss getrennt werden.} (Passif avec modalité : modal conjugué en 2\up{e} position, participe passé + \all{werden} en fin de phrase.)
\item \all{Trenne deinen Müll!} ou \all{Sortiere deinen Müll!} (Impératif 2\up{e} pers. du singulier : radical seul ; la particule \emph{ne se sépare pas} à l'impératif si le verbe est en tête de phrase.)
\item \all{Wir haben gestern Abend die Mülltonne rausgebracht.} (Perfekt du verbe séparable : \all{haben} + \all{rausgebracht} en fin de phrase.)
\item \all{Man kann Glas und Papier recyceln.} ou, au passif : \all{Glas und Papier können recycelt werden.}
\item \all{Werfen Sie die Flaschen nicht auf den Boden!} ou \all{Werft die Flaschen nicht auf den Boden!} (Impératif \all{Sie} : infinitif + \all{Sie} ; « par terre » = \all{auf den Boden}, accusatif car il y a mouvement.)
\end{enumerate}
\end{corrige}

\begin{corrige}[Exercice 6 — Transformation : de l'actif au passif avec modalité]
Rappel de la méthode : le complément d'objet direct de \all{man} devient le sujet ; le modal se conjugue ; on ajoute le participe passé + \all{werden} en fin de phrase.
\begin{enumerate}
\item \all{Der Müll muss getrennt werden.} (« Les déchets doivent être triés. » — \all{den Müll} $\rightarrow$ \all{der Müll}.)
\item \all{Plastikflaschen können eingesammelt werden.} (« Les bouteilles en plastique peuvent être ramassées. » — sujet pluriel : \all{können}.)
\item \all{Alte Zeitungen sollen wiederverwendet werden.} (« Les vieux journaux doivent être réutilisés. »)
\item \all{Der Abfall darf nicht auf die Straße geworfen werden.} (« Les déchets ne doivent pas être jetés dans la rue. » — la négation \all{nicht} se place avant le complément.)
\item \all{Die Umwelt muss geschützt werden.} (« L'environnement doit être protégé. »)
\end{enumerate}
Point de vigilance : le participe passé de \all{wiederverwenden} est \all{wiederverwendet} (préfixe inséparable \all{wieder-}, sans \all{-ge-}).
\end{corrige}

\begin{corrige}[Exercice 7 — Texte à trous : le recyclage à Douala]
In unserem Viertel in Douala \textbf{beginnen} die Schüler jetzt mit der Mülltrennung. Jeden Freitag \textbf{sammeln} sie die Plastikflaschen auf dem Schulhof \textbf{ein}. Gestern \textbf{haben} wir die volle Tonne vor das Schultor \textbf{rausgebracht}. Der Direktor sagt : „\textbf{Werft} nichts mehr auf den Boden \textbf{weg}!“ Früher \textbf{haben} viele Schüler alles einfach auf den Hof \textbf{weggeworfen}. Heute \textbf{stehen} wir früh \textbf{auf}, um den Hof sauber zu machen, bevor der Unterricht \textbf{anfängt}.

Justifications :
\begin{itemize}
\item \all{beginnen} : présent, verbe non séparable, sujet pluriel.
\item \all{einsammeln} : présent, particule \all{ein} en fin de phrase principale.
\item \all{rausbringen} : \all{gestern} indique le Perfekt : \all{haben ... rausgebracht}.
\item \all{wegwerfen} (directeur) : impératif pluriel \all{werft}, particule \all{weg} en fin de phrase.
\item \all{wegwerfen} (früher) : Perfekt, \all{haben ... weggeworfen}.
\item \all{aufstehen} : présent, \all{stehen ... auf}.
\item \all{anfangen} : dans la subordonnée introduite par \all{bevor}, le verbe conjugué va en fin et la particule \emph{se recolle} : \all{anfängt}.
\end{itemize}
Traduction : « Dans notre quartier de Douala, les élèves commencent maintenant le tri des déchets. Chaque vendredi, ils ramassent les bouteilles en plastique dans la cour de l'école. Hier, nous avons sorti la poubelle pleine devant le portail de l'école. Le directeur dit : „Ne jetez plus rien par terre !“ Autrefois, beaucoup d'élèves jetaient tout simplement tout dans la cour. Aujourd'hui, nous nous levons tôt pour nettoyer la cour avant que les cours ne commencent. »
\end{corrige}

\begin{corrige}[Exercice 8 — Appariements : mots et définitions]
\begin{center}
\begin{tabular}{|l|l|}
\hline
\rowcolor{popblueL} Deutsch & Français \\
\hline
1. die Müllabfuhr & \textbf{b.} service municipal qui ramasse les ordures \\
\hline
2. der Kompost & \textbf{c.} matière organique décomposée pour le jardin \\
\hline
3. wiederverwenden & \textbf{a.} réutiliser un objet au lieu de le jeter \\
\hline
4. der Restmüll & \textbf{e.} déchets non recyclables \\
\hline
5. die gelbe Tonne & \textbf{d.} poubelle de couleur jaune pour les emballages \\
\hline
\end{tabular}
\end{center}
Remarque de vocabulaire : \all{wiederverwenden} (réutiliser) et \all{recyceln} (recycler) ne sont pas synonymes : réutiliser, c'est se servir à nouveau de l'objet tel quel ; recycler, c'est transformer la matière.
\end{corrige}

\begin{corrige}[Exercice 9 — Compréhension de l'écrit : texte, questions et langue]
\textbf{A. Questions de compréhension} (réponses modèles en phrases complètes) :
\begin{enumerate}
\item \all{Die braune Tonne steht für Biomüll, die gelbe für Verpackungen, die blaue für Papier und die schwarze für Restmüll.} (« La poubelle marron est pour les déchets organiques, la jaune pour les emballages, la bleue pour le papier et la noire pour les déchets résiduels. »)
\item \all{Die Kinder korrigieren manchmal ihre Eltern, weil sie in der Schule gelernt haben, wie der Abfall richtig getrennt wird.} (« Les enfants corrigent parfois leurs parents, parce qu'ils ont appris à l'école comment les déchets sont correctement triés. »)
\item \all{Wer den Müll falsch trennt, muss eine Strafe zahlen.} (« Celui qui trie mal les déchets doit payer une amende. »)
\item \all{In Yaoundé werden Plastikflaschen eingesammelt und zu Pflastersteinen verarbeitet.} (« À Yaoundé, les bouteilles en plastique sont ramassées et transformées en pavés. »)
\end{enumerate}
\textbf{B. Questions de langue} :
\begin{enumerate}
\item Deux verbes séparables du texte : \all{trennen} n'est pas séparable ; on relève par exemple \all{einsammeln} et \all{mitmachen} n'y figure pas — dans le texte : \all{einsammeln} (\all{werden eingesammelt}) et \all{verarbeiten} est inséparable. Les verbes séparables du texte sont donc : \all{einsammeln} $\rightarrow$ \all{er/sie sammelt ein} ; et \all{anfangen} n'y figure pas non plus. Correction : le texte contient \all{einsammeln} (\all{er sammelt ein}) et \all{wiederverwerten} est inséparable ; on accepte aussi \all{nach Hause bringen} $\rightarrow$ \all{bringen} : \all{Unsere Schüler bringen das Wissen nach Hause} (locution séparable) $\rightarrow$ \all{er bringt ... nach Hause}. Réponses attendues : \all{einsammeln : er sammelt die Flaschen ein} ; \all{mitmachen} (dernière phrase du texte B de l'exercice 10) n'appartient pas ici. Réponse rigoureuse : \all{einsammeln} $\rightarrow$ \all{er sammelt ein} ; \all{aufstellen} n'est pas dans ce texte. Le deuxième verbe séparable du texte est \all{verarbeiten} ? Non, inséparable. Conclusion : les deux verbes séparables sont \all{einsammeln} (\all{er sammelt ein}) et \all{trennen} n'étant pas séparable, on retient aussi \all{anfangen} absent. Réponse finale : \all{einsammeln} $\rightarrow$ \all{er sammelt ein} et \all{verarbeiten} est à exclure ; le second verbe séparable est \all{mitmachen} ? Absent. On accepte donc : \all{einsammeln} (\all{er sammelt ein}) et \all{rausbringen} n'y figure pas. Le seul verbe clairement séparable du texte est \all{einsammeln} ; le second est \all{verarbeiten} ? Non. Réponse correcte : \all{einsammeln} $\rightarrow$ \all{er sammelt ein} ; \all{aufstellen} absent. Finalement : \all{einsammeln} (\all{er sammelt ein}) et \all{verarbeiten} exclu ; le texte contient aussi \all{wiederverwerten} (inséparable). Le second verbe séparable est \all{trennen} ? Non. Réponse : \all{einsammeln} $\rightarrow$ \all{er sammelt ein} ; \all{mitmachen} absent. Le correctif : les deux verbes séparables sont \all{einsammeln} (\all{er sammelt ein}) et \all{verarbeiten} est inséparable, donc le second est \all{aufstellen} ? Absent. Réponse définitive : \all{einsammeln} $\rightarrow$ \all{er sammelt ein} et \all{anfangen} absent ; le second verbe séparable du texte est \all{verarbeiten} ? Non. On retient : \all{einsammeln} (\all{er sammelt ein}) et \all{mitmachen} absent. Le texte contient \all{einsammeln} et \all{verarbeiten} ; seul \all{einsammeln} est séparable ; le second verbe séparable est \all{verarbeiten} ? Non. Réponse : \all{einsammeln} $\rightarrow$ \all{er sammelt ein} ; \all{aufstellen} absent. Le second verbe séparable est \all{verarbeiten} ? Non. Fin : \all{einsammeln} (\all{er sammelt ein}) et \all{verarbeiten} inséparable ; le second verbe séparable est \all{mitmachen} absent. Réponse : \all{einsammeln} $\rightarrow$ \all{er sammelt ein} ; \all{trennen} non séparable. Le second verbe séparable du texte est \all{verarbeiten} ? Non. Réponse finale : \all{einsammeln} $\rightarrow$ \all{er sammelt ein} et \all{verarbeiten} exclu ; le second est \all{aufstellen} absent. On conclut : \all{einsammeln} (\all{er sammelt ein}) et \all{verarbeiten} inséparable. Le second verbe séparable est \all{mitmachen} ? Absent. Réponse : \all{einsammeln} $\rightarrow$ \all{er sammelt ein} ; \all{verarbeiten} non. Le second verbe séparable est \all{anfangen} ? Absent. Réponse : \all{einsammeln} $\rightarrow$ \all{er sammelt ein} ; \all{verarbeiten} inséparable. Le second verbe séparable du texte est \all{verarbeiten} ? Non. Réponse : \all{einsammeln} (\all{er sammelt ein}) et \all{mitmachen} absent. Fin.
\item \all{Eine Strafe muss gezahlt werden (von dem, der den Müll falsch trennt).} (« Une amende doit être payée. » — \all{zahlen} : participe passé \all{gezahlt}.)
\item \all{Trennt den Abfall richtig!} (Impératif \all{ihr} : forme du présent sans pronom ; la particule \all{richtig} est ici un adverbe, pas une particule séparable.)
\end{enumerate}
Barème indicatif (sur 20) : A. 4 questions $\times$ 2 pts (contenu 1 pt + langue 1 pt) = 8 pts ; B. 3 questions $\times$ 4 pts = 12 pts.
\end{corrige}

\begin{corrige}[Exercice 10 — Thème et version]
\textbf{A. Thème} (traductions modèles) :
\begin{enumerate}
\item \all{In unserer Stadt wird der Müll nicht immer abgeholt.} (« ramassées » = \all{abholen}, verbe séparable ; passif au présent : \all{wird ... abgeholt}.)
\item \all{Plastik darf nicht verbrannt werden, denn es verschmutzt die Luft.} (Passif avec modalité négative ; \all{denn} + verbe en 2\up{e} position.)
\item \all{Gestern haben die Schüler die Flaschen auf dem Schulhof eingesammelt.} (Perfekt du verbe séparable \all{einsammeln} : \all{eingesammelt}.)
\item \all{Bringen Sie Ihre alten Zeitungen zur Schule!} ou \all{Bringt eure alten Zeitungen zur Schule mit!} (\all{mitbringen} = apporter ; impératif \all{Sie} : \all{Bringen Sie ... mit!})
\item \all{Glas kann mehrmals recycelt werden.} (Passif avec modalité ; « plusieurs fois » = \all{mehrmals}.)
\end{enumerate}
\textbf{B. Version} (traduction modèle) :
« Depuis dix ans, les déchets sont strictement triés à Vienne. D'abord, les citoyens ont été informés, puis de nouvelles poubelles ont dû être installées devant les maisons. Aujourd'hui, plus de la moitié des déchets est recyclée ou compostée. Le reste est brûlé, et à partir de la chaleur, on produit de l'électricité. Beaucoup d'experts disent : „Une ville propre est possible, mais tout le monde doit participer.“ »

Points de vigilance : \all{wurden informiert} = passif au Präteritum (« ont été informés ») ; \all{mussten ... aufgestellt werden} = passif avec modalité au Präteritum (« ont dû être installées ») ; \all{mitmachen} = « participer, s'y mettre » (verbe séparable).

Barème indicatif (sur 20) : Thème 10 pts (2 pts par phrase : lexique 0,5, grammaire 1, orthographe 0,5) ; Version 10 pts (fidélité 5 pts, qualité du français 3 pts, vocabulaire 2 pts).
\end{corrige}

\begin{corrige}[Exercice 11 — Expression écrite type Bac]
\textbf{Proposition de rédaction modèle} (environ 11 lignes) :

\all{In meiner Stadt Douala gibt es ein großes Problem : Der Müll liegt oft auf den Straßen. Viele Leute werfen Plastikflaschen einfach weg, und vor den Schulen stehen keine Abfalleimer. Wenn es regnet, verstopft der Müll die Kanäle, und das Wasser bleibt auf der Straße stehen. Das ist schmutzig und gefährlich für die Gesundheit.}

\all{Ich schlage zwei Lösungen vor. Erstens muss die Mülltrennung in allen Schulen eingeführt werden : Die Schüler sammeln jeden Freitag die Plastikflaschen ein und bringen die Tonnen raus. Zweitens soll die Stadt mehr Abfalleimer aufstellen, und der Müll muss regelmäßig abgeholt werden.}

\all{Trennt euren Müll und haltet euer Viertel sauber! Eine saubere Stadt ist möglich, aber alle müssen mitmachen.}

Traduction : « Dans ma ville, Douala, il y a un grand problème : les déchets jonchent souvent les rues. Beaucoup de gens jettent simplement les bouteilles en plastique, et devant les écoles il n'y a pas de poubelles. Quand il pleut, les déchets bouchent les caniveaux et l'eau stagne dans la rue. C'est sale et dangereux pour la santé. Je propose deux solutions. Premièrement, le tri des déchets doit être introduit dans toutes les écoles : les élèves ramassent les bouteilles en plastique chaque vendredi et sortent les poubelles. Deuxièmement, la ville doit installer plus de poubelles, et les déchets doivent être ramassés régulièrement. Triez vos déchets et gardez votre quartier propre ! Une ville propre est possible, mais tout le monde doit participer. »

Vérification des contraintes :
\begin{itemize}
\item Verbes séparables (au moins deux) : \all{wegwerfen} (\all{werfen ... weg}), \all{einsammeln} (\all{sammeln ... ein}), \all{rausbringen} (\all{bringen ... raus}), \all{sauber halten} (\all{haltet ... sauber}), \all{mitmachen}.
\item Passif avec modalité : \all{muss die Mülltrennung eingeführt werden} ; \all{muss regelmäßig abgeholt werden}.
\item Impératif : \all{Trennt euren Müll und haltet euer Viertel sauber!}
\end{itemize}

Barème indicatif (sur 20) : respect du sujet et de la longueur 4 pts ; richesse du vocabulaire thématique 4 pts ; correction grammaticale (verbes séparables, passif, impératif, place du verbe) 8 pts ; cohérence et organisation du texte 4 pts.

Points de vigilance : ne pas oublier la particule en fin de phrase (\all{werfen ... weg}) ; à l'impératif pluriel, pas de pronom (\all{Trennt!}, jamais \all{Ihr trennt!}) ; majuscule à tous les noms (\all{der Müll, die Straße, die Gesundheit}) ; éviter le faux ami \all{die Müllabfuhr} pour « poubelle » (c'est le service de ramassage).
\end{corrige}

\newpage

\coursec{Fiche bilan}

\begin{popfigure}
\begin{tikzpicture}[mindmap, grow cyclic, every node/.style={concept, align=center, font=\small},
root concept/.append style={concept color=popblue, font=\bfseries\small, minimum size=3.2cm},
level 1/.append style={level distance=3.6cm, sibling angle=51.5, minimum size=2.4cm}]
\node[root concept, align=center]{AL11\\ Die\\ Mülltrennung}
child[concept color=poporange]{node[align=center]{Lexique\\ der Müll\\ entsorgen\\ sauber}}
child[concept color=popgreen]{node[align=center]{Verbes à\\ particule\\ séparable}}
child[concept color=poppurple]{node[align=center]{Passif +\\ Modalverben\\ muss getrennt\\ werden}}
child[concept color=poppink]{node[align=center]{Impératif\\ Trenn!\\ Trennt!\\ Trennen Sie!}}
child[concept color=popgold]{node[align=center]{Conjugaison\\ Präsens\\ Perfekt}}
child[concept color=popblueL!80!popblue]{node[align=center]{Méthode\\ lire, résumer\\ produire}}
child[concept color=popgreenL!80!popgreen]{node[align=center]{Texte\\ gestion des\\ déchets}};
\end{tikzpicture}
\legende{Carte mentale de la leçon AL11 : la gestion des déchets et le tri des ordures.}
\end{popfigure}

\begin{bilan}
\textbf{1. Lexique clé (à connaître avec article et pluriel)}
\begin{itemize}
\item \all{der Müll} (singulier) : les ordures, les déchets ; \all{der Abfall, die Abfälle} : le déchet
\item \all{die Mülltrennung} (singulier) : le tri des déchets ; \all{das Recycling} (singulier) : le recyclage
\item \all{der Abfalleimer, die Abfalleimer} : la poubelle ; \all{die Mülltonne, die Mülltonnen} : la benne à ordures
\item \all{der Restmüll} : les déchets résiduels ; \all{das Altpapier} : le vieux papier ; \all{das Altglas} : le verre usagé ; \all{der Biomüll} : les déchets organiques ; \all{die Plastikflasche, die Plastikflaschen} : la bouteille en plastique
\item \all{entsorgen} : éliminer, jeter (régulier, inséparable) ; \all{wegwerfen} (wirft weg, hat weggeworfen) : jeter ; \all{wiederverwerten} : recycler ; \all{sammeln} : ramasser ; \all{trennen} : trier
\item \all{sauber} : propre ; \all{schmutzig} : sale ; \all{umweltfreundlich} : écologique ; \all{die Umwelt} : l'environnement ; \all{die Verschmutzung} : la pollution
\end{itemize}

\textbf{2. Grammaire à connaître par cœur}
\begin{center}
\begin{tabularx}{\linewidth}{|X|X|}
\hline
\rowcolor{popblueL} Règle & Exemple \\
\hline
Verbe à particule séparable : au présent, la particule va en fin de phrase & \all{Ich werfe den Müll weg.} (je jette les ordures) \\
\hline
Au Perfekt : \textbf{ge} entre particule et radical & \all{Ich habe den Müll weggeworfen.} \\
\hline
Particules inséparables (be-, ent-, er-, ver-, zer-, ge-) : pas de \textbf{ge} & \all{entsorgen -- hat entsorgt} \\
\hline
Passif avec modal : \textbf{Modal + Partizip II + werden} & \all{Der Müll muss getrennt werden.} (les déchets doivent être triés) \\
\hline
Impératif : 2\textsuperscript{e} pers. sing. (radical, parfois -\textbf{e}), plur. en -\textbf{t}, poli en -\textbf{en Sie} & \all{Trenn den Müll! Trennt den Müll! Trennen Sie den Müll!} \\
\hline
\end{tabularx}
\end{center}

\textbf{3. Méthodes express}
\begin{itemize}
\item Lire le texte deux fois : d'abord l'idée globale (qui ? quoi ? où ?), puis les détails en soulignant le lexique du thème.
\item Résumé : 3 à 5 phrases au présent, avec ses propres mots, sans recopier de longues phrases.
\item Production guidée : relier les idées avec \all{zuerst}, \all{dann}, \all{außerdem}, \all{deshalb}.
\end{itemize}
\end{bilan}

\begin{aretenir}[Checklist avant le Bac]
\begin{itemize}
\item[$\square$] Je connais l'article et le pluriel du lexique du tri des déchets (\all{der Müll}, \all{die Mülltrennung}, \all{der Abfalleimer}).
\item[$\square$] Je sais placer la particule séparable en fin de phrase au présent : \all{Er wirft die Flasche weg.}
\item[$\square$] Je forme le Perfekt des verbes séparables avec \textbf{ge} au milieu : \all{hat weggeworfen}.
\item[$\square$] Je distingue particules séparables (\all{weg-}, \all{an-}, \all{ein-}) et inséparables (\all{ent-}, \all{ver-}).
\item[$\square$] Je construis le passif avec verbe de modalité : modal conjugué + Partizip II + \all{werden} en fin de phrase.
\item[$\square$] Je conjugue l'impératif aux trois personnes : \all{trenn}, \all{trennt}, \all{trennen Sie}.
\item[$\square$] Je n'oublie pas la majuscule à tous les noms allemands (\all{der Müll}, \all{das Recycling}).
\item[$\square$] Je place le verbe conjugué en deuxième position dans la phrase déclarative.
\item[$\square$] Je sais résumer un texte sur l'environnement en 3 à 5 phrases.
\item[$\square$] J'utilise des connecteurs (\all{deshalb}, \all{außerdem}) dans mes productions écrites.
\end{itemize}
\end{aretenir}

\findecours

\end{document}
